\documentclass[11pt]{amsart}
\usepackage[margin=1.05in]{geometry}
\usepackage{amsmath,amssymb,amsthm,longtable,array}
\usepackage[hidelinks]{hyperref}
\newcommand{\toprule}{\hline}
\newcommand{\midrule}{\hline}
\newcommand{\bottomrule}{\hline}
\raggedbottom
\numberwithin{equation}{section}
\newcommand{\AutoTag}[1]{\refstepcounter{equation}\tag{\theequation}\label{#1}}
\newcommand{\unnumtable}{\addtocounter{table}{-1}}
\newtheorem{theorem}{Theorem}[section]
\newtheorem{lemma}[theorem]{Lemma}
\newtheorem{proposition}[theorem]{Proposition}
\newtheorem{rowproposition}[theorem]{Proposition}
\newtheorem{corollary}[theorem]{Corollary}
\theoremstyle{definition}
\newtheorem{definition}[theorem]{Definition}
\theoremstyle{remark}
\newtheorem{remark}[theorem]{Remark}
\newcommand{\rad}{\operatorname{rad}}
\newcommand{\tauDiv}{\tau_{\mathrm{div}}}
\newcommand{\e}{\varepsilon}
\newcommand{\Om}{\Omega}
\newcommand{\cN}{\mathcal N}
\newcommand{\cR}{\mathcal R}
\newcommand{\cM}{\mathcal M}
\newcommand{\lls}{\ll_{\!\e}}
\newcommand{\1}{\mathbf 1}
\title{A power saving below two-thirds for equal radicals of consecutive products}
\author{Ruiyi Zhang}
\address{Independent Researcher, Columbia, Maryland, USA}
\email{ryzfsu@gmail.com}
\subjclass[2020]{Primary 11A25; Secondary 11D25, 11N37}
\keywords{radicals, consecutive integers, quadratic fibres, determinant method}
\hypersetup{pdftitle={A power saving below two-thirds for equal radicals of consecutive products},pdfauthor={Ruiyi Zhang}}
\begin{document}
\begin{abstract}
Let \(F(X)\) count positive integer pairs \(m<n\le X\) with
\(\operatorname{rad}(m(m+1))=\operatorname{rad}(n(n+1))\). We prove
\(F(X)\ll_\varepsilon X^{109/165+\varepsilon}\) for every
\(\varepsilon>0\), improving the earlier exponent \(2/3\). A four-corner
parametrization divides these pairs into thirteen classes. Uniform estimates
for those classes combine with a finite exponent calculation and one choice
of error terms.
\end{abstract}

\maketitle
\section{Introduction}

For a positive integer \(r\), let
\(\operatorname{rad}(r)=\prod_{p\mid r}p\). We count ordered pairs
\[
 F(X)=\#\{(m,n)\in\mathbb Z_{>0}^2:m<n\le X,\
       \operatorname{rad}(m(m+1))=\operatorname{rad}(n(n+1))\}.
\]

\begin{theorem}\label{thm:main}
For every \(\varepsilon>0\) and \(X\ge2\),
\[
                         F(X)\ll_\varepsilon X^{109/165+\varepsilon}.
\]
\end{theorem}

The exponent is \(1/165\) below \(2/3\). Section~\ref{sec:notation}
introduces the four-corner gcd coordinates. It also records the
source-specific identities used here. The new part has three pieces: a
first-match partition into thirteen disjoint populations, uniform
estimates for its new rows, and a finite optimization of the long-shift
remainder.

The minimum \(S_0\) below is the limiting minimum of the row savings.
After the analytic loss \(\beta_0\) and a conservative long-row charge
\(g_0/4\),
the proof retains the smaller saving displayed below.

The earlier result \cite[Theorem~1]{Paper1} gives exponent \(2/3\) for
this counting function. Its support-fibre lemma and low-radical estimate
are used below; Appendix~\ref{app:sources} states them and records the
substitutions. The quotient-tail and Farey-spacing estimates in that paper
are context only. They are not used. Two related papers by the author
\cite{Zhang712,ZhangE20} count primitive \(abc\) triples, not the pairs
counted here, and do not enter this proof.

At the chosen rational point, the minimum row saving is
\[
 S_0=\frac{30869}{5000000},\qquad
 g_0=S_0-\frac1{165}=\frac{18677}{165000000}.
\]
The source-uniform estimates spend a total analytic loss
\(\beta_0=g_0/16\). The long row has three finite corrections: geometric,
interval, and high-degree. Beyond the common threshold each is at most
\(g_0/4\). We charge their maximum, not their sum. The remaining margin is
\[
 S_0-\beta_0-\frac{g_0}{4}
 =\frac1{165}+\frac{11g_0}{16}>\frac1{165}.
\]

Sections~\ref{sec:notation}--\ref{sec:partition} give the canonical source
and the exact row partition. Sections~\ref{sec:outer}--\ref{sec:intermediate}
treat the outer, joint-corner, short-shift and intermediate rows.
Section~\ref{sec:long} handles the long row. Sections~\ref{sec:ledger}--\ref{sec:conclusion}
allocate losses and sum the dyadic shells. The appendices contain the
self-contained square-fibre proof, the application of the cited estimates, and the
branch arithmetic. The use of AI tools is described after the appendices.

We use Heath--Brown's published Lemma~2 under its stated bihomogeneous
hypotheses. The fixed cubic branch uses the unnumbered theorem of
Bombieri--Schmidt with its 1989 correction \cite{BS1989}. Appendix~\ref{app:sources}
states both inputs and records the applications.
Reuss's Theorem~1 \cite{Reuss} is the determinant-method input behind the low-radical
estimate in \cite{Paper1}. That estimate is used here. Tao's
Lemma~2.10 gives a comparison with the square-fibre estimate of
Appendix~\ref{app:fibre}. The proof uses that appendix theorem.

\section{Canonical source and parameters}\label{sec:notation}

For an integer \(Y\ge2\), let
\[
 \Omega_Y=\{(m,n)\in\mathbb Z_{>0}^2:
       m<n\le Y,\ Y/2<n,\
       \operatorname{rad}(m(m+1))=\operatorname{rad}(n(n+1))\}.
\]
An element of \(\Omega_Y\) is called a source pair.
For each source pair define its four corner gcds by
\[
A=\gcd(m,n),\quad B=\gcd(m,n+1),\quad
C=\gcd(m+1,n),\quad D=\gcd(m+1,n+1).
\]
These four integers are pairwise coprime. For example, a prime dividing
\(A\) and \(B\) would divide both \(n\) and \(n+1\). A prime
dividing \(A\) and \(C\) would divide both \(m\) and \(m+1\). The
same two consecutive pairs settle the other four pairs. Each corner
divides the two entries in its definition. Hence, by pairwise
coprimality, \(AB\mid m\), \(CD\mid m+1\), \(AC\mid n\), and \(BD\mid n+1\).
Define the positive integer residuals
\[
a=\frac m{AB},\qquad b=\frac{m+1}{CD},\qquad
c=\frac n{AC},\qquad d=\frac{n+1}{BD}.
\]
\begin{samepage}
The residuals are pairwise coprime. First, \(\gcd(a,b)=\gcd(c,d)=1\),
since each pair divides consecutive integers. For the four cross pairs,
divide the defining corner gcds. This gives
\[
\gcd(Ba,Cc)=\gcd(Aa,Dd)=\gcd(Db,Ac)=\gcd(Cb,Bd)=1,
\]
so \(\gcd(a,c)=\gcd(a,d)=\gcd(b,c)=\gcd(b,d)=1\) as well.
\end{samepage}
The coordinates and mask are
\[
\begin{gathered}
m=ABa,\quad m+1=CDb,\quad n=ACc,\quad n+1=BDd,\\
Z=ABCD,\quad R=\operatorname{rad}(Z),\quad M=abcd,\\
K_{\rm res}=M/\operatorname{rad}(M),\quad L=cd-ab,\\
\alpha=ac,\quad\beta=bd,\quad\gamma=ad,\quad\delta=bc,\\
u=(m+n+1)/(BC),\qquad t=(n-m)/(AD),\\
CDb-ABa=1,\qquad BDd-ACc=1.
\end{gathered}\AutoTag{eq:old:2_1}
\]
The four cross-gcd identities above will be used below. The radical
condition is exactly \(\operatorname{rad}(M)\mid Z\). If \(p\mid a\), then
\(p\mid m\). Equality of the original radicals forces \(p\mid n\) or
\(p\mid n+1\). Hence \(p\mid A\) or \(p\mid B\). The other residuals give
the corresponding corner. Conversely, every prime of \(Z\) divides both
consecutive products. The mask ensures that no residual introduces a
prime outside \(Z\).
Thus the common radical is
\[
\operatorname{rad}(m(m+1))=\operatorname{rad}(n(n+1))=R.
\]
The identities
\[
Z^2M=m(m+1)n(n+1),\quad ZL=(n-m)(m+n+1)
\]
give \(ut=L>0\) once the shifts are integral. Their integrality follows
from the corner identities
\[
m+n+1=C(Ac+Db)=B(Aa+Dd),\qquad
n-m=A(Cc-Ba)=D(Bd-Cb).
\]
Since \(\gcd(B,C)=\gcd(A,D)=1\), the first numerator is divisible by
\(BC\) and the second by \(AD\). Thus \(u,t\) are positive integers
with product \(L\).
The shifts are units at the opposite corner pairs:
\[
\gcd(u,AD)=\gcd(t,BC)=1.
\]
Modulo a prime dividing \(A\), the integer \(m+n+1\) is
congruent to \(1\). Modulo a prime dividing \(D\), it is congruent to
\(-1\). Modulo a prime dividing \(B\), \(n-m\) is congruent to
\(-1\). Modulo a prime dividing \(C\), it is congruent to \(1\).
The denominator factors \(BC\) and \(AD\) are
coprime to these opposite corner pairs.

The two conics, scale identities, and inverse maps from the same
parametrization are
\[
\begin{aligned}
\beta D^2-\alpha A^2&=u,&
\delta C^2-\gamma B^2&=t,\\
\alpha A^2&=u\frac{mn}{m+n+1},&
\beta D^2&=u\frac{(m+1)(n+1)}{m+n+1},\\
\gamma B^2&=t\frac{m(n+1)}{n-m},&
\delta C^2&=t\frac{n(m+1)}{n-m},\\
B&=(cA+bD)/u,& C&=(aA+dD)/u,\\
A&=(dB-bC)/t,& D&=(cC-aB)/t.
\end{aligned}\AutoTag{eq:old:2_2}
\]
We use only source points. They satisfy positivity, integrality, the
four corner gcds and the cross gcds. They also satisfy the shell
condition and the radical mask. We assert no converse for arbitrary
positive solutions of either conic.

The selected chart is deterministic:
\[
(J,K;s)=
\begin{cases}(A,D;u),&u\le t,\\(B,C;t),&t<u.\end{cases}\AutoTag{eq:old:2_3}
\]
Thus \(u=t\) belongs to the first chart. Write \((f,g)=(\alpha,\beta)\)
in the first chart and \((f,g)=(\gamma,\delta)\) in the second. For a
positive integer \(x\), let \(v_p(x)\) be the exponent of \(p\) in its
prime factorization. For a finite set of primes \(\mathcal P\), put
\[
x_{\mathcal P}=\prod_{p\in\mathcal P}p^{v_p(x)},\qquad
x_{\varnothing}=1.
\]
Then
\[
f_J=f_{\{p:p\mid J\}},\quad g_K=g_{\{p:p\mid K\}},\quad
v=f/f_J,\quad w=g/g_K,\quad
r_J=\operatorname{rad}(f_J),\quad r_K=\operatorname{rad}(g_K),
\]
\[
\Pi=\min(r_J,v)\min(r_K,w).\AutoTag{eq:old:2_4}
\]
These supported parts are functions of the canonical source. They are
not extra summation variables.

Set
\[
\begin{gathered}
S_0=\frac{30869}{5000000},\quad
\epsilon_0=k_0=S_0,\quad l_0=2S_0,\quad j_0=p_0=3S_0,\\
a_0=\frac6{125},\quad b_0=\frac{4166537}{30000000},\quad
\Delta_{\rm exp}=3S_0,\\
q_{\min}=\frac{9907393}{15000000}=\frac23-S_0,\quad
q_{\max}=\frac{5092607}{7500000}=\frac23+2S_0,\quad
\tau_0=1-\frac32S_0 .
\end{gathered}\AutoTag{eq:old:2_5}
\]
\begin{samepage}
We use the following notation.
\begin{center}
\small
\begin{tabular}{@{}ll@{}}
Symbol & Meaning\\
\hline
\(Y\) & dyadic shell endpoint; \(Y/2<n\le Y\)\\
\(A,B,C,D\) & four corner gcds\\
\(a,b,c,d\) & positive residual factors\\
\(Z,R,M,L\) & \(ABCD,\ \operatorname{rad}(Z),\ abcd,\ cd-ab\)\\
\(\alpha,\beta,\gamma,\delta\) & \(ac,\ bd,\ ad,\ bc\)\\
\(u,t\) & positive shifts, with \(ut=L\)\\
\(J,K;s\) & selected corner pair and its smaller shift\\
\(f_J,g_K,v,w,\Pi\) & supported coefficients, complements, and their mask\\
\(S_0,g_0,\beta_0\) & minimum saving, reserve, and allocated loss (values in the Introduction)\\
\end{tabular}
\end{center}
\end{samepage}
Define the high sector
\[
\mathcal C_{\rm hi}=\{(m,n)\in\Omega_Y:
R>Y^{5/3-\epsilon_0},\
M>Y^{q_{\min}},\
\operatorname{rad}(M)>Y^{2/3-k_0}\}.
\AutoTag{eq:old:2_6}
\]
The high sector is the complement of the first three predicates in
\eqref{eq:old:3_1}. Those rows test the complementary inequalities.
For sources with \(R>Y^{5/3-\epsilon_0}\), the source identities give
\[
M<\kappa_Y^2Y^{q_{\max}},\qquad
L<\kappa_YY^{1/3+\epsilon_0},\qquad \kappa_Y=1+Y^{-1}.
\AutoTag{eq:old:2_7}
\]
We have \(m(m+1)n(n+1)\le\kappa_Y^2Y^4\) and
\(ZL=(n-m)(m+n+1)<\kappa_YY^2\). Also
\(Z>Y^{5/3-\epsilon_0}\). On \(\mathcal C_{\rm hi}\), the strict lower
cuts on \(R\) and \(M\) also give
\[
Y^{4-3\epsilon_0}<Z^2M=m(m+1)n(n+1)
                       \le2\kappa_Ym^2Y^2,
\]
so \(m^2>Y^{2-3\epsilon_0}/(2\kappa_Y)
\ge Y^{2-3\epsilon_0}/3\). Therefore
\[
m>Y^{\tau_0}/\sqrt3.\AutoTag{eq:old:2_8}
\]
Here \(2\kappa_Y\le3\) for \(Y\ge2\). This last inequality does not
use the radical cut in \(\mathcal C_{\rm hi}\).

After the \(Q0\) row of \eqref{eq:old:3_1} has failed,
\(K_{\rm res}=M/\operatorname{rad}(M)
\le Y^{\Delta_{\rm exp}}\). After \(Q1\) has failed,
\(L\ge Y^{1/3-l_0}\). After both \(Q2\) chart tests have failed, the
selected corners satisfy \(J,K\ge Y^{1/3-j_0}\). These are the
non-strict complements used in the following partition.
Display~\eqref{eq:old:3_1} states the predicates and their order.

\section{The first-match partition}\label{sec:partition}

Each predicate below is tested only after every earlier row has failed.
Here \(s=\min(u,t)\) in the chart selected by \eqref{eq:old:2_3}.
\[
\begin{array}{c|l}
\text{row}&\text{predicate}\\ \hline
O_R&R\le Y^{5/3-\epsilon_0}\\
O_M&M\le Y^{2/3-\epsilon_0}\\
O_{\operatorname{rad}M}&\operatorname{rad}(M)\le Y^{2/3-k_0}\\
Q0&K_{\rm res}>Y^{\Delta_{\rm exp}}\\
Q1&L<Y^{1/3-l_0}\\
Q2u&u\le t,\ \min(A,D)<Y^{1/3-j_0}\\
Q2t&t<u,\ \min(B,C)<Y^{1/3-j_0}\\
3u&u\le t,\ u\le Y^{a_0}\\
3t&t<u,\ t\le Y^{a_0}\\
4E&Y^{a_0}<s<Y^{b_0},\ \min(v,w)\le Y^{p_0}\\
4H&Y^{a_0}<s<Y^{b_0},\ v,w>Y^{p_0},\ \Pi\le Y^{p_0}\\
4P&Y^{a_0}<s<Y^{b_0},\ v,w>Y^{p_0},\ \Pi>Y^{p_0}\\
5W&\text{the remaining source.}
\end{array}\AutoTag{eq:old:3_1}
\]

\begin{theorem}[Exact partition]\label{thm:partition}
The thirteen first-match sets in \eqref{eq:old:3_1} are pairwise disjoint and have union
\(\Omega_Y\). On \(5W\),
\[
s\ge Y^{b_0},\qquad K_{\rm res}\le Y^{\Delta_{\rm exp}},\qquad
L\ge Y^{1/3-l_0},\qquad J,K\ge Y^{1/3-j_0}.
\AutoTag{eq:old:3_2}
\]
\end{theorem}

\begin{proof}
The quantities in \eqref{eq:old:3_1} are deterministic functions of the canonical
source. Each of the first twelve rows is intersected with the complements of
the preceding predicates. So these rows are disjoint. Their complement is
exactly the last row, which proves exhaustiveness. Failure of \(Q0\) gives
\(K_{\rm res}\le Y^{\Delta_{\rm exp}}\). Failure of \(Q1\) gives
\(L\ge Y^{1/3-l_0}\). On the complement of both chart-specific \(Q2\)
tests, the selected corners satisfy \(J,K\ge Y^{1/3-j_0}\).
Failure of the \(3u\) and \(3t\) tests gives \(s>Y^{a_0}\).
If \(s<Y^{b_0}\), trichotomy gives exactly one of the \(4E,4H,4P\)
predicates. Therefore their complement has \(s\ge Y^{b_0}\).
\end{proof}

The following boundary conventions apply after all earlier rows have
failed. The first three rows use non-strict upper cuts. Equality at
\(K_{\rm res}=Y^{\Delta_{\rm exp}}\),
\(L=Y^{1/3-l_0}\), and \(\min(J,K)=Y^{1/3-j_0}\) survives the
corresponding strict test. The chart equality \(u=t\) selects the
\(u\)-chart. The full first-match order determines its row. The
\(3u,3t\) tests include \(s=Y^{a_0}\), while \(5W\) contains
\(s=Y^{b_0}\). In the open middle interval, \(4E\) contains the boundary
\(\min(v,w)=Y^{p_0}\). Also \(4H\) contains \(\Pi=Y^{p_0}\) when
\(v,w>Y^{p_0}\).

\begin{remark}[Counting populations]
The estimates below count the following populations:
\begin{enumerate}
\item For \(O_R\), fix \(n\) first. Then apply the
      support-fibre lemma in \cite{Paper1} to the partners \(m\).
\item \(O_M,O_{\operatorname{rad}M},Q0\) use the complete fixed-\(M\)
      quadratic fibre, including all ordered factorizations
      \(M=(ab)(cd)\). \(Q1\) instead fixes \((P,L)\), with
      \(P=ab\) and \(Q_{\rm prod}=cd=P+L\).
\item \(Q2u,Q2t\) are disjoint chart populations. Their two analytic
      contributions apiece are additive bounds for each population.
\item \(3u,3t,4E,4H,4P\) retain the corresponding earlier-failure
      predicates and the canonical inverse filters. The \(4P\) estimate is
      not applied to the \(5W\) complement.
\item \(5W\) is exactly the remaining first-match set in \eqref{eq:old:3_1}.
\end{enumerate}
\end{remark}

Fix the ordered residual factors \((a,b,c,d)\), both shifts \((u,t)\),
the selected chart and its conic coordinates. Then the inverse identities
\eqref{eq:old:2_2} reconstruct the remaining two corners. This is the
full canonical label used for inverse reconstruction. The outer-row
divisor bound includes ordered residual factorizations. The joint-row sums
include their coefficient factorizations. The later estimates retain the
stated divisor or factor multiplicities. A larger set used to estimate one
row may intersect the set used for another row. The first-match partition
makes the original populations disjoint. Each row is bounded separately.

\section{Outer rows and inherited fibres}\label{sec:outer}

The small-radical input is Theorem~3 of \cite{Paper1}.
Appendix~\ref{app:sources} states it, and Appendix~\ref{app:lowrad}
records the shell substitution. With \(H=Y/2\) and
\(V=Y^{5/3-\epsilon_0}\), it gives
\[
\#\{n:Y/2<n\le Y,\ \operatorname{rad}(n(n+1))\le
Y^{5/3-\epsilon_0}\}
\ll_\eta Y^{2/3-S_0+\eta}.
\AutoTag{eq:old:4_2}
\]
The threshold in that application is large enough that
\(Y\ge4\) and \(Y^{1/3+\epsilon_0}>4\). These are exactly the required
conditions \(H\ge2\) and \(V<H^2\).

Let \(Q_{\rm sf}\) be a positive squarefree integer. For a positive
integer \(Q\), let \(\omega(Q)\) denote the number of distinct prime
divisors of \(Q\). Apply the support-fibre estimate
\cite[Lemma~2.1]{Paper1} with \(X=n\):
\[
\#\{1\le m<n:\operatorname{rad}(m(m+1))\mid Q_{\rm sf}\}
 \ll_\eta Y^{2\eta}\log(2Y).
\AutoTag{eq:old:4_3}
\]
Here \(Q_{\rm sf}=\operatorname{rad}(n(n+1))\) is squarefree and at most
\(Y(Y+1)\). The source equality of radicals gives the divisibility
condition for every partner. The bound follows from
\(2^{\omega(Q_{\rm sf})}\ll_\eta Q_{\rm sf}^\eta\). Choose the
internal exponents smaller than the requested final \(\varepsilon\).
Absorbing the logarithm then proves the \(O_R\) row bound.

\begin{proposition}[Five outer rows]\label{prop:outerrows}
For every \(\varepsilon>0\),
\[
\begin{aligned}
\#O_R,\ \#O_M,\ \#O_{\operatorname{rad}M},\ \#Q1
  &\ll_\varepsilon Y^{2/3-S_0+\varepsilon},\\
\#Q0&\ll_\varepsilon Y^\varepsilon.
\end{aligned}\AutoTag{eq:old:4_4}
\]
\end{proposition}

\begin{proof}
The \(O_R\) estimate is \eqref{eq:old:4_2} followed by the partner-fibre bound \eqref{eq:old:4_3}.
We apply the latter with \(X=n\) and with the squarefree common radical
as \(Q\). The defining equality of radicals implies the required
divisibility \(\operatorname{rad}(m(m+1))\mid Q\). Each logarithm and divisor factor uses a
separate portion of the arbitrary \(\varepsilon\).

For \(O_M\), fix \(M\le Y^{2/3-\epsilon_0}\). Write
\(P=ab\), \(Q_{\rm prod}=cd\), so \(PQ_{\rm prod}=M\),
\(P<Q_{\rm prod}\), and \(Q_{\rm prod}-P=L>0\).
There are at most \(d_4(M)\ll_\eta M^\eta\) ordered choices of
\((a,b,c,d)\), where
\[
d_4(M)=\#\{(a,b,c,d)\in\mathbb Z_{>0}^4:abcd=M\}.
\]
For each, the source identity
\[
Q_{\rm prod}m(m+1)=Pn(n+1)
\]
is equivalent, after setting \(x=2n+1\) and \(y=2m+1\), to
\[
P x^2+(Q_{\rm prod}-P)=Q_{\rm prod}y^2.
\AutoTag{eq:old:4_5}
\]
Appendix~\ref{app:fibre} bounds the number of such pairs by \(Y^\eta\).
The coefficients and nonzero shift have height bounded by a fixed power
of \(Y\). For the theorem's variable bound take \(x_0=2Y+1\), since
\(2n+1\le x_0\). The source identity gives
\(M\le Y^2(Y+1)^2\le4Y^4\), so
\(P,Q_{\rm prod},|Q_{\rm prod}-P|=L\le4Y^4\le x_0^4\), and the
fixed height exponent is \(A_{\rm ht}=4\).
Thus the implied constant is uniform in the moving labels. Summing over
\(M\) proves the \(O_M\) bound.

For \(O_{\operatorname{rad}M}\), \eqref{eq:old:2_7} gives \(M\le4Y^4\). If
\(R_0=Y^{2/3-k_0}\), then
\[
\#\{M\le4Y^4:\operatorname{rad}(M)\le R_0\}
\le R_0(4Y^4)^\sigma
 \sum_{M\ge1}\frac1{\operatorname{rad}(M)M^\sigma}.
\AutoTag{eq:old:4_6}
\]
For fixed \(\sigma>0\), the series is the convergent product
\[
\prod_q\left(1+\frac1{q(q^\sigma-1)}\right).
\]
Fix the final target loss. Then choose a fixed \(\sigma>0\) so small
that \(4\sigma\) is below its allocated part. Choose the divisor and
fibre exponents independently, each within its own remaining allocation.
The resulting label count is
\(O_\sigma(Y^{2/3-k_0+4\sigma})\).
For each \(M\), sum the at most \(d_4(M)\) residual factorizations and
apply the same fixed-height fibre \eqref{eq:old:4_5}. Since \(k_0=S_0\), this gives
the asserted exponent.

On \(Q0\), failure of the first three rows and the \(Q0\) predicate give
\[
Y^{q_{\max}}<M<\kappa_Y^2Y^{q_{\max}}.
\]
The interval length is at most
\((\kappa_Y^2-1)Y^{q_{\max}}\le3Y^{q_{\max}-1}<1\) for \(Y\ge256\).
Thus at most one integer \(M\) occurs. Its complete ordered
factorization sum and the fibre \eqref{eq:old:4_5} cost \(Y^\varepsilon\). We do
not declare the row empty.

On \(Q1\), enumerate \(P=ab\) and \(L=cd-ab\), and put
\(Q_{\rm prod}=cd=P+L\). Since \(M=PQ_{\rm prod}\ge P^2\),
\eqref{eq:old:2_7} gives \(P<\kappa_YY^{q_{\max}/2}\ll Y^{1/3+S_0}\).
The \(Q1\) predicate and the positivity of the integer \(L\) give
\(1\le L<Y^{1/3-l_0}=Y^{1/3-2S_0}\). Hence the number of \((P,L)\)
labels is
\[
\ll Y^{1/3+S_0}Y^{1/3-2S_0}=Y^{2/3-S_0}.
\]
Each label fixes \(Q_{\rm prod}=P+L\). Splitting \(P,Q_{\rm prod}\)
into their ordered factors costs at most
\(\tauDiv(P)\tauDiv(Q_{\rm prod})\ll_\eta Y^\eta\). Then the fixed-height fibre
\eqref{eq:old:4_5} is used once. This row carries no additional divisor sum
over \(M\).
\end{proof}

\begin{remark}
The source and theorem hypotheses in this section hold above fixed
thresholds. Their applications are recorded in
Appendix~\ref{app:lowrad}.
The finitely many remaining \(Y\) are absorbed into the implied constants
in Proposition~\ref{prop:outerrows}.
\end{remark}

\section{Joint-corner and short-shift rows}\label{sec:joint}

The two \(Q2\) populations are disjoint. This is because the chart rule
\eqref{eq:old:2_3} assigns each source to exactly one of \(u\le t\) or
\(t<u\). In a dyadic sector for the selected corners, write
\[
J\asymp J_0,\qquad K\asymp K_0,\qquad
 P_0=J_0K_0,\qquad \mu_{\rm sec}=\min(J_0,K_0).
\]
Here each dyadic sector means \(J_0\le J<2J_0\) and
\(K_0\le K<2K_0\). These half-open sectors are disjoint. The \(Q2\) predicate gives
\(\mu_{\rm sec}<Y^{1/3-j_0}\).

\begin{lemma}[Sector count]\label{lem:q2-sector}
Fix a positive selected shift \(s\) and a nonempty sector as above. Count
the canonical-source coefficient-corner triples with
\(\gcd(f,K)=\gcd(g,J)=\gcd(s,JK)=1\), with the second coefficient
forced by the conic. Suppose that
\[
fJ^2\le H,\qquad gK^2\le H.
\]
Then their number is
\[
\ll_\eta Y^\eta H\left(\frac1{P_0}+\frac{\mu_{\rm sec}^3}{P_0^2}\right)
\]
for every \(\eta>0\), with an implied constant depending only on \(\eta\).
In the \(Q2u\) application take \(H=uY\). In the \(Q2t\) application
we assert the natural cutoff \(H\ll tY^{1+\epsilon_0+l_0}\) only on
the first-match population. There the earlier \(O_R\) and \(Q1\) rows
have failed. That source-uniform bound is proved in \eqref{eq:old:5_2}.
\end{lemma}

\begin{proof}
In the \(u\) chart the conic is \(gK^2-fJ^2=u\). The imposed cutoff
gives \(fJ^2,gK^2\le H\). For the \(Q2u\) application, the source scale
identities permit the choice \(H=uY\).
The full source filter gives \(\gcd(f,K)=\gcd(g,J)=1\) and
\(\gcd(u,JK)=1\). These are the unit conditions for either orientation.
Thus the quadratic congruence has at most \(2^{\omega(K)+2}\) roots
modulo \(K^2\), or \(2^{\omega(J)+2}\) roots modulo \(J^2\).
For odd prime powers each unit-square congruence has at most two roots.
The prime \(2\) contributes at most four roots modulo a power of \(2\).
So the factor \(2^{\omega(\cdot)+2}\) is a uniform
overcount, including the even modulus.
Suppose first \(J_0\ge K_0\). For fixed \(f\) and \(K\),
\[
fJ^2\equiv-u\pmod{K^2}.
\]
This congruence has at most \(2^{\omega(K)+2}\ll_\eta Y^\eta\)
root classes modulo \(K^2\). Each class contains at most
\(J_0/K^2+1\) integers \(J\) in the sector. There are
\(O(K_0)\) choices of \(K\) and \(O(H/J_0^2)\) choices of \(f\).
The resulting count is
\[
\ll_\eta Y^\eta H
\left(\frac1{J_0K_0}+\frac{K_0}{J_0^2}\right)
=Y^\eta H\left(\frac1{P_0}+\frac{\mu_{\rm sec}^3}{P_0^2}\right).
\]
If \(K_0\ge J_0\), reduce the conic modulo \(J^2\). This gives the
same expression with \(J_0,K_0\) interchanged.

For the \(t\) chart the full filter gives
\(\gcd(g,J)=\gcd(f,K)=1\) and \(\gcd(t,JK)=1\). Use
\(gK^2-fJ^2=t\), now with \((J,K)=(B,C)\).
Combine the scale identities in \eqref{eq:old:2_2}, the source bounds
\eqref{eq:old:2_7} and the failed \(Q1\) inequality. They give
\[
fJ^2,\ gK^2\ll tY^{1+\epsilon_0+l_0}.
\AutoTag{eq:old:5_2}
\]
The source identity \(Zcd=n(n+1)\) and \(Z\ge R>Y^{5/3-\epsilon_0}\)
give
\[
cd=\frac{n(n+1)}Z<\kappa_Y Y^{1/3+\epsilon_0}.
\]
The denominator satisfies
\[
cd=ab+L\ll Y^{1/3+\epsilon_0},\quad
n-m=\frac{L\,n(n+1)}{cd(m+n+1)}
\gg Y^{1-\epsilon_0-l_0}.
\]
The second estimate uses \(n>Y/2\), \(L\ge Y^{1/3-l_0}\),
\(m+n+1\le2Y\), and the first estimate. Substitution in the last two
scale identities of \eqref{eq:old:2_2} proves \eqref{eq:old:5_2}. The same
congruence argument applies, using \(\delta C^2\equiv t\pmod{B^2}\) or
\(\gamma B^2\equiv -t\pmod{C^2}\). It gives the asserted sector bound.
\end{proof}

\begin{proposition}[Joint-corner rows]\label{prop:q2}
For every \(\varepsilon>0\),
\[
\begin{aligned}
\#Q2u&\ll_\varepsilon
Y^{2/3-2222179/15000000+\varepsilon}
+Y^{2/3-30869/1250000+\varepsilon},\\
\#Q2t&\ll_\varepsilon
Y^{2/3-972179/7500000+\varepsilon}
+Y^{2/3-S_0+\varepsilon}.
\end{aligned}\AutoTag{eq:old:5_3}
\]
The terms in each line are additive contributions to one chart
population. They are not separate first-match rows.
\end{proposition}

\begin{proof}
On \(Q2u\) or \(Q2t\), \eqref{eq:old:2_7} and the failed earlier tests give
\[
\begin{gathered}
Y^{q_{\min}}<M<\kappa_Y^2Y^{q_{\max}},\qquad
Y^{1/3-l_0}\le L<\kappa_YY^{1/3+\epsilon_0},\\
1\le s<\sqrt{\kappa_Y}\,Y^{1/6+\epsilon_0/2}.
\end{gathered}\AutoTag{eq:old:5_4}
\]
The last inequality follows from \(ut=L\) and
\(s=\min(u,t)\le\sqrt L\).

We also need a lower bound for the product of the selected corners.
Again \(Zcd=n(n+1)\) and \(Z\ge R>Y^{5/3-\epsilon_0}\), so
\(cd<\kappa_Y Y^{1/3+\epsilon_0}\). The identity
\[
\frac{n-m}{L}=\frac{n(n+1)}{cd(m+n+1)}
\]
and \(n>Y/2\) imply
\[
AD=u\frac{n-m}{L}\gg uY^{2/3-\epsilon_0}.
\]
Reducing the exponent by \(\tfrac32\epsilon_0\) absorbs the fixed
constant for all sufficiently large \(Y\). A sector has
\(AD<4P_0\). So its nonemptiness implies
\[
P_0>\frac{uY^{2/3-5\epsilon_0/2}}{32}
\quad\text{in the \(u\) chart.}\AutoTag{eq:old:5_5}
\]
In the \(t\) chart, \(BC=t(m+n+1)/L\), so \eqref{eq:old:5_4} gives
\[
BC\gg tY^{2/3-\epsilon_0},\qquad
P_0>\frac{tY^{2/3-5\epsilon_0/2}}{32}.
\AutoTag{eq:old:5_6}
\]
We reduce the exponent by \(3\epsilon_0/2\) from the scale bound
\(BC\gg tY^{2/3-\epsilon_0}\). This absorbs fixed dyadic constants uniformly.
This gives the exponent used in \eqref{eq:old:5_8}.

Apply Lemma~\ref{lem:q2-sector} in each occupied sector. In the \(u\)
chart, its two terms are bounded by \eqref{eq:old:5_5} and
\(\mu_{\rm sec}<Y^{1/3-j_0}\) as
\[
Y^\eta\left(
Y^{1/3+5\epsilon_0/2}
+\frac{Y^{2/3-3j_0+5\epsilon_0}}u
\right).
\AutoTag{eq:old:5_7}
\]
For the \(t\) chart, \eqref{eq:old:5_2}, \eqref{eq:old:5_6}, and the same corner bound give
\[
Y^\eta\left(
Y^{1/3+7\epsilon_0/2+l_0}
+\frac{Y^{2/3-3j_0+6\epsilon_0+l_0}}t
\right).
\AutoTag{eq:old:5_8}
\]
The source has \(s\le\sqrt{\kappa_Y}Y^{1/6+\epsilon_0/2}\).
Sum the first terms of \eqref{eq:old:5_7}--\eqref{eq:old:5_8} over this many
shifts. This adds \(1/6+\epsilon_0/2\) to their exponents. The second terms are harmonic
in the selected shift and add only \(O(\log Y)\). There are
\(O(\log^2Y)\) corner sectors.

Fix a coefficient \(f\), a sector corner pair and a shift. Then the
conic forces the other coefficient \(g\). The number of ordered residual
splits of \(f,g\) is at most \(\tauDiv(f)\tauDiv(g)\ll_\eta Y^\eta\).
The inverse maps \eqref{eq:old:2_2} then determine the remaining corners. The full
canonical filter is retained. Choose each internal exponent first. Then
the divisor factors, root factors, dyadic logarithms and the harmonic sum
use a total arbitrarily small loss.

Substitute \(\epsilon_0=k_0=S_0\), \(l_0=2S_0\) and \(j_0=3S_0\).
This gives the four exponents
\[
\begin{array}{ll}
1/2+3S_0=2/3-2222179/15000000,&
2/3-9S_0+5S_0=2/3-30869/1250000,\\
1/2+4S_0+2S_0=2/3-972179/7500000,&
2/3-9S_0+6S_0+2S_0=2/3-S_0.
\end{array}
\AutoTag{eq:old:5_9}
\]
This proves \eqref{eq:old:5_3}.
\end{proof}

\subsection*{A determinant estimate for the short-shift rows}

Write
\[
 e=S_0=\epsilon_0=k_0=\frac{30869}{5000000},\quad
 a_0=\frac6{125},\quad \nu=\frac1{64},\quad \xi=\frac1{10000},
\]
\[
 q_-=\frac23-e,\quad q_+=\frac23+2e,\quad
 \tau_0=1-\frac32e .
\AutoTag{eq:old:5_10}
\]
Thus \(q_-=9907393/15000000\), \(q_+=5092607/7500000\), and
\(\tau_0=9907393/10000000\). The high-source inequalities give
\[
 m>\frac{Y^{\tau_0}}{\sqrt3},\qquad
 Y^{q_-}<M<\kappa_Y^2Y^{q_+},\qquad \kappa_Y=1+Y^{-1}.
\AutoTag{eq:old:5_11}
\]
The first inequality follows from \eqref{eq:old:2_8}. The lower bound for \(M\) follows
from failure of the \(O_M\) predicate in the first-match partition
\eqref{eq:old:3_1}. Its upper bound follows from failure of \(O_R\) and
\eqref{eq:old:2_7}.

We first give the determinant estimate for each fixed shift and
coefficient box. Suppose \(r_0,s_0\ge1\),
\(K_{\rm det}=(r_0+1)(s_0+1)\), and
\(K_{\rm det}\) source points lie in one cell. In both applications below, scale
the two projective coordinates by the dyadic box centers. Then a determinant
of the \(K_{\rm det}\) monomial rows has the form
\[
 |\Delta_{\rm det}|\le C_{r_0,s_0}(H_1H_2)^{r_0K_{\rm det}/2}(EF)^{s_0K_{\rm det}/2}
 \max_{\substack{n_\ell\ge0\\\sum_\ell n_\ell=K_{\rm det}}}
 \vartheta^{\sum_\ell \ell n_\ell}
 \lambda^{\sum_\ell n_\ell(n_\ell-1)/2}.
\AutoTag{eq:old:5_12}
\]
Here \(E,F\) are coefficient scales and \(H_1,H_2\) are corner scales.
The cell width in normalized coordinates is
\(\lambda\), and \(\vartheta\) bounds the thickness of the source curve.
In the \(u\)-chart, the remaining functions in a level-\(\ell\) minor are
\(x^{i-2j}\), \(x\in[1/2,2]\). Here
\(\vartheta\) bounds \(\theta\) in \((1-\theta)^j\). In the \(t\)-chart,
put
\[
 \sigma_t=H_B/H_C,\qquad x=(B/C)/\sigma_t,\qquad
 \kappa_t=\sigma_t^2E/F,\qquad y=(E/F)t/(\gamma C^2).
\]
Then \(x\in[1/2,2]\), \(1/16<\kappa_t<8\), and the normalized
coefficient ratio is
\[
 (E/F)(\delta/\gamma)=\kappa_t x^2+y.
\]
The level-\(\ell\) functions are therefore
\[
 x^i(\kappa_t x^2+y)^j
 =\sum_{\ell=0}^j\binom j\ell\kappa_t^{j-\ell}y^\ell
     x^{i+2j-2\ell}.
\]
Fix the determinant degree. The derivatives in \(x\) of these
finitely many monomials are then uniformly bounded, since \(x\in[1/2,2]\)
and \(\kappa_t\in[1/16,8]\). The thickness is the normalized term \(y\).
It is comparable to \(\Xi_t^{-1}\) in the \(3t\) application below.
This normalization gives the same level and cell-width gains in \eqref{eq:old:5_12}.
To verify \eqref{eq:old:5_12}, expand each monomial column as a sum of
level columns and use determinant multilinearity in the columns. In each
resulting determinant, group columns with the same level and apply Laplace
expansion over row subsets of the corresponding sizes. In a level-\(\ell\)
minor, factor the \(\ell\)-th power of the thickness coordinate from each
evaluation row. The resulting thickness exponent is \(\sum_\ell \ell n_\ell\).
If two evaluation points or two powers agree,
its remaining determinant is zero. Otherwise take \(n\) evaluation points
in an interval of length \(\lambda\). Take \(n\) functions with bounded
derivatives through order \(n-1\). Successive divided differences factor
out the Vandermonde product. The remaining determinant is at most a
constant times the derivative bounds. This gives
\(O(\lambda^{n(n-1)/2})\). Thus this group contributes
\(C_{r_0,s_0}\lambda^{n_\ell(n_\ell-1)/2}\). Here all derivatives of
the finitely many powers are bounded on the fixed compact interval. Row
and column scaling gives the first factor in \eqref{eq:old:5_12}.

Put \(b_{\rm cell}=\log_Y(1/\lambda)\) and
\(\mu_{\rm th}=\log_Y(1/\vartheta)\). In the source boxes these exact
quantities equal their limiting exponents up to \(O(1/\log Y)\). Then
\[
\sum_{\ell\ge0}\left(\mu_{\rm th}\ell n_\ell+
 \frac {b_{\rm cell}}2 n_\ell(n_\ell-1)\right)
\ge \frac{2\sqrt2}{3}\sqrt{b_{\rm cell}\mu_{\rm th}}\,K_{\rm det}^{3/2}
-\left(\mu_{\rm th}+\frac {b_{\rm cell}}2\right)K_{\rm det}.
\AutoTag{eq:old:5_13}
\]
For every \(\Lambda>0\), termwise completion of the square gives
\[
\sum_\ell\left(\mu_{\rm th}\ell n_\ell+\frac {b_{\rm cell}}2n_\ell^2\right)
\ge \Lambda K_{\rm det}-\frac1{2b_{\rm cell}}\sum_{\ell\ge0}(\Lambda-\mu_{\rm th}\ell)_+^2
\ge \Lambda K_{\rm det}-\frac{\Lambda^2}{2b_{\rm cell}}-\frac{\Lambda^3}{6b_{\rm cell}\mu_{\rm th}}.
\]
Taking \(\Lambda=(2b_{\rm cell}\mu_{\rm th}K_{\rm det})^{1/2}\) gives the leading
term in \eqref{eq:old:5_13}. The square-completion loss is
\(\Lambda^2/(2b_{\rm cell})=\mu_{\rm th}K_{\rm det}\). Replacing \(n_\ell^2\) by
\(n_\ell(n_\ell-1)\) costs a separate \((b_{\rm cell}/2)K_{\rm det}\).
For \(r_0=\lceil s_0 q/p\rceil\) or \(r_0=\lfloor s_0 q/p\rfloor\),
\[
\frac{r_0p+s_0q}{2}\le s_0q+\frac p2,\qquad
K_{\rm det}^{1/2}\ge (r_0s_0)^{1/2}.
\]
Here \(p\) and \(q\) are the logarithmic scale exponents of the corner
product and coefficient product in the chosen chart. The \(3u\) and \(3t\)
applications below specify their normalizations. Hence a fixed
positive margin in
\[
b_{\rm cell}\mu_{\rm th}>\frac98pq
\AutoTag{eq:old:5_14}
\]
makes every integerized determinant have absolute value less than \(1\).
Here \(p,q,\mu_{\rm th}\) range over fixed compact intervals in both
applications. The strict margin is uniform over the full source boxes.
Fix that margin and those compact intervals, and then choose \(s_0\).
The degree-rounding and lower-order terms are then smaller than the margin.
This one fixed degree works for every sufficiently large \(Y\).
The integerized matrix has rank below \(K_{\rm det}\). Let \(r<K_{\rm det}\)
be the largest order of a nonzero minor. If \(r=0\), a coordinate vector
is a kernel vector. If \(r>0\), choose a maximal nonzero \(r\times r\)
minor and one additional column. Cramer's rule gives a nonzero integral
kernel vector. Its coordinates are minors of order at most \(r\). The selected rows
span the row space, so this vector annihilates every row. Thus the kernel
vector has polynomially bounded height and gives a nonzero integral
auxiliary polynomial.
If a cell contains fewer than \(K_{\rm det}\) points, those points are counted
directly. The degrees are fixed. The coefficient height of the
polynomial is a fixed power of \(Y\). This holds because it is a minor of
a matrix whose entries have fixed polynomial height.

Heath--Brown's Lemma~2~\cite[Lemma~2, p.~8]{HB2012}, stated in
Appendix~\ref{app:sources}, has no size restriction on the second pair.
In \(3u\) and \(3t\), the
primitive coefficient pair is the first pair and the primitive corner
pair is the second. For an irreducible factor, let \(a_c\) be its
degree in the corner pair. Let \(b_q\) be its degree in the coefficient
pair. In the theorem's notation this is \((a,b)=(b_q,a_c)\). Thus a
factor with \(a_c\ge2\) has second-pair degree at least two.
Any application of Heath--Brown below has \(b_q\ge1\). So both
bidegrees are positive. The theorem's exponent is \(2/a_c\le1\). Here
\(T_1+1\) and the factor height are bounded by fixed
powers of \(Y\). Choose the theorem's \(\varepsilon\) after those height
exponents. This gives \(O_\eta(Y^\eta(T_1+1))\).

Let a bounded-degree integer polynomial have height \(Y^C\). Then each
primitive integer factor has height \(Y^{O_C(1)}\). For the proof, apply
a fixed Kronecker substitution with base larger than the total degree.
Then use Gauss's lemma and the univariate Landau--Mignotte factor bound. The
substitution has no collisions among the bounded-degree monomials.

\subsection*{The short-\(u\) row \(3u\)}

In the \(u\)-chart put \(\alpha=ac\), \(\beta=bd\). The source identities
\eqref{eq:old:2_2} give
\[
 \beta D^2-\alpha A^2=u,\qquad
 \alpha A^2=(1-\theta)\beta D^2,\qquad
 \theta=\frac{m+n+1}{(m+1)(n+1)}<\frac2{m+1}.
\AutoTag{eq:old:5_15}
\]
The residual factors are pairwise coprime, and so are the corner
factors. Hence
\[
 \gcd(\alpha,\beta)=\gcd(A,D)=1.
\AutoTag{eq:old:5_16}
\]
Also
\[
 \frac{um}{3}\le\alpha A^2<um,\qquad
 \frac{um}{3}\le\beta D^2\le u(m+1),
\AutoTag{eq:old:5_17}
\]
because \(n/(m+n+1)\ge1/3\). Multiply these two scale identities
and use \(\alpha\beta=M\). This gives
\[
 AD\ge \frac{um}{3\sqrt M}
 \gg uY^{2/3-5e/2}.
\AutoTag{eq:old:5_18}
\]
The implied constant is absolute for sufficiently large \(Y\). The unit
conditions needed below are also source consequences. If a prime divides
\(A\), it divides \(m,n\). So it divides neither \(b\mid m+1\) nor
\(d\mid n+1\). Hence \(\gcd(\beta,A)=1\). Also \(uBC=m+n+1\) and
\(\gcd(A,BC)=1\), so \(\gcd(u,A)=1\). Therefore
\(\gcd(\beta u,A)=1\). Every corner is at most \(Y+1\), and
\(\alpha,\beta\le u(Y+1)\).

Fix a positive shift \(u\) and a dyadic sector. The counted tuple
is \((\alpha,A,D)\), with \(A,D\) in that sector. The conic determines
the second coefficient \(\beta\), subject to the integrality filter.
Fix \(u\le Y^{a_0}\) and dyadic scales
\[
 A\asymp H_A,\quad D\asymp H_D,\quad
 \alpha\asymp E,\quad\beta\asymp F.
\]
There are \(O((\log Y)^4)\) such boxes. Put
\[
 \chi=\frac{A H_D}{D H_A},\quad
 \mu_b=\log_Y\frac{E H_A^2}{u},\quad
 p=\log_Y(H_AH_D),\quad q=\log_Y(EF),\quad z=\log_Yu.
\]
The box and \eqref{eq:old:5_17} imply, uniformly,
\[
 \tau_0-O(1/\log Y)\le\mu_b\le1+O(1/\log Y),\quad
 q_--O(1/\log Y)\le q\le q_++O(1/\log Y),
\]
\[
 p\le z+\mu_b-\frac q2+O(1/\log Y),\qquad 0\le z\le a_0.
\AutoTag{eq:old:5_19}
\]
They also give
\[
 \frac{q_-}{a_0+2/3+e/2}-O(1/\log Y)
 \le \frac qp \le
 \frac{q_+}{2/3-5e/2}+O(1/\log Y).
\AutoTag{eq:old:5_20}
\]
The exact endpoints in \eqref{eq:old:5_20} are
\(19814786/21532607\) and \(20370428/19536965\), both strictly between
\(0.9\) and \(1.1\).

Set
\[
 B_u=\frac98q_+\left(1+a_0-\frac{q_+}{2}\right)
 =\frac{54121135983551}{100000000000000},\qquad
 N_u=\left\lceil Y^{B_u+\xi}\right\rceil .
\AutoTag{eq:old:5_21}
\]
Let \(f(q,z,\mu)=q(1+(z-q/2)/\mu)\). Its derivatives in \(z,\mu,q\)
have signs \(+,+,+\) on \eqref{eq:old:5_19}: for the last sign,
\[
 1+\frac{z-q}{\mu}\ge
 \frac{9351751}{29722179}+O(1/\log Y)>0,
\]
and for the \(\mu\)-derivative use \(q_-/2-a_0>0\). Hence
\[
 \frac98\frac{pq}{\mu_b}\le B_u+O(1/\log Y).
\AutoTag{eq:old:5_22}
\]
Put \(\delta_0=\xi/4\). From \eqref{eq:old:5_19}, \(\mu_b\ge\tau_0-O(1/\log Y)\),
and \(p,q<1+O(1/\log Y)\). Therefore \eqref{eq:old:5_22} gives the direct margin
\[
(B_u+\xi)\mu_b-\frac98(1+\delta_0)pq
\ge \xi\tau_0-\frac{9\xi}{32}-O(1/\log Y)>\frac{\xi}{2}
\]
for all sufficiently large \(Y\).
The thickness in \eqref{eq:old:5_15} satisfies
\(\theta\le2Y^{-\mu_b}\). The strict margin absorbs the factor
\(2\).

Split \(\chi\in(1/2,2)\) into \(N_u\) half-open intervals of length
\(3/(2N_u)\). For \(K_{\rm det}=(r_0+1)(s_0+1)\) source points in one interval,
form the integer matrix whose \(k,(i,j)\)-entry is
\[
 A_k^iD_k^{\,r_0-i}\alpha_k^j\beta_k^{\,s_0-j},
\]
where \(0\le i\le r_0\) and \(0\le j\le s_0\). Choose a fixed
integer \(s_0\) sufficiently large in terms of \(\delta_0\) and the
compact ranges in \eqref{eq:old:5_20}. Put \(r_0=\lceil s_0 q/p\rceil\).
By \eqref{eq:old:5_20}, the possible \(r_0\) form a finite set and
\(r_0/s_0=q/p+O(1/s_0)+O(1/\log Y)\) uniformly. The leading terms
in \eqref{eq:old:5_13} have a uniform positive margin by \eqref{eq:old:5_22}. The ceiling and all
remaining terms in \eqref{eq:old:5_13} are uniformly \(O(1/s_0)+O(1/\log Y)\)
relative to those leading terms. First choose \(s_0\), then the common
threshold for \(Y\). Then these errors use less than half the margin.
Apply the level estimate \eqref{eq:old:5_12}--\eqref{eq:old:5_14} with
\(b_{\rm cell}=B_u+\xi+O(1/\log Y)\) and
\(\mu_{\rm th}=\mu_b+O(1/\log Y)\). It makes every
\(K_{\rm det}\)-row determinant zero for one common sufficiently large \(Y\).
Each entry is an integer of size at most
\[
 (Y+1)^{r_0+(21/20)s_0},
\]
since \(u\le Y^{a_0}<Y^{1/20}\). Hence cofactors give a nonzero
integer polynomial of degree at most \((r_0,s_0)\). Its coefficient
height is \(Y^{C_\xi}\) with fixed \(C_\xi\). This gives the required
integer auxiliary polynomial.

We next count the lattice radii over the moving intervals. Write
\[
\rho=\frac{\alpha/E}{\beta/F},\qquad
\rho=\kappa(1-\theta)\chi^{-2},\qquad
\kappa=\frac{F H_D^2}{E H_A^2}.
\]
Since each normalized coefficient lies in \([1,2)\), actual \(\rho\)
lies in \([1/2,2]\). For each interval, crop its image in \(\rho\) to
\((1/2,2)\). Empty intervals contain no source. The map
\(\chi\mapsto\kappa\chi^{-2}\) has derivative bounded above and below on
\([1/2,2]\), uniformly in every nonempty box. We have
\(1/8\le\kappa\le16\) by the dyadic box inequalities and \eqref{eq:old:5_15}.
Also \(\theta\le2Y^{-\mu_b}<1/N_u\). Hence the retained coefficient directions
have angular width \(O(1/N_u)\). Any angular window of width \(w\)
meets \(O(1+N_uw)\) of them.

For a retained direction with unit vector \(\ell\) and unit normal \(n\),
define
\[
\mathcal A(\alpha,\beta)=
 \left(\ell\cdot(\alpha/E,\beta/F),
       \frac{N_u}{c_0}n\cdot(\alpha/E,\beta/F)\right),
\]
where the fixed \(c_0\) makes the images under \(\mathcal A\)
of the source coefficient vectors lie in the disk of radius
\(5\). This linear map has determinant \(N_u/(c_0EF)\). Its image of the
integer coefficient lattice has a
Euclidean reduced basis \(b_1,b_2\), with angle at least \(60^\circ\).
The inverse images form a unimodular integer basis. For each source pair,
write \((v_1,v_2)\) for its coordinates in this reduced basis. These are
distinct from the original coefficient coordinates \((\alpha,\beta)\). Put
\(T_i=5\sqrt2/|b_i|\). Then \(T_2\le T_1\), every source coefficient
pair has lattice coordinates \(|v_i|\le T_i\), where
\(\mathbf v=(v_1,v_2)\), and
\[
 T_1T_2\gg \frac{EF}{N_u},\qquad T_1^2\gg\frac{EF}{N_u}.
\AutoTag{eq:old:5_23}
\]
The determinant of the lattice and the operator norms of \(\mathcal A\)
and \(\mathcal A^{-1}\) are bounded by fixed powers of \(Y\). Every
nonzero integer preimage has image length at least
\(\|\mathcal A^{-1}\|^{-1}\). Minkowski's theorem bounds the
shortest image vector by a constant times the square root of the
covolume. The reduced-basis angle and determinant then bound the second
basis vector by a fixed power of \(Y\) as well. Apply
\(\mathcal A^{-1}\). Then both unimodular preimage vectors have
polynomially bounded coordinates. Thus the coefficient-lattice change
keeps the auxiliary polynomial's height at \(Y^{O_\xi(1)}\).

We now bound the complete interval sum. Put
\(L_*=\max(1,(EF/N_u)^{1/2})\). Intervals with \(T_1\le2L_*\) contribute
\(O(N_u+(EFN_u)^{1/2})\). For larger \(T_1\), place it in a dyadic range
\((T,2T]\). Let \(\mathbf x_{\min}=(x_1,x_2)\) be the primitive
shortest integer preimage. The definition of \(\mathcal A\) gives
\[
 |(x_1/E,x_2/F)|\ll T^{-1},\qquad
 \operatorname{dist}((x_1/E,x_2/F),\ell)\ll (N_uT)^{-1}.
\]
A nonaxis dyadic annulus of normalized radius \(r\) contains
\(O(r^2EF)\) integer vectors, with
\(r\ge\max(E^{-1},F^{-1})\). A fixed vector is compatible with
\(O(1+(Tr)^{-1})\) interval directions. Sum
\(EF(r^2+r/T)\) over the annuli up to \(r\ll T^{-1}\). This gives
\(O(EF/T^2)\) intervals in this \(T\)-range.
Their contribution to \(\sum(T_1+1)\) is \(O(EF/T)\). The vertical and
horizontal primitive vectors are the two axes. The retained source
directions have coordinate ratio \(\rho\in[1/2,2]\). So they stay a fixed
positive angle from both axes. The corresponding vertical and horizontal
cases cost \(O(F+N_u)\) and \(O(E+N_u)\), respectively. Finally
\(T_1\gg(EF/N_u)^{1/2}\) by \eqref{eq:old:5_23}. Summing the dyadic ranges gives
\[
 \sum_{\text{intervals}}(T_1+1)
 \ll_\eta Y^\eta\{N_u+\sqrt{EFN_u}\log Y+E+F\}.
\AutoTag{eq:old:5_24}
\]

For each interval, factor its auxiliary polynomial over \(\mathbb Q\).
There are \(O_\xi(1)\) factors, each of fixed degree and height \(Y^{O_\xi(1)}\).
Make the unimodular coefficient-lattice change. Then let an irreducible
factor have degree \(a_c\) in the corner pair and \(b_q\) in the
coefficient pair. Write the original coefficients \(\alpha(\mathbf v),\beta(\mathbf v)\)
as homogeneous linear forms in the primitive
lattice coordinates \(\mathbf v=(v_1,v_2)\). The following cases account for
every factor.

\begin{enumerate}
\item If it is not absolutely irreducible, its distinct conjugate
components have no common curve component. By bihomogeneous Bezout they have \(O_\xi(1)\)
rational common points. Each rational point gives at
most four signed primitive pairs.
\item If \(a_c=0\), the factor fixes \(O_\xi(1)\) primitive coefficient
ratios. For each, \eqref{eq:old:5_15} is a square-fibre equation. It has positive
coefficients \(\alpha,\beta\) and nonzero shift \(u\). The coefficients and
shift are at most \((Y+1)^{21/20}\), and \(A\le Y+1\). Theorem~\ref{thm:fibre}
with \(x=2Y\) and fixed height exponent \(A_{\rm ht}=2\) gives \(Y^\eta\)
solutions uniformly.
\item If \(b_q=0\), the factor fixes \(O_\xi(1)\) primitive corner pairs.
The solutions in \((\alpha,\beta)\) to
\(\beta D^2-\alpha A^2=u\) form a progression with step
\((D^2,A^2)\). Its length in the \(\alpha\)-box is at most
\[
 1+\frac{E}{D^2}\ll 1+\frac{uY}{A^2D^2}
 \ll 1+\frac{Y^{-1/3+5e}}u=O(1)
\]
uniformly for \(u\ge1\), by \eqref{eq:old:5_17}--\eqref{eq:old:5_18}.
\item If \(a_c\ge2\) and \(b_q\ge1\), apply Heath--Brown's
Lemma~2. The first pair is the primitive coefficient-lattice pair,
bounded by \(T_1+1\). The second is the primitive corner pair. Its degree
in the second pair is \(a_c\ge2\). Choose the theorem's epsilon after
the fixed polynomial height exponent. Hence the count is
\(O_\eta(Y^\eta(T_1+1))\).
\item If \(a_c=1\) and \(b_q\ge1\), write the factor as
\(A Q(\mathbf v)+D P(\mathbf v)\), with coprime nonzero homogeneous forms
\(P,Q\). The source equation forces \(P(\mathbf v)=-A\zeta\),
\(Q(\mathbf v)=D\zeta\) for a nonzero integer \(\zeta\). The Sylvester
identities put
\(\operatorname{Res}(P,Q)v_i^{2b_q-1}\) in the ideal \((P,Q)\).
Primitivity of \(\mathbf v\) therefore gives
\(\zeta\mid\operatorname{Res}(P,Q)\).
The resultant has fixed polynomial height, so its divisors cost
\(Y^\eta\). Fix a divisor \(\zeta\mid\operatorname{Res}(P,Q)\). Then the
source conic gives the nonzero homogeneous
level equation
\[
\beta(\mathbf v) Q(\mathbf v)^2-\alpha(\mathbf v) P(\mathbf v)^2=u\zeta^2.
\]
The left side is a nonzero binary form of degree \(d_\Phi=2b_q+1\). If it
vanished identically, coprimality of \(P,Q\) would force \(P^2\) to
divide the linear form \(\beta(\mathbf v)\). This is impossible for \(b_q\ge1\).
Denote this form by \(\Phi\). Its level is \(h=u\zeta^2\ne0\). Define
\(\mathcal E=\{v_1\in\mathbb Z:\Phi(v_1,v_2)-h\text{ is identically zero in }v_2\}\).
Writing \(\Phi(v_1,0)=a_\Phi v_1^{d_\Phi}\), every \(v_1\in\mathcal E\) satisfies
\(a_\Phi v_1^{d_\Phi}=h\), so \(|\mathcal E|\le d_\Phi\). If \(a_\Phi=0\), then \(\mathcal E=\varnothing\).
This includes the pure-power case. There every slice is constant, and at most
\(d_\Phi\) slices equal the nonzero target \(h\). There are at most
\(2\lfloor T_1\rfloor+1\) other first coordinates. For each, the nonzero
polynomial in \(v_2\) has degree at most \(d_\Phi\). Hence it has at most \(d_\Phi\)
integer roots. Each exceptional coordinate contributes at most
\(2\lfloor T_2\rfloor+1\) points. Thus
\[
\begin{aligned}
\#\{(v_1,v_2)\in\mathbb Z^2:|v_i|\le T_i,\ \Phi(v_1,v_2)=h\}
&\le d_\Phi(2\lfloor T_1\rfloor+1)+d_\Phi(2\lfloor T_2\rfloor+1)\\
&\le4d_\Phi(T_1+1),\qquad T_2\le T_1.
\end{aligned}
\]
The corners are then forced.
\end{enumerate}

Thus every full interval contributes \(Y^\eta(T_1+1)\). The underfull
intervals cost \(O_\xi(N_u)\). Fix the ordered residual factorizations
\(\alpha=ac,\ \beta=bd\). The inverse formulas \eqref{eq:old:2_2} determine \(B,C\),
then \(m=ABa,n=ACc\). The positivity, integrality, gcd, mask, shell and
first-match filters only discard candidates. The inverse multiplicity is
at most \(\tauDiv(\alpha)\tauDiv(\beta)\ll_\eta Y^\eta\), since
\(\alpha,\beta\le(Y+1)^{21/20}\). The opposite shift is recovered from
\(ut=L\).
Combine this with \eqref{eq:old:5_24}. For one fixed \(u\) and one box the count is
\[
 \ll_\eta Y^\eta\{N_u+\sqrt{EFN_u}\log Y+E+F\}.
\AutoTag{eq:old:5_25}
\]

Split the \(3u\) population counted by \eqref{eq:old:5_25} into
\(\min(A,D)\ge Y^{1/3-\nu}\) and its complement. In the first part,
\eqref{eq:old:5_17} gives
\[
 E+F\ll uY^{1/3+2\nu}.
\AutoTag{eq:old:5_26}
\]
In the second part, suppose \(H_A\le H_D\). The other case is symmetric.
For fixed \(A,\beta\), the congruence
\(\beta D^2\equiv u\pmod{A^2}\) has
\(O(2^{\omega(A)})=Y^\eta\) roots modulo \(A^2\), since
\(\gcd(\beta u,A)=1\). Each root class contributes
\(O(H_D/A^2+1)\) values of \(D\). Also
\(\beta\ll uY/H_D^2\). Summing over \(A,\beta\) gives
\[
    \ll_\eta Y^\eta\left\{\frac{uY}{H_AH_D}
                 +\frac{uYH_A}{H_D^2}\right\}.
\AutoTag{eq:old:5_27}
\]
For odd prime powers the unit-square congruence has at most two roots.
At powers of \(2\) it has at most four. The absolute factor \(4\) is
absorbed in \(O(2^{\omega(A)})\).
By \eqref{eq:old:5_18} and \(H_A\le Y^{1/3-\nu}\), the two terms in \eqref{eq:old:5_27} are at
most
\[
 Y^{1/3+5e/2+o(1)},\qquad
 \frac{Y^{2/3+5e-3\nu+o(1)}}u.
\AutoTag{eq:old:5_28}
\]
The second term sums harmonically over \(u\). The first costs
\(\sum_{u\le Y^{a_0}}1\ll Y^{a_0}\).

Finally \(EF\le M<\kappa_Y^2Y^{q_+}\). Sum \eqref{eq:old:5_25}--\eqref{eq:old:5_28} over
\(u\le Y^{a_0}\), all boxes and the logarithmic factors. Before the final
loss allocation, this gives five exponents and their savings from $2/3$
\[
\begin{array}{@{}ll@{}}
\text{term and exponent} & \text{saving from }2/3\\[8pt]
\text{square-root term: }a_0+(q_++B_u+\xi)/2 &
  \dfrac{5102312049347}{600000000000000}\ (0.00850385\ldots)\\[7pt]
\text{interval term: }a_0+B_u+\xi &
  \dfrac{23206592049347}{300000000000000}\\[7pt]
\text{core axes: }1/3+2\nu+2a_0 & \dfrac{2473}{12000}\\[7pt]
\text{small corner, first: }a_0+1/3+5e/2 & \dfrac{1619393}{6000000}\\[7pt]
\text{small corner, second: }2/3+5e-3\nu & \dfrac{8003}{500000}
\end{array}\AutoTag{eq:old:5_29}
\]
The smallest saving is \(5102312049347/600000000000000>S_0\).
Fix the degree and height constants. Then all fixed-height, divisor,
dyadic and logarithmic losses can be assigned arbitrarily small
exponents. Therefore, for every \(\varepsilon>0\),
\[
 \#3u\ll_\varepsilon
 Y^{2/3-5102312049347/600000000000000+\varepsilon}.
\AutoTag{eq:old:5_30}
\]
The next proposition records this display as the complete first-match row
estimate.
\begin{rowproposition}[\(3u\) row estimate]\label{prop:row-3u}
For every \(\varepsilon>0\), the complete first-match row satisfies
\[
 \#3u\ll_\varepsilon
 Y^{2/3-5102312049347/600000000000000+\varepsilon}.
\]
\end{rowproposition}
\begin{proof}
The full row count is \eqref{eq:old:5_30}. Its five additive terms are bounded in
\eqref{eq:old:5_25}--\eqref{eq:old:5_29}. The least saving is the one displayed there.
\end{proof}

\subsection*{The short-\(t\) row \(3t\)}

In the \(t\)-chart use \(\gamma=ad,\delta=bc\). The exact source equations
are
\[
 \delta C^2-\gamma B^2=t,\qquad
 \frac{\delta}{\gamma}=\left(\frac BC\right)^2(1+\beta_{\rm rem}),\qquad
 \beta_{\rm rem}=\frac{n-m}{m(n+1)}\in(0,1/m).
\AutoTag{eq:old:5_31}
\]
By \eqref{eq:old:5_11}, \(\beta_{\rm rem}\ll Y^{-\tau_0}\). The source heights and the
\(O_R\), \(O_M\), \(Q1\) and \(Q2t\) complements give
\[
 \frac t4Y^{q_-}<BC\le3tY^{c_t},\quad
 \gamma,\delta\ll tY^{g_t},\quad
 \frac{\gamma}{C^2}\le \frac{6Y^{-d_0}}t<1
 \quad\text{for all sufficiently large }Y,
\AutoTag{eq:old:5_32}
\]
where
\[
 c_t=\frac23+l_0=q_+,\quad
 g_t=\frac13+e+l_0+2j_0=\frac{5833463}{15000000},\quad
 d_0=\frac13-2e>0.
\]
For the \(\gamma/C^2\) bound, start from the exact source identity
\[
 n-m=\frac{ZL}{m+n+1}.
\]
Then use
\(BC=(m+n+1)t/L\), \(m+n+1>Y/2\), and
\(L<\kappa_Y Y^{1/3+e}\). Since \(2\kappa_Y\le3\), this gives
\(BC>(t/3)Y^{2/3-e}>(t/4)Y^{q_-}\). For the upper bound use
\(L\ge Y^{1/3-l_0}\) and \(m+n+1\le2Y+1\le3Y\). This gives
\(BC\le3tY^{2/3+l_0}\). Also
\[
 \gamma B^2=\frac{t\,m(n+1)}{n-m},\qquad
 \gamma/C^2=\frac{m(n+1)}{ZBC};
\]
using \(Z>Y^{5/3-e}\) and the lower bound for \(BC\) proves the last
inequality in \eqref{eq:old:5_32}. The bound for \(\delta\) is symmetric. It uses
\(\delta C^2=t\,n(m+1)/(n-m)\) and \(B,C\ge Y^{1/3-j_0}\).

Fix \(t\le Y^{a_0}\) and a dyadic box
\[
 B\asymp H_B,\quad C\asymp H_C,\quad
 \gamma\asymp E,\quad\delta\asymp F.
\]
Put \(\mathcal L=\ln Y\), \(p=\log_Y(8H_BH_C)\),
\(q=\log_Y(128EF)\), \(z=\log_Yt\), and
\[
 \Xi_t=\frac{H_BH_C\sqrt{EF}}{8t},\quad
 \mu_0=\log_Y \Xi_t,\quad
 \ell_1=\ln(8H_BH_C),\quad \ell_2=\ln(128EF).
\]
From \eqref{eq:old:5_11}--\eqref{eq:old:5_32}, with \(\kappa_Y\le3/2\),
\[
 q_-\mathcal L-\ln2<\ell_1\le(c_t+a_0)\mathcal L+\ln24,\quad
 q_-\mathcal L+\ln32<\ell_2<q_+\mathcal L+\ln288,
\]
\[
 \tau_0\mathcal L-\ln256<\ln \Xi_t<(1+e+l_0)\mathcal L,\qquad
 p-z\ge q_--O(1/\mathcal L).
\AutoTag{eq:old:5_33}
\]
Dyadicity gives \(BC/4<H_BH_C\le BC\) and
\(EF\le M<4EF\). With \eqref{eq:old:5_32} these imply
\(8H_BH_C>tY^{q_-}/2\),
\(8H_BH_C\le24tY^{c_t}\),
\(128EF>32Y^{q_-}\), and \(128EF<288Y^{q_+}\).
They also give
\[
 \frac{Y^{\tau_0}}{256}<\Xi_t
 \le\frac{3\kappa_Y}{8}Y^{3q_+/2}
 \le\frac9{16}Y^{1+e+l_0},
\]
which proves the two \(\Xi_t\) bounds in \eqref{eq:old:5_33}.
In particular \(\ell_1,\ell_2,\ln \Xi_t\) lie in
\([\mathcal L/5,2\mathcal L]\) for \(Y\ge2048\).

Set \(\sigma_t=H_B/H_C\) and \(x=(B/C)/\sigma_t\). The source has
\(x\in(1/2,2)\). Partition \(x\in(1/2,2)\) into
\[
 I_k=\left(\frac{k}{N},\frac{k+1}{N}\right].
\]
Equivalently, these are the intervals
\((k\sigma_t/N,(k+1)\sigma_t/N]\) for \(B/C\). Here
\(N/2-1\le k\le2N-1\), with \(N\ge40\) to be fixed below. Let
\(x_4=\lfloor k^2/N\rfloor\) and
\(r_k=x_4\sigma_t^2/N\). For \(x\in I_k\), the width of the interval
and \eqref{eq:old:5_31} give
\[
 |\delta-r_k\gamma|\le18\sigma_t^2E/N,\qquad 0<\gamma<2E.
\AutoTag{eq:old:5_34}
\]
The square of \(B/C\) differs from \(r_k\) by at most
\(5\sigma_t^2/N\). The additional term
\((B/C)^2\beta_{\rm rem}\) is at most \(4\sigma_t^2/m\), and the choice below has
\(N<m\). The endpoint \(k\le2N-1\) follows from \(x<2\). This is the
range used in the \(5\sigma_t^2/N\) bound. A nonempty box also gives
\[
 F/16<\sigma_t^2E<8F.
\AutoTag{eq:old:5_35}
\]
Define the linear map
\[
 \mathcal A_k(y_1,y_2)=
 \left(\frac{N}{\sigma_t^2E}(y_1-r_ky_2),\frac{y_2}{E}\right).
\]
Its image of \(\mathbb Z^2\) has covolume
\(N/(\sigma_t^2E^2)\asymp N/(EF)\). By \eqref{eq:old:5_34}, coefficient pairs
\((\delta,\gamma)\) in the box map into the disk of radius \(19\).
Choose a Euclidean reduced basis \(b_1,b_2\) for this lattice. Let
\(h_1,h_2\) be its unimodular integer preimages. Put
\(T_i=19\sqrt2/|b_i|\). The reduced-basis angle is at least \(60^\circ\).
Hence \(T_2\le T_1\), every source pair has coordinates
\(|v_i|\le T_i\), where \(\mathbf v=(v_1,v_2)\), and \eqref{eq:old:5_35} gives
\[
 T_1T_2
 =\frac{722}{|b_1||b_2|}
 \ge\frac{361\sqrt3}{\operatorname{covol}(\mathcal A_k(\mathbb Z^2))}
 =\frac{361\sqrt3\,\sigma_t^2E^2}{N}
 >\frac{361\sqrt3}{16}\frac{EF}{N}
 >39\frac{EF}{N},\qquad T_1^2\ge T_1T_2.
\AutoTag{eq:old:5_36}
\]
The matrix \(\mathcal A_k\), its inverse, and the covolume have polynomial
height. Every nonzero integer preimage has image length at least
\(\|\mathcal A_k^{-1}\|^{-1}\). Minkowski's theorem bounds the shortest
image vector by a constant times the square root of the covolume. The
reduced-basis angle and determinant then give a polynomial upper bound
for the second vector. Apply \(\mathcal A_k^{-1}\). Then both
unimodular integer preimages have polynomially bounded coordinates.

\begin{samepage}
We now count the moving grid. The integer \(x_4\) lies in
\([N/5,4N]\). Each \(x_4\) comes from at most two \(k\), since
\[
 \sqrt{Nx_4}\le k<\sqrt{N(x_4+1)}
\]
has length at most \(\sqrt5/2<2\) for \(x_4\ge N/5\). Therefore it
contains at most two integers. Take a dyadic range \(T_1\in(T,2T]\). Write
\(\mathcal A_k(y)=b_1\) for the primitive shortest vector
\(y=(y_1,y_2)\). Then
\[
 |y_1-r_ky_2|\ll \sigma_t^2E/(NT),\qquad
 |y_2|\ll E/T.
\AutoTag{eq:old:5_37}
\]
If \(y_2=0\), primitivity gives \(y_1=\pm1\). The contribution
over all \(k\) is then \(O(F+N)\). If \(y_1=0\), then \(y_2=\pm1\).
Since \(x_4\ge N/5\), \eqref{eq:old:5_37} bounds its total contribution by
\(O(E+N)\). For \(y_1y_2\ne0\), \eqref{eq:old:5_37} gives
\(T\ll E\) and \(|y_1|\ll F/T\). Multiply the first inequality in
\eqref{eq:old:5_37} by \(N/\sigma_t^2\). Then, for each fixed \(y_1\), the nonzero
integer \(y_2x_4\) lies in an interval of length \(O(E/T)\). Since
\(E/T\gg1\), there
are \(O(E/T)\) such integer values. Their absolute values satisfy
\(|y_2x_4|\ll NE/T\), a fixed power of \(Y\). So each has \(Y^\eta\) factorizations
as \(y_2x_4\). Each \(x_4\) has at most two \(k\). Thus the number of
intervals in this \(T\)-range is
\(O_\eta(Y^\eta EF/T^2)\). This includes the endpoint cases
\(|y_1|\le C F/T\) and \(E/T\ge c>0\). Their contribution to
\(\sum_k(T_1(k)+1)\) is \(O_\eta(Y^\eta EF/T)\). By \eqref{eq:old:5_36},
\(T\gg(EF/N)^{1/2}\). Sum the dyadic \(T\)-ranges and pay the
\(+1\) terms. This gives
\[
 \sum_k(T_1(k)+1)\ll_\eta
 Y^\eta\{N+\sqrt{EFN}\log Y+E+F\}.
\AutoTag{eq:old:5_38}
\]
\end{samepage}

For each interval \(I_k\), apply the determinant construction \eqref{eq:old:5_12}
to the normalized monomials \(x^i\zeta^j\), where
\(\zeta=(E/F)(\delta/\gamma)=\kappa_t x^2+y\). Scaling rows and columns
back to the integer monomials gives \eqref{eq:old:5_42}. We choose the grid
size \(N\) used above as follows. Put
\[
B_t=\frac34(q_++a_0)
 =\frac{5452607}{10000000},\quad
 \nu_t=B_t+\xi=\frac{5453607}{10000000},\quad
 N_t=\lceil Y^{\nu_t}\rceil .
\AutoTag{eq:old:5_39}
\]
The exact gap
\[
 \tau_0-\nu_t=\frac{2226893}{5000000}>0
\]
shows that \(N_t<m\) for all sufficiently large \(Y\). Here
\(m>Y^{\tau_0}/\sqrt3\) and \(N_t\le2Y^{\nu_t}\). In
\eqref{eq:old:5_34}--\eqref{eq:old:5_38} take \(N=N_t\).
Since \(EF>M/4>Y^{q_-}/4\), we have
\[
\frac{EF}{N_t}\gg Y^{q_- -\nu_t}
=Y^{690793/6000000}.
\]
Thus \eqref{eq:old:5_36} puts every dyadic radius \(T\) above
\(cY^{690793/12000000}>1\) for all sufficiently large \(Y\).
The separately counted \(+1\) per grid interval is already included in
the \(O(N_t)\) term of \eqref{eq:old:5_38}.

The determinant thickness \(\vartheta\) in the normalized coefficient
ratio is comparable to \(\Xi_t^{-1}\). The conic gives
\[
 \frac{\delta}{\gamma}
 =\left(\frac BC\right)^2+\frac{t}{\gamma C^2}.
\]
Normalize the coefficient ratio by \(E/F\). Then the thickness is
\((E/F)t/(\gamma C^2)\asymp t/(F H_C^2)\). From \eqref{eq:old:5_35},
\(F H_C^2\asymp E H_B^2\asymp H_BH_C\sqrt{EF}\). Hence this thickness is
\(\asymp \Xi_t^{-1}\). More precisely, the dyadic bounds
\(E\le\gamma<2E\), \(H_C\le C<2H_C\), and
\(1/16<\sigma_t^2E/F<8\) imply
\[
 \frac{\Xi_t^{-1}}{256}\le\vartheta\le\Xi_t^{-1}.
\]
Thus \(0\le\mu_{\rm th}-\mu_0\le\ln256/\mathcal L\). The normalized
cell width is exactly \(1/N_t\), so
\(\nu_t\le b_{\rm cell}=\log_Y N_t\le\nu_t+\ln2/\mathcal L\).

The box bounds also give
\[
 \frac{q_-}{c_t+a_0}+O(1/\mathcal L)
 \le\frac{\ell_2}{\ell_1}
 \le\frac{q_+}{q_-}+O(1/\mathcal L).
\]
The limiting endpoints are \(9907393/10905214>0.9\) and
\(10185214/9907393<1.1\). So this ratio belongs to \([0.9,1.1]\) for
all sufficiently large \(Y\). Choose a fixed integer \(s_0\)
sufficiently large in terms of the strict margin below and this compact
range. Set
\[
 r_0=\left\lfloor s_0\frac{\ell_2}{\ell_1}\right\rfloor .
\]
The possible \(r_0\) form a finite set. The degree rounding and the
lower-order terms in \eqref{eq:old:5_13} are uniformly \(O(1/s_0)\) relative to the
leading terms. Choose \(s_0\) first. Then enlarge the threshold for \(Y\)
to absorb the \(O(1/\mathcal L)\) errors.

The normalized scale parameters in \eqref{eq:old:5_12} are
\(p=\ell_1/\mathcal L\), \(q=\ell_2/\mathcal L\), and
\(\mu_{\rm th}=p+q/2-z+O(1/\mathcal L)\). Precisely,
\[
 \mu_0=p+q/2-z-\frac{\log(64\sqrt{128})}{\mathcal L}.
\]
The displayed constant and the additional \(O(1/\mathcal L)\) from the
thickness comparison are among the finite errors accounted for below.
The limiting function
\[
 (p,q,z)\longmapsto \frac98\frac{pq}{p+q/2-z}
\]
is increasing in \(p,q,z\) on the source range. The relevant derivative
numerators are \(q(q/2-z)\), \(p(p-z)\) and \(pq\). With
\(D=p+q/2-z\), its three partial derivatives are
\[
\frac98\frac{q(q/2-z)}{D^2},\qquad
\frac98\frac{p(p-z)}{D^2},\qquad
\frac98\frac{pq}{D^2}.
\]
On the stated source domain each has absolute value at most
\((9/8)\cdot4\cdot25=112.5<120\). Here
\[
 q_-/2-a_0=\frac{8467393}{30000000}>0,\qquad
 q_--a_0=\frac{9187393}{15000000}>0.
\]
The bound \(BC\le3tY^{c_t}\) gives the one-sided endpoint estimates
\[
p\le c_t+z+\frac{\ln24}{\mathcal L},\qquad
q\le q_++\frac{\ln288}{\mathcal L}.
\]
First put \(p=c_t+z+O(1/\mathcal L)\). Then the denominator is
\(c_t+q/2+O(1/\mathcal L)\), and \(q/(c_t+q/2)\) increases. Thus the
maximum is \(B_t+O(1/\mathcal L)\), attained at \(q=q_+\), \(z=a_0\).
For \(\delta_0=\xi/4=1/40000\),
\[
 \nu_t-(1+\delta_0)B_t
 =\xi-\delta_0B_t
 =\frac{34547393}{400000000000}>0.
\AutoTag{eq:old:5_40}
\]
Now the uniform error. The upper endpoints in \eqref{eq:old:5_33} change \(p,q\)
from \(c_t+z,q_+\) by at most \(\ln24/\mathcal L,\ln288/\mathcal L\).
The exact normalization of \(\mu_0\) contributes
\(\ln(64\sqrt{128})/\mathcal L\). The thickness comparison contributes
at most \(\ln256/\mathcal L\). Take the source range
\(0\le p,q,z\le2\) and \(p+q/2-z\ge1/5\). There every first derivative of
\(\Phi(p,q,z)=(9/8)pq/(p+q/2-z)\) has absolute value at most \(120\).
Since \(\ln24<4\), \(\ln288<6\), \(\ln(64\sqrt{128})<7\), and
\(\ln256<6\), the total perturbation is at most \(2760/\mathcal L\).
We may therefore take \(C_t=4000\). Hence, whenever
\[
(\xi-\delta_0B_t)\ln Y>2C_t(1+\delta_0),
\AutoTag{eq:old:5_41}
\]
the determinant condition has the strict surplus required in \eqref{eq:old:5_14}.
Take one interval with \(K_{\rm det}=(r_0+1)(s_0+1)\) source rows. Use columns
indexed by \(0\le i\le r_0,\ 0\le j\le s_0\), and multiply each row by
\(C^{r_0}\gamma^{s_0}\). The resulting entry is
\[
 C^{r_0}\gamma^{s_0}(B/C)^i(\delta/\gamma)^j
 =B^iC^{r_0-i}\delta^j\gamma^{s_0-j}\in\mathbb Z.
\AutoTag{eq:old:5_42}
\]
We have \(B,C\le Y+1\) and \(\gamma,\delta\ll Y^{a_0+g_t}\), and the degrees are
fixed. So every matrix entry and cofactor has height \(Y^{O(1)}\). By
\eqref{eq:old:5_12}--\eqref{eq:old:5_14}, every integerized \(K_{\rm det}\times K_{\rm det}\) determinant has absolute
value less than one. Hence it vanishes. If an interval has fewer than
\(K_{\rm det}\) source rows, count those points directly. Otherwise the integer
matrix has a nonzero rational kernel vector. Clearing its denominator
gives a nonzero integral auxiliary polynomial on that interval. It has
bounded bidegree and height \(Y^{O_\xi(1)}\).

In this \(t\)-chart, \(\gcd(\gamma,\delta)=\gcd(B,C)=1\). The
transformed auxiliary factors have height \(Y^{O_\xi(1)}\) by the
fixed-degree factor-height argument above. These are the primitivity and
height hypotheses used in the square-fibre and Heath--Brown counts.
The coefficient-lattice coordinates \(\mathbf v=(v_1,v_2)\) are primitive,
since \((\gamma,\delta)\) is primitive and the basis change is
unimodular. The original coefficients \(\gamma(\mathbf v),\delta(\mathbf v)\)
are homogeneous linear forms in these coordinates.
Factor an auxiliary form over \(\mathbb Q\). Let \(a_c\)
and \(b_q\) be its degrees in the corner and coefficient pairs. There are
five branches. The conjugate-component Bezout argument above handles a
factor that is not absolutely irreducible. If \(a_c=0\),
it fixes \(O_\xi(1)\) primitive coefficient pairs. Then the
square-fibre theorem (Theorem~\ref{thm:fibre} in
Appendix~\ref{app:fibre}) applies to
\(\delta C^2-\gamma B^2=t\), with \(x=2Y\) and fixed height exponent
\(A_{\rm ht}=2\). If \(b_q=0\), it fixes a primitive corner pair
\((B,C)\). Reducing the conic modulo \(C^2\) fixes \(\gamma\). Also
\(0<\gamma<C^2\) by \eqref{eq:old:5_32}. So there is at most one coefficient pair.
If \(a_c\ge2\) and \(b_q\ge1\), apply Heath--Brown's Lemma~2 with the
primitive coefficient pair first. It gives \(O_\eta(Y^\eta(T_1+1))\).
Finally, if \(a_c=1\) and \(b_q\ge1\), write the factor as
\(B Q(\mathbf v)+C P(\mathbf v)\), with coprime nonzero homogeneous forms
\(P,Q\). The conic forces \(P(\mathbf v)=-B\zeta\),
\(Q(\mathbf v)=C\zeta\). The Sylvester
identities and primitivity give \(\zeta\mid\operatorname{Res}(P,Q)\).
There are \(Y^\eta\) choices for \(\zeta\). For each, the conic gives
\[
 \delta(\mathbf v) Q(\mathbf v)^2-\gamma(\mathbf v) P(\mathbf v)^2=t\zeta^2.
\]
  The left side is a nonzero binary form of degree \(2b_q+1\). An identity
  would force \(P^2\) to divide the linear form \(\delta(\mathbf v)\), contrary
  to \(b_q\ge1\). The preceding complete slice count applies with the
  nonzero target \(h=t\zeta^2\). It counts \(O_\eta(T_1+1)\) coefficient
  pairs. Then \(B,C\) are forced. Thus every full interval
  contributes \(Y^\eta(T_1+1)\).

Fix the ordered factorizations \(\gamma=ad,\delta=bc\). The inverse
formulas \eqref{eq:old:2_2} determine
\[
 A=(dB-bC)/t,\qquad D=(cC-aB)/t,
\]
then \(m=ABa,n=ACc\). Also \(u=(cd-ab)/t\). These are forced, not
summed. The factor multiplicity is
\(\tauDiv(\gamma)\tauDiv(\delta)\ll_\eta Y^\eta\). Thus, for fixed \(t\) and
box, \eqref{eq:old:5_38} gives
\[
 \ll_\eta Y^\eta\{N_t+\sqrt{EFN_t}\log Y+E+F\}.
\AutoTag{eq:old:5_43}
\]
The row \(3t\) is contained in this larger population with
\(1\le t\le Y^{a_0}\), \(L\ge Y^{1/3-l_0}\) and \(B,C\ge Y^{1/3-j_0}\).
The restrictions \(t<u\) and \(K_{\rm res}\le Y^{\Delta_{\rm exp}}\) only reduce it.
There are \(O((\log Y)^4)\) boxes. By \eqref{eq:old:5_32}, \(EF\le M\ll Y^{q_+}\)
and \(E+F\ll tY^{g_t}\). Sum \eqref{eq:old:5_43} once over the positive shifts
\(t\le Y^{a_0}\). This gives
\[
 \#3t\ll_\eta Y^\eta\left\{
 Y^{a_0+\nu_t}
 +Y^{a_0+(q_++\nu_t)/2}
 +Y^{g_t+2a_0}\right\}.
\AutoTag{eq:old:5_44}
\]
The three exponents and their savings are
\[
\begin{array}{c|c|c}
\text{term}&\text{exponent}&\text{saving from }2/3\\ \hline
\text{interval count}&5933607/10000000&2199179/30000000\\
\text{square-root count}&39611249/60000000&388751/60000000\ (0.00647918\ldots)\\
\text{coefficient axes}&7273463/15000000&2726537/15000000 .
\end{array}\AutoTag{eq:old:5_45}
\]
The least saving is \(388751/60000000>S_0\). Fix the uniform degree and height constants. Then the fixed loss terms,
divisor factors and logarithms fit within any prescribed
\(\varepsilon>0\). Hence
\[
 \#3t\ll_\varepsilon
 Y^{2/3-388751/60000000+\varepsilon}.
\AutoTag{eq:old:5_46}
\]
\begin{rowproposition}[\(3t\) row estimate]\label{prop:row-3t}
For every \(\varepsilon>0\), the complete first-match row satisfies
\[
 \#3t\ll_\varepsilon
 Y^{2/3-388751/60000000+\varepsilon}.
\]
\end{rowproposition}
\begin{proof}
The complete row is bounded in \eqref{eq:old:5_44}. Its three contributions and
their savings are listed in \eqref{eq:old:5_45}--\eqref{eq:old:5_46}.
\end{proof}

\section{The intermediate rows}\label{sec:intermediate}

Throughout this section abbreviate
\[
 e=\epsilon_0,\qquad k=k_0,\qquad l=l_0,\qquad
 j=j_0,\qquad p=p_0,\qquad \mu_{\rm HP}=\frac{16}{25},\qquad
 S_0=\frac{30869}{5000000}.
\]
The first five and $S_0$ are the fixed parameters from
\eqref{eq:old:2_5}; $\mu_{\rm HP}=16/25$ is set here for the 4P grid.

In this section the chart and all source predicates are those of
Section~\ref{sec:partition}.  If the selected chart is the $u$-chart, then
$(J,K;f,g;s)=(A,D;ac,bd;u)$.  If it is the $t$-chart, then
$(J,K;f,g;s)=(B,C;ad,bc;t)$.  In either case
\begin{equation}
 gK^2-fJ^2=s>0,\qquad fg=M,\qquad ut=L.
\AutoTag{eq:old:6_1}
\end{equation}
The chart choice is unique, including equality $u=t$.  Fix the residual
factors, the chart, the shift and the two selected corners.  Then the
inverse formulas in \eqref{eq:old:2_2} determine the original pair.  All
counts below keep the positivity, integrality, gcd, mask, shell and
first-match tests when a relaxed tuple is reconstructed.

\subsection{The supported-part row $4E$}

Let $f_J$ be the full prime-power part of $f$ supported on primes dividing
$J$.  Let $g_K$ be the full prime-power part of $g$ supported on primes
dividing $K$.  Write
\[
 \begin{aligned}
 f&=d_fv,\qquad g=d_gw,\qquad d_f=f_J,\qquad d_g=g_K,\\[-2pt]
 \operatorname{rad}(d_f)&\mid J,\qquad
 \operatorname{rad}(d_g)\mid K.
 \end{aligned}                                          \AutoTag{eq:old:6_2}
\]
Thus $d_f,d_g$ contain all powers of their supported primes.  For fixed
\((f,J)\) and \((g,K)\), each is uniquely determined.

Explicitly, \(d_f\) is the unique largest divisor of \(f\) all of whose
prime divisors divide \(J\), and \(v=f/d_f\).  Similarly, \(d_g\) is the
unique largest divisor of \(g\) supported on \(K\), and \(w=g/d_g\).  Once
\(g\) and \(K\) are recovered in the \(v\)-tag, \(d_g\) is determined by
\((g,K)\).  The same holds for \(d_f\) in the \(w\)-tag.  The mask assigns
every supported prime power to these determined parts.

In the \(u\)-chart, \(JK=AD=(n-m)/t\), so \(JK\le Yu/L\).  In the
\(t\)-chart, \(JK=BC=(m+n+1)/u\), so \(JK\le2Yt/L\).  We have
\(s=u\) in the first chart, \(s=t\) in the second, and \(ut=L\).  Hence
both inequalities give \(JK\le2Ys/L\).  The failed \(Q1\) and
selected-corner \(Q2\) tests give the weak lower bounds below.  Thus
\[
 JK\le \frac{2Ys}{L},\qquad
 L\ge Y^{1/3-l},\qquad J,K\ge Y^{1/3-j}.
\AutoTag{eq:old:6_3}
\]
Hence
\[
 J,K\le 2sY^{1/3+l+j}.                                  \AutoTag{eq:old:6_4}
\]
The row has $Y^{a_0}<s<Y^{b_0}$ and
\[
 \min(v,w)\le Y^p.
\AutoTag{eq:old:6_5}
\]
It has at most four disjoint tags.  They are the selected chart and either
$v\le Y^p$, or $v>Y^p$ and $w\le Y^p$.  By \eqref{eq:old:6_5}, the second
alternative covers the complement of the first.

Each tag uses a uniform supported-power bound.
For fixed $C,D,\alpha>0$, if $J\le Y^C$, then
\[
 \#\{d\le Y^D:\operatorname{rad}(d)\mid J\}
 \ll_{C,D,\alpha}Y^\alpha.                           \AutoTag{eq:old:6_6}
\]
To prove it, put $\sigma=\alpha/(2D)$.  Rankin's inequality gives
\[
 \#\{d\le Y^D:\operatorname{rad}(d)\mid J\}
 \le Y^{D\sigma}\prod_{q\mid J}(1-q^{-\sigma})^{-1}.
\]
Fix any $\delta>0$.  All primes above a fixed threshold satisfy
$-\log(1-q^{-\sigma})\le\delta\log q$.  The product over the finitely many
smaller primes is a constant.  The remaining product is at most
$J^\delta$.  Taking $\delta=\alpha/(2C)$ proves \eqref{eq:old:6_6},
uniformly in the moving label $J$.

In the tag $v\le Y^p$, enumerate $(s,v,J,d_f)$ and put
\[
 N=vd_fJ^2+s=fJ^2+s.
\]
For an actual source, $N=gK^2>0$.  Thus $K^2\mid N$, and each divisor
choice for $K^2$ forces $g=N/K^2$.  There are at most $\tauDiv(N)$ choices.
In the other tag enumerate $(s,w,K,d_g)$ and put
\[
 N=wd_gK^2-s=gK^2-s.
\]
Keep only \(N>0\).  On an actual source \(N=fJ^2\), so \(J^2\mid N\) and
$f=N/J^2$ is forced.  Discarding \(N\le0\) is safe, since an actual source
has \(N=fJ^2>0\).

Once $f,g$ are known, the ordered residual splittings are
\[
 (a,c)\text{ with }ac=f,\quad (b,d)\text{ with }bd=g
 \quad\text{or}\quad
 (a,d)\text{ with }ad=f,\quad (b,c)\text{ with }bc=g,
\]
according to the chart.  Their total number is at most $\tauDiv(f)\tauDiv(g)$.
Each splitting and $(s,J,K)$ determine the remaining corners, and then
$(m,n)$ by \eqref{eq:old:2_2}.  Candidates failing an original-source
predicate are discarded.  Once the ordered residuals are fixed, the
reconstruction has multiplicity at most one.  Then $ut=L$ determines the
opposite shift.

Before the divisor estimates, discard formal candidates that violate the
original source-height conditions.  Every retained source has
$J,K\le Y+1$, $f,g,d_f,d_g\le(Y+1)^2$, and
$1\le N\le(Y+1)^4$.  Apply \eqref{eq:old:6_6} to the one supported part.
Apply the standard estimate $\tauDiv(q_0)\ll_\rho q_0^\rho$ to each of $N,f,g$.
Fix a target loss $\eta>0$ and put $T=Y+1$. Apply \eqref{eq:old:6_6}
with $T$ in place of its displayed base variable $Y$, and take
$C=1$, $D=2$, and output exponent $\alpha=\eta/4$.  Its Rankin
parameters may be $\sigma=\eta/16$ and $\delta=\eta/8$.  For the divisor
bounds, take $\rho_N=\eta/16$ for $\tauDiv(N)$ and
$\rho_f=\rho_g=\eta/8$ for $\tauDiv(f),\tauDiv(g)$.  Since
$N\le T^4$ and $f,g\le T^2$, each of the four multiplicative factors is
$O_\eta(T^{\eta/4})$.  Their product is
$O_\eta(T^\eta)=O_\eta(Y^\eta)$, uniformly in the retained source
labels.  Thus each fixed $(s,v,J)$, or its counterpart $(s,w,K)$,
contributes at most $Y^\eta$ complete source records.  Sum over the small
endpoint, then over the selected corner using \eqref{eq:old:6_4}, then over
$s<Y^{b_0}$.  This gives
\[
 \#4E\ll_\eta
 Y^{p+1/3+l+j+\eta}\sum_{1\le s<Y^{b_0}}s
 \ll_\eta Y^{1/3+2b_0+l+j+p+\eta}.
\AutoTag{eq:old:6_7}
\]
The direct sum satisfies $\sum_{s<Y^{b_0}}s\ll Y^{2b_0}$.  At the parameter
values in \eqref{eq:old:2_5}
\[
 \frac13+2b_0+l+j+p=\frac{9907393}{15000000}=\frac23-S_0.
\AutoTag{eq:old:6_8}
\]
This equality, with $\eta$ arbitrary, is the claimed $4E$ bound.
\begin{rowproposition}[4E row estimate]\label{prop:row-4e}
For every \(\eta>0\), the complete first-match row satisfies
\[
 \#4E\ll_\eta Y^{2/3-S_0+\eta}.
\]
\end{rowproposition}
\begin{proof}
The full count is \eqref{eq:old:6_7}, and \eqref{eq:old:6_8} gives its exponent.
\end{proof}

\subsection{The opposite-shift cubic row $4H$}

For each prime dividing a residual, the radical mask places that prime in
one of the four corners.  The adjacent coprimalities determine which
corner.  Primes in \(a,b,c,d\) go respectively to one of
\((A,B),(C,D),(A,C),(B,D)\).  These choices are unique, since the two
corners in each listed pair are coprime.  Apply the same assignment to the
full prime powers.  This gives four pairwise coprime parts, assigned to
corners, whose product is \(M\):
\[
 (ac)_A,\quad (bd)_D,\quad (ad)_B,\quad (bc)_C.
\]
Set
\[
 G_u=\operatorname{rad}((ac)_A(bd)_D),\qquad
 G_t=\operatorname{rad}((ad)_B(bc)_C).
\]
Then \(G_uG_t=\operatorname{rad}(M)\).  In the \(u\)-chart, the supported
prime-power parts are \((ac)_A\) and \((bd)_D\).  Their radicals have
product \(G_u\).  In the \(t\)-chart the corresponding parts are
\((ad)_B\) and \((bc)_C\).  Their radicals have product \(G_t\).  Put
\[
 \Pi=\min(\operatorname{rad}(d_f),v_{\rm sel})
      \min(\operatorname{rad}(d_g),w_{\rm sel}).
\]
In \(4H\), \(v_{\rm sel},w_{\rm sel}>Y^p\) and \(\Pi\le Y^p\).  Hence
neither minimum can equal \(v_{\rm sel}\) or \(w_{\rm sel}\).  So
\(\Pi=\operatorname{rad}(d_f)\operatorname{rad}(d_g)\le Y^p\).  This is the
selected chart's assigned-corner radical product.  Since
\(\operatorname{rad}(M)>Y^{2/3-k}\), the opposite assigned-corner product
satisfies
\[
 G_{\rm opp}>R_0:=Y^{2/3-k-p}.                            \AutoTag{eq:old:6_9}
\]
The two-sided prime assignment in the mask makes \(G_{\rm opp}\) the
complementary assigned-corner product.

We use the following published theorem of Bombieri and Schmidt
\cite{BS1987}.  Let \(F\in\mathbb Z[X,Z]\) be an irreducible binary form
of degree \(r\ge3\), and let \(N\ge1\).  Then the number of primitive
integer pairs \((x,z)\) with \(|F(x,z)|=N\) is at most
\(C_{\rm BS} r^{1+\omega(N)}\), with an absolute \(C_{\rm BS}\).  Here
primitive means \(\gcd(x,z)=1\).  This is the unnumbered theorem on
pp.~69--70 of \emph{Inventiones Mathematicae} 88 (1987), 69--81, with the
correction on p.~445 of \emph{Inventiones Mathematicae} 97 (1989)
\cite{BS1989}. Its
bound is uniform in the coefficients.  It has no discriminant-coprimality
or coefficient-height hypothesis.

Hence, for each fixed \(C_0>0\) and every \(\eta>0\), uniformly for
positive integers \(c_1,c_2\le Y^{C_0}\) and
\(0<|h|\le Y^{C_0}\),
\[
 \#\{x,z>0:c_1x^3-c_2z^3=h\}\ll_{C_0,\eta}Y^\eta,
\AutoTag{eq:old:6_10}
\]
For the proof, let \(c=\gcd(c_1,c_2)\).  If \(c\nmid h\) there are no
solutions.  Otherwise divide by \(c\).  This gives coprime positive
coefficients \(c'_1,c'_2\) and \(h'=h/c\ne0\).  For each
\(g=\gcd(x,z)\), write \(x=gx_0,z=gz_0\).  Then \(g^3\mid h'\),
\((x_0,z_0)=1\), and
\(c'_1x_0^3-c'_2z_0^3=h'/g^3\).  If this cubic form is irreducible,
Bombieri--Schmidt gives at most
\(C3^{1+\omega(|h'|/g^3)}\) primitive pairs for that \(g\).  There are at
most \(\tauDiv(|h'|)\) choices of \(g\), and
\(3^{\omega(n)}\le\tauDiv(n)^{\log_2 3}\).  The divisor bound
\(\tauDiv(n)\ll_\rho n^\rho\), with \(\rho\) chosen after \(\eta\), makes
their total \(O_{C_0,\eta}(Y^\eta)\).

If the normalized cubic is reducible, \(c'_1/c'_2\) is a rational cube.
Coprimality then forces \(c'_1=A_{\rm cube}^3,c'_2=B_{\rm cube}^3\) for
coprime positive integers \(A_{\rm cube},B_{\rm cube}\).  Put
\(X=A_{\rm cube}x,Z=B_{\rm cube}z\).  Then
\[
 (X-Z)(X^2+XZ+Z^2)=h'.
\]
The nonzero integer $X-Z$ is a signed divisor of $h'$.  For each such
divisor the resulting quadratic has at most two integer roots.  The
divisor bound absorbs the content, gcd and reducible branches into
$Y^\eta$.  Here \(h'\ne0\), so the right side is nonzero.

For the opposite-chart conic count, write
\[
\begin{aligned}
 r_{\rm opp}&=\operatorname{rad}(d_{f,\rm opp}),&
 q_{\rm opp}&=\operatorname{rad}(d_{g,\rm opp}),\\
 d_{1,\rm opp}&=d_{f,\rm opp}/r_{\rm opp},&
 d_{2,\rm opp}&=d_{g,\rm opp}/q_{\rm opp},\\
 J_{\rm opp}&=r_{\rm opp}j_1,&K_{\rm opp}&=q_{\rm opp}z_1.
\end{aligned}
\]
The actual conic becomes the integral equation
\[
 (w_{\rm opp}d_{2,\rm opp}z_1^2)q_{\rm opp}^3
 -(v_{\rm opp}d_{1,\rm opp}j_1^2)r_{\rm opp}^3=s>0.
\AutoTag{eq:old:6_11}
\]
Here \(s\) is the shift of the opposite, counted chart. The selected
shift is then \(L/s\). The opposite chart has data
\[
(d_{f,\rm opp},d_{g,\rm opp},v_{\rm opp},w_{\rm opp})
\]
These are its supported parts and complementary factors from
\eqref{eq:old:2_4}. They differ from \((v_{\rm sel},w_{\rm sel})\) in
the \(4H\) predicate. Suppose the endpoint squares are at most
\(H_{4H}\) and \(r_{\rm opp}q_{\rm opp}\ge R_0\).  Fix
\(s,v_{\rm opp},w_{\rm opp},d_{1,\rm opp},d_{2,\rm opp},j_1,z_1\).
Apply \eqref{eq:old:6_10} to dyadic boxes
\(r_{\rm opp}\asymp R\), \(q_{\rm opp}\asymp Q\).  There are only
\(O(\log^2Y)\) such boxes, and \(RQ\ge R_0/4\).  The endpoint conditions
yield
\[
\begin{aligned}
 j_1&\le\sqrt{H_{4H}/(v_{\rm opp}d_{1,\rm opp}R^3)},\qquad
 z_1\le\sqrt{H_{4H}/(w_{\rm opp}d_{2,\rm opp}Q^3)},\\
 v_{\rm opp}w_{\rm opp}d_{1,\rm opp}d_{2,\rm opp}
 r_{\rm opp}q_{\rm opp}&<X_{\rm cap}.
\end{aligned}
\]
The displayed coefficient product is
\(v_{\rm opp}w_{\rm opp}d_{1,\rm opp}d_{2,\rm opp}r_{\rm opp}q_{\rm opp}=M\).
Failure of the \(O_R\) cut and \eqref{eq:old:2_7} give
\(M<\kappa_Y^2Y^{q_+}=\kappa_Y^2Y^{2/3+2e}\).  Define
\(X_{4H}:=\kappa_Y^2Y^{2/3+2e}\) here and take \(X_{\rm cap}=X_{4H}\).
If either upper bound is below one, that class is empty.  Otherwise put
\(N_{\rm opp}=v_{\rm opp}w_{\rm opp}\).  Sum over its at most
\(\tauDiv(N_{\rm opp})\) factorizations.  The number of \((j_1,z_1)\) is
at most the product of their upper bounds, so
\[
 \frac{H_{4H}}{\sqrt{d_{1,\rm opp}d_{2,\rm opp}}(RQ)^{3/2}}
 \sum_{N_{\rm opp}<X_{\rm cap}/(d_{1,\rm opp}d_{2,\rm opp}RQ)}
 \frac{\tauDiv(N_{\rm opp})}{\sqrt {N_{\rm opp}}}
 \ll \frac{H_{4H}\sqrt{X_{\rm cap}}\log(2X_{\rm cap})}
 {d_{1,\rm opp}d_{2,\rm opp}(RQ)^2}.
\]
Here we use $\sum_{N\le Z}\tauDiv(N)N^{-1/2}\ll\sqrt Z\log(2Z)$.  It
follows from $\tauDiv(N)=\sum_{ab=N}1$ by summing first over $b$.  Summing
\((d_{1,\rm opp},d_{2,\rm opp})\) harmonically and over the dyadic radical
boxes costs only \(Y^\eta\).  Since \(RQ\ge R_0/4\), the fixed-shift count
is
\[
 N_s\ll_\eta H_{4H}\sqrt{X_{4H}}R_0^{-2}Y^\eta.\AutoTag{eq:old:6_12}
\]
Fix an outer tuple
\(s,v_{\rm opp},w_{\rm opp},d_{1,\rm opp},d_{2,\rm opp},j_1,z_1\).
Then the cubic coefficients in \eqref{eq:old:6_11} are the fixed positive
integers \(c_1=w_{\rm opp}d_{2,\rm opp}z_1^2\) and
\(c_2=v_{\rm opp}d_{1,\rm opp}j_1^2\).
The cubic unknowns are \(q_{\rm opp}\) and \(r_{\rm opp}\).  So
\eqref{eq:old:6_10} applies uniformly in every dyadic radical box.
The coefficients may vary between outer tuples.  The uniformity of
\eqref{eq:old:6_10} permits the preceding sum over those tuples and boxes.
All canonical-source filters are kept on reconstruction.

In the following displays, \(J,K,s\) denote the counted chart and its
shift; this chart may be opposite to the first-match chart. The variables
in \eqref{eq:old:6_11} come from the canonical source. In the
$u$-chart they are
\[
 (f,g,J,K,s)=(ac,bd,A,D,u),
 \qquad (d_{f,\rm opp},d_{g,\rm opp},v_{\rm opp},w_{\rm opp})
 =(f_J,g_K,f/f_J,g/g_K),
\]
and in the $t$-chart they are
\[
 (f,g,J,K,s)=(ad,bc,B,C,t),
 \qquad (d_{f,\rm opp},d_{g,\rm opp},v_{\rm opp},w_{\rm opp})
 =(f_J,g_K,f/f_J,g/g_K).
\]
The two ordered residual factorizations are the only inverse choices.  In
the $u$-chart, fix a solution of \eqref{eq:old:6_11} and choose $ac=f$ and
$bd=g$.  Then
\[
 B=\frac{cA+bD}{u},\quad C=\frac{aA+dD}{u},\quad
 t=\frac{cd-ab}{u},\quad m=ABa,\quad n=ACc.                 \AutoTag{eq:old:6_12a}
\]
In the $t$-chart choose $ad=f$ and $bc=g$.  Then
\[
 A=\frac{dB-bC}{t},\quad D=\frac{cC-aB}{t},\quad
 u=\frac{cd-ab}{t},\quad m=ABa,\quad n=ACc.                 \AutoTag{eq:old:6_12b}
\]
Only positive integral reconstructions passing the original mask, gcd,
shell and first-match predicates are kept.  Hence
$\tauDiv(f)\tauDiv(g)$ bounds the inverse multiplicity, and $ut=L$
determines the other shift.

For the selected chart, $ut=L$ determines the opposite shift.
The source height and $R>Y^{5/3-e}$ give
\[
U=\kappa_Y Y^{1/3+e-a_0},\qquad
H_{4H}=12\kappa_Y^2Y^{4/3+e-a_0},\quad R_0=Y^{2/3-k-p},
\]
where \(\kappa_Y=1+Y^{-1}\le3/2\) for \(Y\ge2\).  If \(t<u\) is selected,
count the opposite $u$ chart with $H=uY$.  Then sum the actual $u<U$ once.
If $u\le t$ is selected, count the opposite $t$ chart with endpoint cap
$H_{4H}$.  Then sum the actual $t<U$ once; the other shift is forced.  For
the endpoint bound in the $t$-chart, the canonical identities give
\[
 (ad)B^2=\frac{m(n+1)}{ AD},\qquad
 (bc)C^2=\frac{n(m+1)}{ AD},\qquad
 (bc)C^2-(ad)B^2=t.                                      \AutoTag{eq:old:6_13a}
\]
The high-height condition \(ABCD>Y^{5/3-e}\) and
\(n(n+1)\le\kappa_Y Y^2\) give
\[
 cd=\frac{n(n+1)}{ ABCD}<\kappa_Y Y^{1/3+e}.
\]
Now \(BC=(m+n+1)/u\) and \(ABCD\,cd=n(n+1)\).  Also
\(n>Y/2\) and \(m+n+1\le3Y\).  Hence
\[
 AD=\frac{n(n+1)u}{ cd(m+n+1)}
   >\frac{uY^{2/3-e}}{12\kappa_Y}.                            \AutoTag{eq:old:6_13b}
\]
Thus both \(t\)-chart endpoint squares in \eqref{eq:old:6_13a} are at most
\(12\kappa_Y^2Y^{4/3+e}/u\).  On the subrow \(u>Y^{a_0}\) this is at most
\(H_{4H}\).  So this endpoint estimate uses both the high-height cut and
the subrow lower bound for \(u\).  The two counts are
\[
 \#4H_t\ll U^2Y\sqrt{X_{4H}}R_0^{-2}Y^\eta,\qquad
 \#4H_u\ll U H_{4H}\sqrt{X_{4H}}R_0^{-2}Y^\eta.
\AutoTag{eq:old:6_13}
\]
Here \(4H_t\) is the population selected in the \(t\)-chart
(\(t<u\)) and counted through the opposite \(u\)-chart.  Similarly,
\(4H_u\) is selected in the \(u\)-chart (\(u\le t\)) and counted through
the opposite \(t\)-chart.  The factor
\(\sqrt{X_{4H}}=\kappa_Y Y^{1/3+e}\) is the square-root factor in the
fixed-shift cubic count \eqref{eq:old:6_12}.  Substitute \(U\),
\(H_{4H}\) and \(R_0\) from above into either expression.  Both exponents
equal
\[
 \frac{2}{3}-2a_0+3e+2k+2p,
\]
so the saving is
\[
 r_H=2a_0-3e-2k-2p=\frac{140441}{5000000}>S_0.             \AutoTag{eq:old:6_14}
\]
The content split, reducible cubic count, residual divisor factors and
inverse formulas each take an arbitrary power loss.  Their number is fixed.
Summation over the shifts is the factor displayed in \eqref{eq:old:6_13}.
\begin{rowproposition}[4H row estimate]\label{prop:row-4h}
For every \(\eta>0\), the complete opposite-shift row satisfies
\[
 \#4H\ll_\eta Y^{2/3-140441/5000000+\eta}.
\]
\end{rowproposition}
\begin{proof}
The two endpoint counts in \eqref{eq:old:6_13} have the common exponent of
\eqref{eq:old:6_14}.  Their finite sum keeps that exponent.
\end{proof}

\subsection{The simultaneous-shift row $4P$}

The count in this row is simultaneous in both shifts.  Its source keeps
$v,w>Y^p$ and the strict defect bound $\Pi>Y^p$.  It also keeps the full
canonical source and all three high cuts.  It keeps the mask, unit and gcd
conditions and the selected-chart rule.  It keeps the outer and \(Q1/Q2\)
complements and the open middle-shift interval.  The 4P estimate below applies to this source.
Section~\ref{sec:long} counts the later $5W$ row separately.

In a dyadic box write
\[
 \begin{aligned}
 f&\asymp E,&g&\asymp F,&J&\asymp H_J,&K&\asymp H_K,\\
 S_h\le s<2S_h,&H_{4P}=2S_h,&X_{4P}&=H_J^2E,&V_*&=X_{4P}/H_K^2,\\
 N&=\lceil Y^{16/25}\rceil.
 \end{aligned}
\]
Here \(E\) and \(F\) are integer powers of two.  So \(EF\) is an integer
common denominator for the coefficient-lattice map below.  We also put
\[
 \begin{aligned}
 q&=\frac{2}{3}+2e,\qquad h_{\rm HP}=\frac{2}{3}+l+b_0,\\
 c&=1+e+l+b_0-2(1/3-j),\qquad
 \tau=\tau_0-\frac{1}{10000},\qquad \tau_0=1-\frac{3e}{2}.
 \end{aligned}                                          \AutoTag{eq:old:6_15}
\]
The same enlarged source gives $EF\ll Y^q$, $H_JH_K\ll Y^{h_{\rm HP}}$,
$E,F\ll Y^c$, and thickness $Y^{-\tau_0}$.  For
$z=J/K$, $w=g/f$, the conic is
\[
 gK^2-fJ^2=s.
\]
In particular \(H_{4P}/X_{4P}\ll Y^{-\tau_0}\), uniformly in the dyadic
box.  This is the thickness estimate used below.
For fixed nonnegative degrees \(A_{\rm det},B_{\rm det}\), use the
bidegree-\((A_{\rm det},B_{\rm det})\) columns
\[
 \{x^r y^j:0\le r\le A_{\rm det},\ 0\le j\le B_{\rm det}\},\qquad
 x=f/g,\quad y=K/J.                                      \AutoTag{eq:old:6_16}
\]
Clearing each row by \(g^{A_{\rm det}}J^{B_{\rm det}}\) gives the
integer entries in \eqref{eq:old:6_17c}.
The normalized determinant has the strict criterion
\[
\mu_{\rm HP}>\frac{9}{8}\frac{h_{\rm HP}q}{\tau},\qquad
\mu_{\rm HP}=\frac{16}{25}.
\AutoTag{eq:old:6_17}
\]
The statement of Lemma~\ref{lem:ratio-cell-determinant} fixes the 5W
parameters in \eqref{eq:old:7_2}--\eqref{eq:old:7_6}.  For 4P we use the
determinant calculation in its proof, with the separate 4P parameters.
The 4P grid exponent is \(a=\mu_{\rm HP}=16/25\), and
\[
 \frac35<a<\min(q_{\min},\tau_0),\qquad
 q_{\max}<\tau_0,\qquad b_0>\frac32e.
\]
The bounds \(H_J\le Y+1\) and \(E\le(Y+1)^2\) follow from the source
identities \eqref{eq:old:2_2}.  The 4P row itself has
\(H_{4P}>Y^{a_0}\).  The failed \(Q1\) and \(Q2\) cuts and the shell lower
bound \eqref{eq:old:2_8} supply the remaining endpoint restrictions in the
4P boxes.  In particular these source bounds give
\(X_{4P}/H_{4P}\le Y^4\) for sufficiently large \(Y\).  The Taylor order
is seven because \(8\mu_{\rm HP}=8(16/25)=128/25>4\).  With
\[
 p_1=\frac q{\mu_{\rm HP}}=\frac{5092607}{4800000}>1,
 \qquad
 p_2=\frac{h_{\rm HP}}\tau=\frac{24536965}{29719179},
\]
we have \(2/3<p_2<1\), \(p_1p_2<8/9\), and
\(p_1-1<3(1-p_2)^2\).  These include the strict trapezoid-branch hypotheses in
the proof of the general determinant lemma.  They give its 4P
specialization
\((a,p_1,p_2)=(16/25,5092607/4800000,24536965/29719179)\).
The determinant criterion needs these strict range and strip-weight
hypotheses as well as its displayed inequality.  The exact collared
determinant margin is positive; it is recorded just below.
Here \(\rho=p_2=h_{\rm HP}/\tau>2/3\), so 4P uses the trapezoid branch.
The stronger sufficient condition \(\mu_{\rm HP}>(9/8)q\rho\)
implies its criterion for \(2/3\le\rho<1\), because
\[
\frac98\rho\bigl(1+3(1-\rho)^2\bigr)-1
=\frac{(1-3(1-\rho))^3}{8}\ge0.
\]
At \eqref{eq:old:2_5}, the collared gap
\(\mu_{\rm HP}-(9/8)q\rho\) equals
$368942136449/39625572000000>0$.  Fix half this strict gap, then choose
the finite degree family.  All maximal-size integer minors have absolute
value less than one.  Hence they vanish, and the integerized matrix
\(\mathcal M\) has rank \(r<D\), where \(D\) is its number of columns.
If an interval contains fewer than \(D\) source rows, count those rows
directly; otherwise the vanishing \(D\)-minors imply rank below \(D\).
If \(r=0\), a coordinate vector is in its kernel.  Otherwise choose \(r\)
pivot columns and \(r\) rows giving a nonsingular minor \(A\), and one
free column \(j\).  Set the \(j\)-coordinate to \(\det A\), the other
free coordinates to zero, and the pivot coordinates to
\(-\operatorname{adj}(A)\mathcal M_{I,j}\).  This nonzero integer vector
annihilates the selected rows.  They span the row space, so it annihilates
every row.  Its entries are minors of bounded order and have polynomial
height.  The same argument covers the fixed anchor, the inverse monomials
and every endpoint branch.

For the normalization, use in each box the exact half-open ratio cells
\[
 r_i=\frac{i}{2N}\frac{H_J}{ H_K},\qquad
 I_i=(r_i,r_{i+1}],\qquad i=N,\ldots,4N-1 .                 \AutoTag{eq:old:6_17a}
\]
For a point in \(I_i\), put \(x=f/g\), \(y=K/J\), \(y_0=r_i^{-1}\),
\(\delta_{\rm cell}=s/(fJ^2)\), and \(\alpha_{\rm cell}=x/y_0^2-1\).
The conic gives the exact
identities
\[
 x=\frac{y^2}{1+\delta_{\rm cell}},\qquad
 \frac{y}{ y_0}=\sqrt{1+\alpha_{\rm cell}}\sqrt{1+\delta_{\rm cell}}.
\]
The cell width and \(s\le H_{4P}\) give
\[
 |\alpha_{\rm cell}|\le \frac{3}{ N},\qquad
 0\le \frac{y}{y_0}-\sqrt{1+\alpha_{\rm cell}}\le \frac{H_{4P}}{2X_{4P}}.       \AutoTag{eq:old:6_17b}
\]
Here \(r_{i+1}/r_i=1+1/i\le1+1/N\) and
\(1-(1+\delta_{\rm cell})^{-1/2}\le\delta_{\rm cell}/2\).  Since
\(X_2=X_{4P}/(2H_{4P})\), the conic error in \eqref{eq:old:6_17b} is
\(H_{4P}/(2X_{4P})=X_2^{-1}/4\).  Also
\(NH_{4P}/X_{4P}\le Y^{\mu_{\rm HP}-\tau_0+o(1)}<1\).  Choose a fixed
Taylor degree \(T\) and write
\(\sqrt{1+\alpha_{\rm cell}}=1+P_T(\alpha_{\rm cell})+\varepsilon_T\).
With \(X_1=N/3\), \eqref{eq:old:6_17b} gives
\(|\alpha_{\rm cell}|\le X_1^{-1}\). Choose \(T\)
large enough that the Taylor remainder is at most \(X_2^{-1}/2\), where
\(X_2=X_{4P}/(2H_{4P})\). Then the combined remainder
\(\beta=(y/y_0)-1-P_T(\alpha_{\rm cell})\) satisfies \(|\beta|\le X_2^{-1}\),
uniformly over all boxes and shifts.

Take columns \(x^r y^j\), \(0\le r\le A_{\rm det}\),
\(0\le j\le B_{\rm det}\), and multiply each row by
\(g^{A_{\rm det}}J^{B_{\rm det}}\).  The resulting entries are the integers
\[
 f^r g^{A_{\rm det}-r}K^jJ^{B_{\rm det}-j}.             \AutoTag{eq:old:6_17c}
\]
Write the polynomial in \eqref{eq:old:6_16} as
\[
C(x,y)=\sum_{r=0}^{A_{\rm det}}\sum_{j=0}^{B_{\rm det}}c_{rj}x^ry^j.
\]
With \(z=J/K=1/y\) and \(w=g/f=1/x\), its reciprocal is
\[
F(z,w)=z^{B_{\rm det}}w^{A_{\rm det}}C(1/w,1/z)
 =\sum_{r,j}c_{rj}z^{B_{\rm det}-j}w^{A_{\rm det}-r}.
\]
It has degrees at most \((B_{\rm det},A_{\rm det})\) in \((z,w)\) and the
same coefficient height.  Its coefficient vector permutes the full
rectangular coefficient array.  The column map is
\((r,j)\mapsto(a,b)=(B_{\rm det}-j,A_{\rm det}-r)\).  In particular,
\[
K^{B_{\rm det}}f^{A_{\rm det}}F(J/K,g/f)
 =J^{B_{\rm det}}g^{A_{\rm det}}C(f/g,K/J),
\]
so the two integer matrices agree after clearing denominators, up to that
column permutation.
After extracting the common anchor, the normalized entries are
\((1+\alpha_{\rm cell})^r(1+P_T(\alpha_{\rm cell})+\beta)^j\).  If
\(D_c=(A_{\rm det}+1)(B_{\rm det}+1)\),
the anchor and row-clearing factors have logarithmic cost at most
\[
 \frac{D_c}{2}\{A_{\rm det}\Lambda_1+B_{\rm det}\Lambda_2\},\qquad
 \Lambda_1=\log(324EV_*),\quad\Lambda_2=\log(8H_JH_K),
\AutoTag{eq:old:6_17d}
\]
The column-anchor factor is \(y_0^{2r+j}\).  Its exponent sum is
\[
\sum_{\substack{0\le r\le A_{\rm det}\\0\le j\le B_{\rm det}}}(2r+j)
=D_c\left(A_{\rm det}+\frac{B_{\rm det}}2\right).
\]
The other anchor bound is
\[
J y_0^{1/2}<2H_J(2H_K/H_J)^{1/2}=(8H_JH_K)^{1/2}.
\]
Also \(V_*=E H_J^2/H_K^2\), so
\[
18V_*\frac{H_K}{H_J}=18\sqrt{EV_*},\qquad
(18\sqrt{EV_*})^2=324EV_*.
\]
These inequalities give the constants \(8\) and \(324\) in the uniform
anchor bounds.
The conic and the thickness give \(EV_*\asymp EF\).  So replacing one by
the other changes only an absolute logarithmic constant.  The displayed
cost follows from \(g<9V_*\), \(y_0\le2H_K/H_J\), and \(J<2H_J\).
It is independent of \(i\) and \(s\).

The proof of Lemma~\ref{lem:ratio-cell-determinant} derives the finite
Cauchy--Binet expansion, the strip-weight lower bound, the branch
integrals and the integer-minor conclusion.  It does so for general
logarithmic parameters satisfying its strict inequalities.  We apply that
calculation with the 4P parameters and effective scales below.
For this row put
\[
 X_1=N/3,\qquad X_2^{\rm eff}=Y^\tau,
\]
and choose fixed constants \(C_1,C_2\) such that, uniformly in every
box,
\[
 \Lambda_1^{\rm eff}=q\log Y+C_1\ge\log(324EV_*),\qquad
 \Lambda_2^{\rm eff}=h_{\rm HP}\log Y+C_2\ge\log(8H_JH_K).
\]
Such constants exist by the source bounds \(EV_*\ll EF\ll Y^q\) and
\(H_JH_K\ll Y^{h_{\rm HP}}\).  Also \(H_{4P}/X_{4P}\ll Y^{-\tau_0}\), where
\(\tau_0>\tau\).  So for all sufficiently large \(Y\),
\(X_2=X_{4P}/(2H_{4P})\ge X_2^{\rm eff}\).  Thus the actual Taylor error
\(|\beta|\le X_2^{-1}\) is at most \((X_2^{\rm eff})^{-1}\).  The
arithmetic cost in \eqref{eq:old:6_17d} is bounded by the one obtained from
\(\Lambda_1^{\rm eff},\Lambda_2^{\rm eff}\).  With
\[
 \ell_1=\log X_1,\quad \ell_2^{\rm eff}=\tau\log Y,\quad
 P_{\rm eff}=\frac{\Lambda_1^{\rm eff}}{\ell_1},\quad
 R_{\rm eff}=\frac{\Lambda_2^{\rm eff}}{\ell_2^{\rm eff}},\quad
 \ell_{\rm det}=\frac{\ell_1}{\ell_2^{\rm eff}},
\]
we have
\[
 P_{\rm eff}=\frac q{\mu_{\rm HP}}+O(1/\log Y),\quad
 R_{\rm eff}=\frac{h_{\rm HP}}\tau+O(1/\log Y),\quad
 \ell_{\rm det}=\frac{\mu_{\rm HP}}\tau+O(1/\log Y).
\]
At the limiting 4P vector,
\[
 \mu_{\rm HP}=\frac{16}{25},\qquad
 p_1=\frac{5092607}{4800000},\qquad
 p_2=\frac{24536965}{29719179},\qquad
 \ell_{\rm det}=\frac{6400000}{9906393}.
\]
The corresponding normalized determinant margins are
\[
 \frac89-p_1p_2=\frac{368942136449}{28530411840000}>0,\qquad
 3(1-p_2)^2-(p_1-1)
 =\frac{14253084474184329457}{471055786898155200000}>0.
\]
Hence the strict margin in \eqref{eq:old:6_17}, with the fixed collar
\(\tau=\tau_0-1/10000\), applies to these uniform majorants
for all sufficiently large \(Y\).  The finite determinant calculation uses
only the displayed upper bounds for the two errors and the arithmetic
cost.  So it gives one degree choice and one polynomial-height bound,
uniformly in the cell, shift and box.

After bihomogenization and factorization, we count each factor.
A factor that is not absolutely irreducible meets a distinct conjugate
in only boundedly many rational projective points at fixed degree.  For
corner-degree-zero factors, first fix the primitive coefficient pair,
then its ordered residual factors.  These fix $L=cd-ab$.  Discard $L\le0$
and enumerate the divisor choices $s\mid L$.  For each fixed
$(f,g,s)$ apply Theorem~\ref{thm:fibre} with
$(a,b,h,n,m)=(f,g,s,J,K)$, $x=2Y$, and fixed height exponent
$A_{\rm ht}=2$.
The residuals are at most $Y+1$, whence $f,g,s\le(Y+1)^2$ and
$J,K\le Y+1$.  This costs $Y^\eta$ per coefficient label.  The selected
shift is enumerated among the divisors of \(L\).  Explicitly, this cost
consists of three parts.  They are the ordered residual splitting factors
\(\tauDiv(f)\tauDiv(g)\), the selected-shift choices \(\tauDiv(L)\), and
the uniform square-fibre bound for each fixed \((f,g,s)\).  Give each
operation its own small enough internal exponent.  Then their product is
\(Y^\eta\).

For a coefficient-degree-zero factor, fix its primitive positive corner
pair $(J,K)$.  For two source coefficient pairs the conic gives
\[
 K^2(g_1f_2-g_2f_1)=s_1f_2-s_2f_1,
 \qquad |g_1f_2-g_2f_1|
 \le \frac{4H_{4P}E}{H_K^2}
 \ll EF\frac{H_{4P}}{X_{4P}}
 \ll Y^{q-\tau_0}=o(1).
\AutoTag{eq:old:6_18}
\]
Here $s_i<H_{4P}$, $f_i<2E$, $K\ge H_K$, $X_{4P}=H_J^2E$, and
 $V_*=X_{4P}/H_K^2\ll F$ by the source box bounds.  The cross determinant
is an integer.  It vanishes for all sufficiently large $Y$, since
$\tau_0-q=9351751/30000000>0$.  Primitivity of $(f_i,g_i)$ then forces
the coefficient pairs to agree, and the conic fixes the shift.  This
branch contributes \(O(N)\) primitive conic-coordinate tuples across the
primary cells.  To return to complete canonical source records, include
the ordered residual factorizations.  They cost
\(\tauDiv(f)\tauDiv(g)\ll_\eta Y^\eta\), so this branch is
\(O_\eta(Y^\eta N)\).

Let the coefficient degree be $b_q\ge1$ and the corner degree $a_c\ge2$.
Then apply Heath--Brown's bihomogeneous Lemma~2~\cite{HB2012} to the
factor after the unimodular coefficient-lattice change of basis.  Its
bidegree is $(b_q,a_c)$.  The first pair is the primitive transformed
coefficient pair of radius $O(T_1+1)$.  The second pair is the original
primitive corner pair $(J,K)$, with no size hypothesis.  The factor height
is $Y^C$ for a fixed $C$.  Thus the count is
$O((T_1+1)^{2/a_c+\varepsilon}Y^{C\varepsilon})$.  Here $2/a_c\le1$ and
$T_1+1\le Y^{C'}$ for a fixed $C'$.  Choose the theorem's internal
$\varepsilon<\eta/(C+C'+1)$; this bounds the count by $Y^\eta(T_1+1)$.
For corner degree one write the factor as
$KP(f,g)+JQ(f,g)$ with coprime homogeneous $P,Q$.  Their nonzero resultant
and homogeneous Bezout identities give
$d_0=\gcd(P(f,g),Q(f,g))\mid\operatorname{Res}(P,Q)$, since
$\gcd(f,g)=1$.  The factor equation forces
$(J,K)=\pm(P(f,g)/d_0,-Q(f,g)/d_0)$, with at most one positive sign.
Write $f_J=d_f$, $g_K=d_g$ for the supported parts and
$v=f/d_f$, $w=g/d_g$.
Write \(\ell_f,\ell_g\) for the endpoint labels in the four patterns
\((P(0,1),Q(1,0))\):
\[
\begin{array}{c|c|c|c}
P(0,1)&Q(1,0)&\ell_f&\ell_g\\ \hline
\ne0&\ne0&\rad(d_f)&\rad(d_g)\\
\ne0&=0&\rad(d_f)&w\\
=0&\ne0&v&\rad(d_g)\\
=0&=0&v&w .
\end{array}
\AutoTag{eq:old:6_18a}
\]
If $P(0,1)\ne0$, reducing $d_0J=P(f,g)$ modulo $\rad(d_f)$ gives
$\rad(d_f)\mid P(0,1)$.  If $P(0,1)=0$, then $f\mid P(f,g)$, and
$\gcd(v,J)=1$ gives $v\mid d_0\mid\operatorname{Res}(P,Q)$.  The same
argument with $Q(1,0)$ proves the two $g$-side alternatives.  The chosen
endpoint labels satisfy \(\ell_f\ell_g\ge\Pi>Y^p\), since each is at
least the corresponding minimum in the definition of \(\Pi\).  Each label
divides a fixed nonzero endpoint coefficient or resultant of polynomial
height.  So there are \(\ll_\eta Y^\eta\) possible pairs of endpoint
labels, for any prescribed $\eta>0$.  Also \(\gcd(\ell_f,\ell_g)=1\),
since \(\gcd(f,g)=1\) and the labels are supported on the respective parts
of \(f\) and \(g\).  Take two distinct primitive coefficient pairs with
the same labels.  Their divisibilities combine, and the nonzero cross
determinant is divisible by \(\ell_f\ell_g\).  Take a coefficient cell of
width $O(F/(EN))$ in $g/f$.  There the ratios are separated by at least
$\ell_f\ell_g/(4E^2)$, since $f_i<2E$.  Thus the number of pairs in one
cell is
\[
 O\left(1+\frac{EF}{ N \ell_f\ell_g}\right).
\]
Summing over the $O(N)$ coefficient cells and the divisor labels gives
\[
 \ll_\eta Y^\eta\left(N+\frac{EF}{ \ell_f\ell_g}\right)
 \ll_\eta Y^\eta\left(N+\frac{EF}{ Y^p}\right),
\]
including unit labels and endpoint equalities. A factor equal to a source
coordinate vanishes only when that coordinate is zero, so it has no points
in the positive source. This is separate from the lattice shortest-vector
axes, which are retained.
This branch counts the primitive coefficient and corner tuples.
Recovering complete source pairs still needs the ordered residual
splittings of \(f\) and \(g\), at most \(\tauDiv(f)\tauDiv(g)\).  That
multiplicity is absorbed in the displayed \(Y^\eta\).

For the high-degree factors use a second, coefficient-ratio grid.  Write
\[
 r_{i_c}=\frac FE\frac{i_c}{N},\qquad
 C_{i_c}=(r_{i_c},r_{i_c+1}],\qquad
 \lfloor N/2\rfloor\le i_c<2N.
\AutoTag{eq:old:6_19a}
\]
The dyadic source has \(g/f>F/(2E)\), since \(g\ge F\) and \(f<2E\).
Thus these \(O(N)\) cells cover its full range
\(F/(2E)<g/f<2F/E\).  Now use
\[
 \frac{g}{ f}=\left(\frac{J}{ K}\right)^2+\frac{s}{ fK^2},
 \qquad \frac{s}{gK^2}\ll \frac{H_{4P}}{X_{4P}}\ll Y^{-\tau_0},
\AutoTag{eq:old:6_19b}
\]
and $\tau_0>\mu_{\rm HP}$.  The relative displacement from the square map is
then $o(1/N)$.  Each grid has relative width $O(1/N)$, and squaring
preserves relative widths on positive intervals.  Hence each coefficient
cell meets $O(1)$ primary $J/K$ cells, and conversely.  Thus the cell
polynomials from \eqref{eq:old:6_17a} give $O(1)$ factors per coefficient
cell.

For a coefficient cell define the rational linear map
\[
 \mathcal T_{i_c}(g,f)=
 \left(\frac{N}{ F}(g-r_{i_c}f),\frac{f}{ E}\right).
\AutoTag{eq:old:6_19c}
\]
For a point in this cell, $f/E\in[1,2)$ and
\[
0\le \frac{N}{ F}(g-r_{i_c}f)
 <\frac{N}{ F}\,2E\,\frac{F}{ EN}=2,
\]
so its image lies in $[0,2]\times[1,2]$.  The matrix of
$\mathcal T_{i_c}$ is
\[
\begin{pmatrix}N/F&-i_c/E\\0&1/E\end{pmatrix},\qquad
\det\mathcal T_{i_c}=N/(EF).
\]
Let $\Lambda_{i_c}=\mathcal T_{i_c}\mathbb Z^2$.  Choose a
Lagrange-reduced basis $v_1,v_2$ of this image lattice.  Its lengths are
$\lambda_1\le\lambda_2$, and its angle is bounded away from zero.  Write
$v_j=\mathcal T_{i_c}h_j$.  Since $v_1,v_2$ are a lattice basis,
$h_1,h_2\in\mathbb Z^2$ form a unimodular basis of $\mathbb Z^2$.
Hence every primitive coefficient pair has a primitive transformed
pair $(x_1,x_2)$ in $(g,f)=x_1h_1+x_2h_2$.  Put
$T_1=1/\lambda_1$, $T_2=1/\lambda_2$.  The determinant and reduction give
\[
 T_1\ge T_2,\qquad T_1T_2\asymp \frac{EF}{ N},\qquad
 T_1\gg\sqrt{EF/N},\qquad
 \max(|x_1|,|x_2|)\ll T_1
\AutoTag{eq:old:6_19d}
\]
for the transformed coordinates $(x_1,x_2)$.  For the last bound, the
bounded image rectangle and the angle bound imply
$|x_j|\ll1/\lambda_j\le T_1$.  The source has
$EF\gg Y^{2/3-e}$ and $N\asymp Y^{16/25}$, so $EF/N$ grows as a positive
power.  A common denominator for the matrix of $\mathcal T_{i_c}$ is
$EF$.  So every nonzero image lattice vector has length at least
$(EF)^{-1}$.  Also $\lambda_1\ll\sqrt{N/(EF)}$ by the two-dimensional
Minkowski bound, and $\lambda_1\lambda_2\asymp N/(EF)$ by reduction.
Hence $\lambda_2\ll N$.  The inverse matrix is
\[
\mathcal T_{i_c}^{-1}(y_1,y_2)
=\bigl((F/N)(y_1+i_cy_2),Ey_2\bigr).
\]
Since $i_c\asymp N$ and $E,F,N\le Y^C$ for a fixed $C$, both $v_j$
and their integer preimages $h_j$ have polynomial height.  So substitution
by the unimodular matrix $[h_1\ h_2]$ keeps every fixed-degree factor at
polynomial coefficient height.  The unimodular change uses no moving
denominator.

The census below includes axes and first points.  At scale
$D_{\rm lat}\le T_1<2D_{\rm lat}$, let $(g_{\rm lat},f_{\rm lat})$ be the
primitive preimage of the shortest basis vector.  From
\eqref{eq:old:6_19c},
\[
 |f_{\rm lat}|\ll E/D_{\rm lat},\quad |g_{\rm lat}|\ll F/D_{\rm lat},\quad
 |i_cf_{\rm lat}-(NE/F)g_{\rm lat}|\ll E/D_{\rm lat}.
\AutoTag{eq:old:6_19e}
\]
If $f_{\rm lat},g_{\rm lat}\ne0$, then $E/D_{\rm lat}$ and $F/D_{\rm lat}$
are at least a fixed positive constant.  There are $O(F/D_{\rm lat})$
choices for $g_{\rm lat}$.  For each, the integer product $i_cf_{\rm lat}$
lies in an interval of length $O(E/D_{\rm lat})$.  This gives that many
possible values.  Each nonzero value has at most
$\tauDiv(|i_cf_{\rm lat}|)\ll_\eta Y^\eta$ factorizations into
$i_c f_{\rm lat}$, since all factors have polynomial height.  Thus
\[
 \#\{i_c:T_1\asymp D_{\rm lat},\ f_{\rm lat}g_{\rm lat}\ne0\}
 \ll_\eta Y^\eta \frac{EF}{D_{\rm lat}^2}.
\AutoTag{eq:old:6_19f}
\]
If $f_{\rm lat}=0$, primitivity gives $g_{\rm lat}=\pm1$ and
$T_1\asymp F/N$.  So the sum over $O(N)$ cells is $O(F+N)$.  If
$g_{\rm lat}=0$, then $f_{\rm lat}=\pm1$, $i_c\asymp N$, and
$T_1\asymp E/N$, giving $O(E+N)$.  Cells with $T_1<1$ contribute $O(N)$.
These cells, both axes and every $+1$ are kept.  Weight
\eqref{eq:old:6_19f} by $T_1\asymp D_{\rm lat}$ and sum geometrically
from $D_{\rm lat}\gg\sqrt{EF/N}$.  This yields
\[
 \sum_{i_c}(T_1(i_c)+1)\ll_\eta
 Y^\eta\{N+\sqrt{EFN}+E+F\}.
\AutoTag{eq:old:6_19}
\]
This count factors the exact integer product \(i_cf_{\rm lat}\). The number
of cells in this dyadic range is
\(O_\eta(Y^\eta EF/D_{\rm lat}^2)\).

For a factor with $b_q\ge1$ and $a_c\ge2$, the Heath--Brown bound above
is $Y^\eta(T_1+1)$.  Summing over coefficient cells by \eqref{eq:old:6_19}
gives $Y^\eta(N+\sqrt{EFN}+E+F)$.  Short matrices, nonabsolute factors and
the two zero-degree branches contribute $Y^\eta N$.  The corner-linear
branch contributes $Y^\eta(N+EF/Y^p)$, by the endpoint-label spacing
argument above.  These cases exhaust the fixed-degree factorization.

Combine the lattice census \eqref{eq:old:6_19} with the corner-linear
count.  The complete bound then has the form
\[
 \#4P\ll_\eta Y^\eta\left(N+\sqrt{EFN}+E+F+\frac{EF}{Y^p}\right).
\]
These local bounds are summed over the \(O(\log^4Y)\) choices of dyadic
scales and the \(O(\log Y)\) shift scales.  The resulting logarithmic
factor is absorbed into the arbitrary \(Y^\eta\) loss.
The four distinct exponent contributions, in the order
\(EF/Y^p\), \(\sqrt{EFN}\), \(N\), and \(E+F\), are
\[
 q-p=\frac{9907393}{15000000},\quad
 \frac{q+\mu}{2}=\frac{9892607}{15000000},\quad
 \mu=\frac{16}{25},\quad
 c=\frac{5277821}{10000000}.
\AutoTag{eq:old:6_20}
\]
Their savings relative to $2/3$ are respectively
\[
S_0,\qquad \frac{107393}{15000000},\qquad \frac{2}{75},\qquad
 \frac{4166537}{30000000}.
\AutoTag{eq:old:6_21}
\]
The minimum is $S_0=30869/5000000$.  We combine the uniform lattice census
\eqref{eq:old:6_19}, the finite factor family and the polynomial-height
estimates.  With the divisor losses and the one simultaneous source
inversion, this gives
\[
\#4P\ll_\eta Y^{2/3-S_0+\eta}.
\AutoTag{eq:old:6_22}
\]
\begin{rowproposition}[4P row estimate]\label{prop:row-4p}
For every \(\eta>0\), the complete simultaneous-shift row satisfies
\[
 \#4P\ll_\eta Y^{2/3-S_0+\eta}.
\]
\end{rowproposition}
\begin{proof}
Equations \eqref{eq:old:6_19}--\eqref{eq:old:6_21} bound every factor
branch and its exponent.  Then \eqref{eq:old:6_22} sums them with the
source-uniform lattice and inverse bounds.
\end{proof}

The $4P$ estimate applies to its enlarged source.  Section~\ref{sec:long}
counts $5W$ from its original-cell source and geometric envelope.

\section{The long-shift row and its geometric cap}\label{sec:long}

The row \(5W\) is the exact first-match remainder after \(4P\). Its source
keeps every earlier complement. In particular,
\[
 s\ge Y^{b_0},\quad K_{\rm res}\le Y^{\Delta_{\rm exp}},\quad
 \Delta_{\rm exp}=2\epsilon_0+k_0=3S_0.
\AutoTag{eq:old:7_1}
\]
The selected corners and both shifts keep their lower bounds from
Section~\ref{sec:partition}. The remainder predicate includes the exact
bound \(K_{\rm res}\le Y^{\Delta_{\rm exp}}\).

With \(e=\epsilon_0\), \(k=k_0\), \(l=l_0\), \(j=j_0\), and \(p=p_0\), write
\[
\begin{aligned}
 S_0=e=k&=\frac{30869}{5000000},&
 l&=2S_0,& j=p&=3S_0,\\
 b_0&=\frac{4166537}{30000000},&
 a_0&=\frac6{125},&
 \mu_{\rm HP}&=\frac{16}{25},\\
 \nu&=\frac1{64},&
 \xi&=\frac1{10000}.
\end{aligned}
\AutoTag{eq:old:7_2}
\]
Set
\[
 \begin{aligned}
 q_{\min}&=\frac23-S_0,\qquad q_{\max}=\frac23+2S_0,\\
 \tau_0&=1-\frac{3e}{2},\qquad
 \lambda=\log_Y(1+Y^{-1}),\qquad d_2=\log_Y2,\\
 \ell_C&=\log_Y(1+\sqrt3Y^{-\tau_0}),\qquad
 \ell_{\rm end}=\log_Y\bigl((1+\sqrt3Y^{-\tau_0})(1+2/Y)\bigr).
 \end{aligned}
\AutoTag{eq:old:7_3}
\]
The source inequalities imply \(q_{\min}<\log_Y M<q_{\max}+2\lambda\).
They also imply all endpoint comparisons used below. Each chart keeps its
own shift interval throughout the count.

\subsection{A polynomial on each ratio cell}

In a local chart box write
\[
 D_0\le J<2D_0,\quad E_0\le K<2E_0,\quad
 U_0\le f<2U_0,\quad S_h\le s<2S_h,\quad H_0=2S_h,
\]
\[
 X=D_0^2U_0,\qquad V_*=\frac{X}{E_0^2}.
\AutoTag{eq:old:7_4}
\]
These symbols are local box scales. On the actual population the source
identities give, uniformly,
\[
 \frac{X}{H_0}>\frac{Y^{\tau_0}}{96},\qquad
 \frac M{18}<U_0V_*\le4M,\qquad
 U_0\le f,\quad V_*\le4g.
\AutoTag{eq:old:7_5}
\]
The selected chart stays fixed as part of the first-match source label.
Independently, either counting orientation
\((J,K;f,g;s)=(A,D;ac,bd;u)\) or \((B,C;ad,bc;t)\) may be used
on that retained population. All constants in \eqref{eq:old:7_5} are
uniform in this orientation and its moving shift box.
Here \(T=fJ^2/s\) equals \(mn/(m+n+1)\ge m/2\) or
\(m(n+1)/(n-m)>m\), respectively. The dyadic definitions give
\(T<16X/H_0\). Hence \(X/H_0>m/32>Y^{\tau_0}/96\).
Also \(X/H_0<(Y+1)^4/Y^{b_0}<Y^4\) for sufficiently large \(Y\).
The conic gives the other comparisons in \eqref{eq:old:7_5} in
both orientations. It needs no lower bound for the unselected corners.

Put \(R_{\rm box}=D_0/E_0\), \(L_{\rm grid}=2N\), where
\[
 \bar m=\frac{q_{\max}}{1+3b_0^2}
       =\frac{203704280000000}{317360030572369},\qquad
 a_{\rm grid}=\bar m+\xi,\qquad N=\lceil Y^{a_{\rm grid}}\rceil.
\AutoTag{eq:old:7_6}
\]
The original grid consists of exactly the \(3N\) half-open cells
\[
 r_i=\frac{i}{2N}R_{\rm box},\qquad I_i=(r_i,r_{i+1}],
 \qquad i=N,\ldots,4N-1,
\AutoTag{eq:old:7_7}
\]
for the ratio \(J/K\). As \(D_0\le J<2D_0\) and
\(E_0\le K<2E_0\), these cells cover the whole local source ratio.
For \(r_0=r_i\), the relative width satisfies
\(r_{i+1}/r_i=1+1/i\le1+1/N\).

The exact determinant cap used to choose \(N\) is
\[
 \phi(\rho)=
 \begin{cases}
  9\rho/8,&0<\rho\le2/3,\\[2pt]
  \bigl(1+3(1-\rho)^2\bigr)^{-1},&2/3\le\rho<1.
 \end{cases}
\AutoTag{eq:old:7_8}
\]
Write \(q_M=\log_Y M\), \(c_0=\log_Y(D_0E_0)\), and
\(\tau=\log_Y(X/H_0)\). The limiting comparison ratio is
\(\rho=c_0/\tau\). The comparison
\(\log_YN>q_M\phi(\rho)+o(1)\) is the leading-order sufficient
condition. The proof below gives the uniform finite-\(Y\) margin after
the source and arithmetic collars. For the source comparison put
\(\sigma=\log_Ym\), \(z_s=\log_Ys\), and
\(\ell_L=\log_YL\). In the second chart also put
\(z_u=\log_Yu\). The exact identities from \eqref{eq:old:2_1}--\eqref{eq:old:2_2} used below are
\[
 AD={n-m\over t},\quad BC={m+n+1\over u},\quad
 L={(n-m)(m+n+1)\over Z},\quad Z^2M=m(m+1)n(n+1),
\]
\[
 A^2(ac)={umn\over m+n+1},\quad
 B^2(ad)={tm(n+1)\over n-m},\quad
 C^2(bc)={tn(m+1)\over n-m}.
\AutoTag{eq:old:7_8a}
\]
In the first chart, \(s=u\). The limiting source coordinates are
\[
 (c^*,\tau^*)=(\sigma+z_s-q_M/2,\sigma),
\qquad \rho_1^*={c^*\over\tau^*}
 =1+{z_s-q_M/2\over\sigma}.
\AutoTag{eq:old:7_8b}
\]
Before the change to local dyadic scales,
\[
 \log_Y(AD)-c^*
 =\log_Y{\sqrt{m(m+1)n(n+1)}\over m(m+n+1)},\quad
 \log_Y(A^2f/u)-\sigma=\log_Y{n\over m+n+1}.
\AutoTag{eq:old:7_8c}
\]
The first ratio factors as
\((1+1/m)^{1/2}[n/(m+n+1)](1+1/n)^{1/2}\). The middle
factor lies in \([1/(2+2/Y),1]\). The source gives
\(\tau_0-2d_2<\sigma\le1\). Both shifts satisfy
\(u,t\ge Y^{b_0}\), and
\(L=ut<\kappa_YY^{1/3+e}\). So for either chart shift \(s\in\{u,t\}\),
\[
 z_s\le q_{\max}/2-b_0+\lambda.
\AutoTag{eq:old:7_8b_common}
\]
On the selected chart, \(s=\min(u,t)\). So \(s^2\le L\) gives the stronger
individual estimate
\[
 z_s\le \frac16+\frac e2+\frac\lambda2.
\]
The alternative chart uses the other shift and the common bound
\eqref{eq:old:7_8b_common}.
The stronger estimate holds only for the selected smaller shift. It is
recorded for that source coordinate. It is not used in \eqref{eq:old:7_8e}.
Both chart caps use the common shift bound, which applies to either shift. The limiting
gap in this comparison is
\[
 \frac{q_{\max}}2-b_0-\left(\frac16+\frac e2\right)
 =\frac{30869}{1000000}>0.
\AutoTag{eq:old:7_8b_gap}
\]
The common shift bound minus the selected-shift bound is exactly
\[
 \frac{30869}{1000000}+\frac{\lambda}{2}>0.
\]
Thus their ordering needs no extra restriction on \(\lambda\). Small
\(\lambda\) is used below for the other strict endpoint comparisons. Its
terms tend to zero and are absorbed by their fixed margins.
Since \(q_M>q_{\min}\),
\[
 q_M/2-z_s\ge b_0-3e/2-\lambda>0
\]
for all sufficiently large \(Y\). Therefore
\[
 \frac{\partial\rho_1^*}{\partial\sigma}
 =\frac{q_M/2-z_s}{\sigma^2}>0,\qquad
 \frac{\partial\rho_1^*}{\partial z_s}=\frac1{\sigma}>0.
\]
Hence the maximum in \(\sigma\) is at \(\sigma=1\).
Fix \(z_s\). On the triangle branch
\[
 {d\over dq_M}\left[{9\over8}q_M(1-q_M/2+z_s)\right]
 ={9\over8}(1-q_M+z_s)>0.
\]
On the trapezoid branch, with \(r=q_M/2-z_s\),
\[
 {d\over dq_M}{q_M\over1+3(q_M/2-z_s)^2}
 ={1-3q_M^2/4+3z_s^2\over(1+3(q_M/2-z_s)^2)^2}>0,
\]
since \(q_M<1\). So this chart's maximum is at
\(q_M=q_{\max}+2\lambda\). At the joint upper endpoints
\(\sigma=1\), \(z_s=q_{\max}/2-b_0+\lambda\), the ratio is
\(\rho_1^*=1-b_0\). This is the alternative \(u\)-chart bound, even on a
source whose selected chart is \(t\). The comparison follows from the two
lower shift bounds and \(ut=L\).

In the second chart, \(s=t\). So the opposite shift satisfies
\(z_u\ge b_0\) on the entire \(5W\) source, and
\(\ell_L\le1/3+e+\lambda\). Its limiting coordinates are
\[
 (c^*,\tau^*)=(1-z_u,1-\ell_L+q_M/2),
 \qquad \rho_2^*={1-z_u\over1-\ell_L+q_M/2}.
\AutoTag{eq:old:7_8d}
\]
At \(\lambda=0\), increasing \(\ell_L\) and decreasing \(z_u\) increase
the ratio. So
\(\rho_2^*\le(1-b_0)/(q_{\min}+q_M/2)\). For finite \(Y\), the
denominator is at least \(D_\lambda=q_{\min}+q_M/2-\lambda\), where
\(D=q_{\min}+q_M/2\). With \(c=1-b_0\), the triangle-branch derivative
of \(q_M\phi(c/D_\lambda)\) is
\[
 \frac98\frac{c(q_{\min}-\lambda)}{D_\lambda^2}>0.
\]
Here \(9/8\) is the value of \(\phi(\rho)/\rho\) on this branch.
On the trapezoid branch the derivative has the sign of
\[
 1+3(1-\rho_2)^2-3(q_M/D_\lambda)\rho_2(1-\rho_2)
 >1-3\rho_2(1-\rho_2)\ge1/4,
\]
since \(q_M<D_\lambda\) for all sufficiently large \(Y\). The maximum
is at \(q_M=q_{\max}+2\lambda\). There \(D_\lambda=1\) and
\(\rho_2^*=1-b_0\). The same endpoint applies to the first chart above.
Thus both chart comparisons give the finite bound
\[
 q_M\phi(\rho_j^*)\le(q_{\max}+2\lambda)\phi(1-b_0)
 \le\bar m+2\lambda.
\AutoTag{eq:old:7_8e}
\]
Here \(j=1,2\) indexes the two local counting orientations and their
starred comparison ratios. Either orientation is available for every
source, whatever its selected-chart tag. The collar below compares the
local ratios with them.
Their common limiting cap as \(Y\to\infty\) is
\(q_{\max}/(1+3b_0^2)=\bar m\). They are alternative counting routes
for the same full \(5W\) source. Each source box is counted once, by its
better route. The choice is made separately for each box.

It remains to bound the finite collars. For
\(\rho\in(0,1)\), \(\phi'\le9/8\). On the triangle branch
\(\phi'=9/8\). On the trapezoid branch, with
\(r=1-\rho\in[0,1/3]\),
\(\phi'=6r/(1+3r^2)^2\le9/8\). Thus the partial derivatives of
\(q_M\phi(c_0/\tau)\) in \((q_M,c_0,\tau)\) are bounded by
\((1,9/4,9/4)\) whenever \(q_M<1\) and \(\tau>1/2\). In the two
charts, \eqref{eq:old:7_8a} and the local scale definitions give
\[
 |c_0-c^*|\le B_Y,\quad |\tau-\tau^*|\le B_Y,
 \qquad B_Y=6d_2+\tfrac12\ell_{\rm end}+\lambda.
\AutoTag{eq:old:7_8f}
\]
For chart one the exact residuals are \eqref{eq:old:7_8c}. For chart two
they are
\[
\begin{aligned}
 \log_Y(BC)-(1-z_u)&=\log_Y\frac{m+n+1}{Y},\\
 \tau^*-\log_Y(B^2f/t)
 &=\log_Y\left(\frac{Y}{m+n+1}
             \sqrt{\frac{m+1}{m}}\sqrt{\frac{n}{n+1}}\right).
\end{aligned}
\AutoTag{eq:old:7_8g}
\]
Here \(Y/2<m+n+1\le2Y\). So the logarithm of
\((m+n+1)/Y\), and also its negative, has magnitude at most \(d_2\).
In \eqref{eq:old:7_8c}, the middle factor \(n/(m+n+1)\) lies between
\(1/(2+2/Y)\) and \(1\). Its logarithm has magnitude at most
\(d_2+\lambda\). The other factors contribute at most
\((h_m+h_n)/2\). The local dyadic replacement of \(AD\) by \(D_0E_0\)
costs at most \(2d_2\). The replacement of \(A^2f/u\) by \(X/H_0\)
costs at most \(4d_2\). In the second chart, the first residual in
\eqref{eq:old:7_8g} costs at most \(d_2\), and the second at most
\(d_2+(h_m+h_n)/2\). Its local scale changes have the same respective
costs \(2d_2\) and \(4d_2\). Also \(h_m+h_n\le\ell_{\rm end}\), where
\(h_m=\log_Y(1+1/m)\), \(h_n=\log_Y(1+1/n)\). The bounds
\(m>Y^{\tau_0}/\sqrt3\) and \(n>Y/2\) imply
\[
 h_m+h_n\le
 \log_Y\bigl((1+\sqrt3Y^{-\tau_0})(1+2/Y)\bigr)
 =\ell_{\rm end}.
\]
The larger of these bounds gives \eqref{eq:old:7_8f}, with room for the
extra \(\lambda\) term in the first chart. This keeps one
\(d_2\) unit of slack beyond the largest bound above.

Put
\[
 E_Y=\rho_Y+\log_Y(96\cdot18\cdot324\cdot8\cdot3\cdot2),
 \qquad \rho_Y=11d_2+2\ell_{\rm end}+2\lambda.
\AutoTag{eq:old:7_8h}
\]
The six factors pay for fixed comparisons. The factor 96 is used
in \(X/H_0>Y^{\tau_0}/96\). The factor 18 is used in \(M<18U_0V_*\).
The factor 324 is used in
the first anchor \(\Lambda_1=\log(324U_0V_*)\), together with
\(U_0V_*\le4M\). The factor 8 is used in the second anchor
\(\Lambda_2=\log(8D_0E_0)\). The factors 3 and 2 are used in \(X_1=N/3\)
and in \(X_2=X/(2H_0)\). The source collars contribute at most
\(2\lambda\) in \(q_M\). Then \eqref{eq:old:7_8f} contributes at most
\((1+9/4+9/4)B_Y+2\lambda\le6\rho_Y\). With the fixed logarithmic
comparisons above, these terms are bounded by \(6E_Y\).
The upper comparison \(U_0V_*\le4M\) contributes \(\log_Y4=2d_2\).
This lies in the \(11d_2\) part of \(\rho_Y\), hence in the
\(6\rho_Y\) collar. Also
\[
 \frac{\Lambda_1}{\log Y}=q_M+\log_Y(324U_0V_*/M),
\]
so the coefficient of \(q_M=\log_YM\) in this anchor is exactly one.
Hence the full finite cap is at most \(\bar m+6E_Y\). Since \(E_Y\to0\),
we may impose \(6E_Y<\xi/2\). This leaves a grid margin of at least
\(\xi/2\) in \eqref{eq:old:7_6}.

On the trapezoid branch the strict determinant surplus
at the limiting source comparison is at least
\[
\begin{aligned}
 \frac{q_{\min}\xi}{\bar m(\bar m+\xi)}
 &=\frac{997846754805373553496571768583986273}
         {6225284767495687836225589800000000000000}\\
 &=0.0001602893\ldots>0.
\end{aligned}
\AutoTag{eq:old:7_9}
\]
The triangle branch has its own strict surplus. The fixed margin \(\xi\)
is charged in \eqref{eq:old:7_6}, before the degree family is chosen. The
closed source-parameter ranges are compact, and \eqref{eq:old:7_9} gives a
positive surplus. Together they give a fixed \(\delta_T>0\) below the
minimum trapezoid-branch gap. The inequalities
\[
 q_{\min}-a_{\rm grid}>0,\quad \tau_0-a_{\rm grid}>0,\quad
 \tau_0-q_{\max}=\frac{9351751}{30000000}>0,\quad
 b_0-\frac{3e}{2}=\frac{972179}{7500000}>0
\AutoTag{eq:old:7_10}
\]
put all ratios in compact subintervals of their required ranges. Thus one
finite set of degrees, one Taylor order, and one threshold work for every
cell and shift.

\begin{lemma}[Ratio-cell determinant]\label{lem:ratio-cell-determinant}
Fix the parameters in \eqref{eq:old:7_2}--\eqref{eq:old:7_6}, with
\(3/5<a_{\rm grid}<\min(q_{\min},\tau_0)\),
\(q_{\max}<\tau_0\), and \(b_0>3e/2\). Let \(Y\) be sufficiently large.
Keep the exact first-match \(5W\) population
\eqref{eq:old:3_1}, with its source masks, original selected-chart labels,
and all preceding-row complements. Take either counting orientation,
any local box satisfying
\eqref{eq:old:7_4}--\eqref{eq:old:7_5} and \(X/H_0\le Y^4\), and any
cell in \eqref{eq:old:7_7}. Then there is a
nonzero polynomial
\(F_i\in\mathbb Z[X,Z]\) of degree at
most \(D_\theta\) and coefficient height at most \(Y^{H_\theta}\). Here
\(D_\theta,H_\theta\) depend only on the fixed parameter vector. It
satisfies \(F_i(J/K,g/f)=0\) for every retained source point in that cell
and every admissible shift in the box. The bounds are uniform in the
counting orientation, original selected chart, box, cell, labels and shift.
\end{lemma}
\begin{proof}
The proof is an integer calculation. We give the form used below. Put
\[
 x={f\over g},\quad y={K\over J},\quad y_0=r_i^{-1},\quad
 \delta_{\rm cell}={s\over fJ^2},\quad
 \alpha_{\rm cell}={x\over y_0^2}-1.
\]
The conic gives \(x=y^2/(1+\delta_{\rm cell})\) and
\(y/y_0=\sqrt{1+\alpha_{\rm cell}}\sqrt{1+\delta_{\rm cell}}\).
As \(J/K\in(r_i,r_{i+1}]\), we have
\(i/(i+1)\le y/y_0<1\), with \(i\ge N\). Also
\[
0<\delta_{\rm cell}\le \frac{H_0}{U_0D_0^2}
=\frac{H_0}{X}\le\frac1N
\]
for sufficiently large \(Y\), since \(X/H_0>Y^{\tau_0}/96\),
\(N\le2Y^{a_{\rm grid}}\), and \(a_{\rm grid}<\tau_0\). Therefore
\[
|\alpha_{\rm cell}|
=1-\frac{(y/y_0)^2}{1+\delta_{\rm cell}}
\le \frac2N+\frac{H_0}{X}
\le\frac3N.
\]
Next,
\[
0\le\frac{y}{y_0}-\sqrt{1+\alpha_{\rm cell}}
=\sqrt{1+\alpha_{\rm cell}}\bigl(\sqrt{1+\delta_{\rm cell}}-1\bigr)
\le\frac{H_0}{2X},
\]
using \(\sqrt{1+\alpha_{\rm cell}}\le1\) and
\(\sqrt{1+\delta}-1\le\delta/2\). Thus
\[
 |\alpha_{\rm cell}|\le3/N,\qquad
 0\le y/y_0-\sqrt{1+\alpha_{\rm cell}}\le H_0/(2X).
 \AutoTag{eq:old:7_10a}
\]
As \(X_2=X/(2H_0)\), the conic error is
\(H_0/(2X)=X_2^{-1}/4\).
Set \(T_{\rm tay}=7\). Let \(P_{T_{\rm tay}}\) be the degree-seven
Taylor polynomial of \(\sqrt{1+\alpha_{\rm cell}}-1\). For
\(c_k=|\binom{1/2}{k}|\), we have
\(c_{k+1}/c_k=(k-1/2)/(k+1)<1\) for \(k\ge1\), and
\(c_8=429/32768\). So, for
\(|\alpha_{\rm cell}|\le1/2\), the remainder is at most
\[
 \frac{c_8|\alpha_{\rm cell}|^8}{1-|\alpha_{\rm cell}|}
 \le\frac{429}{16384}|\alpha_{\rm cell}|^8
 <2|\alpha_{\rm cell}|^8.
\]
The source gives
\(X/H_0\le Y^4\) and \(a_{\rm grid}>3/5\). Hence, once
\(4\cdot3^8Y^{-8a_{\rm grid}}\le Y^{-4}\), the Taylor remainder is at
most \(X_2^{-1}/2\), where \(X_1=N/3\), \(X_2=X/(2H_0)\). To check this,
\(X/H_0\le Y^4\) gives \(X_2^{-1}\ge2Y^{-4}\). So the
Taylor bound is at most \(Y^{-4}/2\le X_2^{-1}/2\). Add the
conic error \(X_2^{-1}/4\). Then
\(\beta_{\rm cell}=y/y_0-1-P_{T_{\rm tay}}(\alpha_{\rm cell})\) satisfies
\(|\alpha_{\rm cell}|\le X_1^{-1}\),
\(|\beta_{\rm cell}|\le X_2^{-1}\), uniformly.

Take columns \(x^r y^j\), \(0\le r\le A_{\rm det}\),
\(0\le j\le B_{\rm det}\). Multiply the row for \((f,g,J,K)\) by
\(g^{A_{\rm det}}J^{B_{\rm det}}\). The entries are the integers
\[
 f^r g^{A_{\rm det}-r}K^jJ^{B_{\rm det}-j}.             \AutoTag{eq:old:7_10b}
\]
After the anchor is removed the entries are
\((1+\alpha_{\rm cell})^r
(1+P_{T_{\rm tay}}(\alpha_{\rm cell})+\beta_{\rm cell})^j\). With
\(D_c=(A_{\rm det}+1)(B_{\rm det}+1)\), \(g<9V_*\),
\(y_0\le2E_0/D_0\), and \(J<2D_0\),
the total row and column arithmetic cost is at most
\[
 {D_c\over2}\{A_{\rm det}\Lambda_1+B_{\rm det}\Lambda_2\},\qquad
 \Lambda_1=\log(324U_0V_*),\quad
 \Lambda_2=\log(8D_0E_0).                                \AutoTag{eq:old:7_10c}
\]
This includes the factor
\(y_0^{D_c(A_{\rm det}+B_{\rm det}/2)}\).
To verify \(g<9V_*\), multiply the conic
\(g/f=(J/K)^2+s/(fK^2)\) by \(f\):
\[
g=f(J/K)^2+s/K^2
 <8U_0D_0^2/E_0^2+H_0/E_0^2
 =8V_*+H_0/E_0^2<9V_*.
\]
Here \(f<2U_0\), \(J/K<2D_0/E_0\), \(s<H_0\), \(K\ge E_0\),
and \(X/H_0>Y^{\tau_0}/96>1\) for sufficiently large \(Y\).

For the finite Cauchy--Binet calculation let
\(\ell_1=\log X_1\), \(\ell_2=\log X_2\),
\(p_1=\Lambda_1/\ell_1\), \(p_2=\Lambda_2/\ell_2\), and
\(\ell_{\rm det}=\ell_1/\ell_2\). Set \(\ell=\ell_{\rm det}\).
Here \(X_1=N/3\), \(X_2=X/(2H_0)\),
\(\ell_1=\log(N/3)\), and \(\ell_2=\log(X/(2H_0))\). So the finite
parameters in the strip comparison are exactly
\[
 p_1=\frac{\log(324U_0V_*)}{\log(N/3)},\qquad
 p_2=\frac{\log(8D_0E_0)}{\log(X/(2H_0))},\qquad
 \ell_{\rm det}=\frac{\log(N/3)}{\log(X/(2H_0))}.
\]
Define the strip weight
\[
 W_{\rm strip}=\min_{\substack{\mathcal S\subset\mathbb Z_{\ge0}\times\{0,\ldots,B_{\rm det}\}\\
                                  |\mathcal S|=D_c}}
                  \sum_{(i,k)\in\mathcal S}(i\ell_1+k\ell_2).
\]
After column normalization the entries are polynomials in
\((\alpha_{\rm cell},\beta_{\rm cell})\) of \(\beta_{\rm cell}\)-degree
at most \(B_{\rm det}\). The Taylor coefficients of
\(\sqrt{1+\alpha_{\rm cell}}-1\) have absolute sum at most one. So each
column's coefficient \(\ell^1\)-norm is at most \(2^{A_{\rm det}}3^{B_{\rm det}}\).
Cauchy--Binet, Hadamard, and the at most
\(2^{(A_{\rm det}+1+B_{\rm det}T_{\rm tay})(B_{\rm det}+1)}\) subsets
of the finite monomial set
then give
\[
 \begin{aligned}
 \log|\Delta_1|&\le \log K_{A_{\rm det},B_{\rm det},T_{\rm tay}}-W_{\rm strip},\\
 K_{A_{\rm det},B_{\rm det},T_{\rm tay}}
 &=2^{(A_{\rm det}+1+B_{\rm det}T_{\rm tay})(B_{\rm det}+1)}(D_c!)^2
 (2^{A_{\rm det}}3^{B_{\rm det}})^{D_c}.
 \end{aligned}
\AutoTag{eq:old:7_10d}
\]
Here \(K_{A_{\rm det},B_{\rm det},T_{\rm tay}}\) is fixed once the degree
and Taylor order are fixed.
It is independent of \(Y\), the box, the cell, and the shift.

We give the strip lower bound explicitly. Put
\(A'=A_{\rm det}+1\), \(B'=B_{\rm det}+1\), so \(D_c=A'B'\).
Suppose
\(A'\ell_1\ge B'\ell_2/2\). To each lattice point
\((i,k)\) attach its unit square \([i,i+1)\times[k,k+1)\). For
\(\Phi(x,y)=x\ell_1+y\ell_2\), the integral over that square is
\(i\ell_1+k\ell_2+(\ell_1+\ell_2)/2\). The union \(E\) of the
\(D_c\) squares has area \(D_c\) and lies in the strip
\(x\ge0,0\le y\le B'\). Among all sets of this area in that strip, the
integral of \(\Phi\) is smallest on the sublevel set
\(E_* = \{\Phi\le w\}\cap\{x\ge0,\ 0\le y\le B'\}\). Points outside \(E_*\)
have \(\Phi\ge w\), and points removed from \(E_*\) have \(\Phi\le w\).
So exchanging equal areas cannot lower the integral. The condition above
makes \(E_*\) a trapezoid with right edge
\(x=(w-y\ell_2)/\ell_1\). Its area is
\[
 \frac{wB'-\ell_2B'^2/2}{\ell_1}=A'B',\qquad
 w=A'\ell_1+B'\ell_2/2.
\]
Integrating first in \(x\), then in \(y\), gives
\[
 \int_{E_*}\Phi
 =\int_0^{B'}\frac{w^2-y^2\ell_2^2}{2\ell_1}\,dy
 =\frac12A'^2B'\ell_1+\frac12A'B'^2\ell_2
   -\frac{B'^3\ell_2^2}{24\ell_1}.
\AutoTag{eq:old:7_10e}
\]
Subtracting the unit-square correction gives
\[
 W_{\rm strip}\ge \frac12A'^2B'\ell_1+
 \frac12A'B'^2\ell_2-\frac{B'^3\ell_2^2}{24\ell_1}
 -\frac12A'B'(\ell_1+\ell_2).
\AutoTag{eq:old:7_10f}
\]

The integer row clearing and the anchor factors have total logarithmic
cost at most
\(C_{\rm arith}=D_c(A_{\rm det}\Lambda_1+B_{\rm det}\Lambda_2)/2\), as in
\eqref{eq:old:7_10c}. As \(A_{\rm det}\le A'\) and \(B_{\rm det}\le B'\),
this is at most
\(D_c(A'\Lambda_1+B'\Lambda_2)/2\). Put \(z=A'/B'\). Division by
\(B'^3\ell_2/2\) gives, on the trapezoid branch,
\[
 W_{\rm strip}-C_{\rm arith}\ge\frac{B'^3\ell_2}{2}
 \left[-\ell_{\rm det}(p_1-1)z^2+(1-p_2)z-\frac1{12\ell_{\rm det}}
       -\frac{z(\ell_{\rm det}+1)}{B'}\right].          \AutoTag{eq:old:7_10g}
\]
The first three terms form the concave quadratic
\[
 G(z)=-\ell_{\rm det}(p_1-1)z^2+(1-p_2)z
       -\frac1{12\ell_{\rm det}}.
\]
If \(p_2\ge2/3\), \(p_1>1\), and
\(p_1-1\le3(1-p_2)^2-\delta_T\), then
\[
 z_* =\frac{1-p_2}{2\ell_{\rm det}(p_1-1)}
 \ge\frac1{2\ell_{\rm det}},\qquad
 G(z_*)=\frac{3(1-p_2)^2-(p_1-1)}
 {12\ell_{\rm det}(p_1-1)}>0.
\]
Here we use \(3(1-p_2)^2\le1-p_2\) for \(p_2\ge2/3\). On the compact
source range, \(z_*\), \(\ell\), and \(p_1-1\) are uniformly bounded,
and \(\ell(p_1-1)\) is bounded away from zero. Choose one fixed \(B'\)
so large that the rounding loss
\(\ell(p_1-1)/B'^2\) and the boundary term
\((z_*+1/B')(\ell+1)/B'\) sum to less than half the minimum of
\(G(z_*)\). Then set \(A'=\lceil z_*B'\rceil\). Since
\(G(z)=G(z_*)-\ell(p_1-1)(z-z_*)^2\), the bracket in \eqref{eq:old:7_10g} is
positive uniformly. The strict source margin in \eqref{eq:old:7_9} gives the
positive minimum. Choose the fixed lower threshold so that the total
\(O(1/\log Y)\) collar is at most half that minimum. The finite comparison
then keeps at least half the limiting gap uniformly.

For \(p_2\le2/3\), use the whole quadrant in place of the strip.
The sublevel set \(\Phi\le w\) in \(x,y\ge0\) is a triangle of area
\[
 \frac{w^2}{2\ell_1\ell_2}=D_c,\qquad
 \int_{E_*}\Phi=\frac{2\sqrt2}{3}D_c^{3/2}\sqrt{\ell_1\ell_2}.
\]
The unit-square error is \(D_c(\ell_1+\ell_2)/2\). Put \(z=A'/B'\). The
leading weight minus the upper bound for \(C_{\rm arith}\), divided by
\(B'^3\ell_2/2\), is
\[
 z\left(\frac{4\sqrt2}{3}\sqrt{\ell z}-p_1\ell z-p_2\right).
\]
Set \(v=\sqrt{\ell z}\). The parenthesis is the concave quadratic
\((4\sqrt2/3)v-p_1v^2-p_2\). Its maximum is at
\(v=2\sqrt2/(3p_1)\), which corresponds to
\[
 z_* =\frac{8}{9p_1^2\ell},\qquad
 A'=\left\lceil\frac{8B'}{9p_1^2\ell}\right\rceil,
\]
and the maximum value is
\((8/9-p_1p_2)/p_1\). Thus the triangle branch has a uniform positive
gap whenever \(p_1p_2<8/9\). A fixed large enough \(B'\) absorbs
the square and rounding errors. At \(p_2=2/3\), the two threshold
conditions agree. In both branches the source inequality
\(N>Y^{q\phi(\rho)}\), with the fixed margin \eqref{eq:old:7_9}, gives
these strict conditions for every cell.

After the fixed collar cutoff, the normalized bracket has a uniform
positive lower bound on both branches. Also
\(\ell_2=\log(X/(2H_0))>\tau_0\log Y-\log192\). Hence
\(W_{\rm strip}-C_{\rm arith}\ge c_\theta\log Y\) for some fixed
\(c_\theta>0\) and all sufficiently large \(Y\). The degrees and
Taylor order range over a fixed finite family. So
\[
 K_{\max}=\max K_{A_{\rm det},B_{\rm det},7}
\]
is fixed. Increase the determinant cutoff until
\(c_\theta\log Y>\log K_{\max}\). Then
\(W_{\rm strip}-C_{\rm arith}>\log K_{\max}\) uniformly. Every
cleared \(D_c\)-minor is an integer of absolute value less than one,
hence zero. The matrix has rank \(r<D_c\). If \(r=0\),
a coordinate vector gives a kernel polynomial. Otherwise choose \(r\)
pivot columns and \(r\) rows giving a nonsingular minor \(A\), then one
free column \(j\). Set the \(j\)-coordinate to \(\det A\), the other
free coordinates to zero, and the pivot coordinates to
\(-\operatorname{adj}(A)M_{I,j}\). The resulting nonzero integer vector
annihilates the selected rows. They span the row space, so it annihilates
every row. Its coordinates are bounded-order minors, so they have
polynomial height. This gives a nonzero integer polynomial \(C(x,y)\)
vanishing at
\((x,y)=(f/g,K/J)\). If a cell has fewer than \(D_c\) rows, then
\(r<D_c\) is automatic and its points are counted directly; this includes
an empty cell. Define
\[
 F_i(R,S)=R^{B_{\rm det}}S^{A_{\rm det}}C(1/S,1/R).
\AutoTag{eq:old:7_10h}
\]
This is a nonzero integer polynomial of degrees at most
\((B_{\rm det},A_{\rm det})\), with
the same coefficient height. It vanishes at
\((R,S)=(J/K,g/f)\). The substitution permutes the coefficients, so the
integral polynomial and its height are unchanged. This proves the
ratio-cell lemma.
\end{proof}

\subsection{One reduced lattice and the complete cell census}

The auxiliary polynomial is bihomogenized in the two primitive pairs.
For its high-degree factors use the single canonical coefficient lattice
\[
 j(i)=\left\lfloor\frac{i^2}{L_{\rm grid}}\right\rfloor,\quad
 t_i=\frac{j(i)}{L_{\rm grid}}\frac{V_*}{U_0},\quad
 W_i(g,f)=\left(\frac{L_{\rm grid}}{V_*}(g-t_if),\frac f{U_0}\right).
\AutoTag{eq:old:7_11}
\]
Its determinant is \(L_{\rm grid}/(U_0V_*)\). The map \(i\mapsto j(i)\)
is injective, because
\(((i+1)^2-i^2)/L_{\rm grid}=(2i+1)/(2N)>1\).
Recall \(R_{\rm box}=D_0/E_0\). So \(V_*/U_0=R_{\rm box}^2\), and
\[
 W_i=\begin{pmatrix}L_{\rm grid}/V_*&-j(i)/U_0\\0&1/U_0\end{pmatrix}.
\]
The source conic gives
\(g/f=(J/K)^2+s/(fK^2)\). Since
\(r_i=iR_{\rm box}/L_{\rm grid}\), \(J/K\le r_i(1+1/i)\), and
\(0\le r_i^2-t_i<R_{\rm box}^2/L_{\rm grid}\), for \(N\ge4\) we get
\[
 0<g/f-t_i\le R_{\rm box}^2\left({10\over N}+{H_0\over X}\right).
\]
As \(f<2U_0\), \(L_{\rm grid}=2N\), and
\(NH_0/X\le2^{-20}\), the first coordinate of \(W_i(g,f)\) lies
between zero and \(41\). Its second lies in \([1,2]\). So all
source images lie in one fixed bounded rectangle. Put
\(\Lambda_i=W_i\mathbb Z^2\). Choose a Lagrange-reduced basis
\(v_{1,i},v_{2,i}\) of this image lattice, with lengths
\(\lambda_{1,i}\le\lambda_{2,i}\) and angle bounded away from zero.
Write \(v_{j,i}=W_i h_{j,i}\). Since \(v_{1,i},v_{2,i}\) are a basis
of \(\Lambda_i\), their integer preimages \(h_{1,i},h_{2,i}\) form a
unimodular basis of \(\mathbb Z^2\). The vector \(v_{1,i}\) is shortest
in the Euclidean norm of the image lattice \(\Lambda_i\). Define
\[
 T_{1,i}={32\sqrt2\over\lambda_{1,i}},\qquad
 T_{2,i}={32\sqrt2\over\lambda_{2,i}}.
\]
Then \(T_{2,i}\le T_{1,i}\). The determinant identity gives
\(T_{1,i}T_{2,i}\asymp U_0V_*/L_{\rm grid}\). A unimodular change
preserves primitivity. By the bounded rectangle and the reduced-basis
angle, every source point has coordinates
\(O(T_{1,i})\) in this basis. The map and its inverse are rational matrices of
polynomial height. Explicitly,
\[
 W_i^{-1}(y_1,y_2)=
 \left({V_*\over L_{\rm grid}}(y_1+j(i)y_2),U_0y_2\right).
\]
A common denominator for \(W_i\) is at most a fixed power of \(Y\). So
every nonzero vector of \(\Lambda_i\) has length at least \(Y^{-C}\).
Its covolume is \(L_{\rm grid}/(U_0V_*)\le Y^C\). Minkowski's bound
and the reduced-angle determinant identity then give
\(\lambda_{1,i},\lambda_{2,i}\le Y^C\). The inverse matrix gives
\(\lVert h_{j,i}\rVert\le Y^{C'}\). Thus the unimodular basis and every
fixed-degree transformed factor have polynomial height, uniformly in the
box, cell and labels. The first-coordinate radius \(T_{1,i}\) satisfies
\[
 \sum_i(T_{1,i}+1)\ll_\eta
 Y^\eta\{N+\sqrt{U_0V_*N}+U_0+V_*\}.
\AutoTag{eq:old:7_12}
\]
Take a dyadic range \(D_{\rm lat}<T_{1,i}\le
2D_{\rm lat}\). Write the primitive shortest preimage in source
coordinates as \((\omega_1,\omega_2)\). It satisfies
\[
 |\omega_1|\ll V_*/D_{\rm lat},\quad |\omega_2|\ll U_0/D_{\rm lat},\quad
 \omega_2j(i)=\frac{L_{\rm grid}U_0}{V_*}\omega_1+O(U_0/D_{\rm lat}).
\AutoTag{eq:old:7_13}
\]
The first coordinate bound follows from the first component of \(W_i\),
the bound \(j(i)/L_{\rm grid}=O(1)\), and
\(|\omega_2|\ll U_0/D_{\rm lat}\). Suppose both coordinates are nonzero.
Fix \(\omega_1\) first. There are \(O(V_*/D_{\rm lat})\) choices. For each such
\(\omega_1\), \eqref{eq:old:7_13} places the integer product
\(w_i=\omega_2j(i)\) in an interval
of length \(O(U_0/D_{\rm lat})\). As \(\omega_2\ne0\),
\(U_0/D_{\rm lat}\gg1\). So this interval contains
\(O(U_0/D_{\rm lat})\) integers. Also \(w_i\ne0\), because
\(j(i)>0\).
For each fixed \(w_i\), the divisor bound gives \(Y^\eta\) choices of
\((\omega_2,j(i))\). Injectivity of \(j\) gives at most one index \(i\) for
each \(j(i)\). So there are
\(O_\eta(Y^\eta U_0V_*/D_{\rm lat}^2)\) indices in this dyadic range.
Weighting by \(D_{\rm lat}\) and summing from the lattice minimum gives
\(O_\eta(Y^\eta\sqrt{U_0V_*N})\). If \(\omega_2=0\),
primitivity gives
\(\omega_1=\pm1\), and the total is \(O(V_*+N)\). If \(\omega_1=0\), it is
\(O(U_0+N)\). The cells with \(T_{1,i}<1\) and all endpoint \(+1\)
terms are included in \(N\). This proves \eqref{eq:old:7_12}.

The nonzero product is at most a fixed power of \(Y\).
So its divisor estimate is uniform over the boxes and cells. The lower
limit \(D_{\rm lat}\gg\sqrt{U_0V_*/N}\) follows from
\(T_{1,i}T_{2,i}\asymp U_0V_*/N\) and \(T_{1,i}\ge T_{2,i}\).
Summing the unweighted number of non-axis indices over the same dyadic
ranges gives \(O_\eta(Y^\eta N)\). On the first axis, \(\omega_2=0\).
There primitivity gives \(\omega_1=\pm1\), and the contribution is
\(V_*+N\). On
the second axis, \(\omega_1=0\). There it gives \(\omega_2=\pm1\), and the
contribution is \(U_0+N\). The \(T_{1,i}<1\) bins contribute another
\(O(N)\). The bound \eqref{eq:old:7_12} keeps all these terms.
The bins with \(T_{2,i}<1\) are included in the \(+1\) terms.

A factor has degree \((a_c,b_q)\): corner degree first,
coefficient degree second. If \(a_c\ge2\) and \(b_q\ge1\), apply
Heath--Brown's published Lemma~2. Put the primitive coefficient-lattice
pair first and the primitive corner pair second. The bound is
\[
 O_{a_c,b_q,\eta}\bigl((T_{1,i}+1)^{2/a_c+\eta}
 \|F\|^\eta\bigr),\qquad 2/a_c\le1,
\AutoTag{eq:old:7_14}
\]
and the second pair is unrestricted. The statement in
Appendix~\ref{app:sources}, cited as
\cite[Lemma~2, p.~8]{HB2012}, applies with the reduced coefficient-lattice
pair first and the corner pair second. Reversing the degree order gives
second-pair degree \(a_c\). The coefficient pair is bounded by
\(X\asymp T_{1,i}+1\), while the corner pair is unrestricted. The source
gcd conditions make both pairs primitive. The factor height is
\(Y^{O(1)}\) before \(\eta\) is chosen. The stated bound therefore gives
\eqref{eq:old:7_14}. Polynomial factor heights and the finite degree family
let the loss be made arbitrary.

The remaining factors are as follows.
\begin{enumerate}
\item A rational factor that is not absolutely irreducible has rational
points only on intersections with distinct conjugate components.
Resultants in the four affine and infinity charts give \(O(1)\) points at
fixed degree.
\item If \(a_c=0\), a coefficient pair fixes the polynomial but not yet
the shift. First choose its ordered residual factors. These fix the
nonzero integer \(L=cd-ab\). Then enumerate \(s\mid L\). Apply the
uniform square-fibre theorem of Appendix~\ref{app:fibre} with
\((a,b,h,n,m)=(f,g,s,J,K)\), \(x=2Y\), and fixed height exponent
\(A_{\rm ht}=2\). This costs \(Y^\eta\), including the shift divisor count.
\item If \(b_q=0\), fix \((J,K)\) and compare two coefficient pairs:
\[
\begin{aligned}
 K^2(g_1f_2-g_2f_1)&=s_1f_2-s_2f_1,\\
 |g_1f_2-g_2f_1|
 &<\frac{3S_hU_0}{E_0^2}
 =\frac{3U_0V_*H_0}{2X}
 \le\frac{6MH_0}{X}\\
 &<576Y^{q_{\max}+2\lambda-\tau_0}=o(1).
\end{aligned}
\AutoTag{eq:old:7_15}
\]
The determinant is an integer, so it vanishes for all sufficiently large
\(Y\). Primitivity forces the coefficient pairs to agree. Then the conic
fixes the shift. The strict gap is
\(\tau_0-q_{\max}=9351751/30000000\). Since
\(s_if_j\in[S_hU_0,4S_hU_0)\), the difference of two such products
has absolute value less than \(3S_hU_0\). Then \eqref{eq:old:7_5} gives
the bound.
\item For \(b_q\ge1\) and \(a_c=1\), write the factor as
\(K\mathcal P(f,g)+J\mathcal Q(f,g)\), with coprime homogeneous
\(\mathcal P,\mathcal Q\). Their resultant is nonzero. By its Bezout
identities,
\(\gcd(\mathcal P(f,g),\mathcal Q(f,g))\) divides a fixed
polynomial-height integer. Write
\(d_0=\gcd(\mathcal P(f,g),\mathcal Q(f,g))\). Homogeneous Bezout
identities and \(\gcd(f,g)=1\) give
\(d_0\mid\operatorname{Res}(\mathcal P,\mathcal Q)\). Primitivity
then determines the positive corner pair
\((J,K)=\pm(\mathcal P(f,g)/d_0,-\mathcal Q(f,g)/d_0)\).

Let \(d_f=f_J\) be the full prime-power part of \(f\) supported on primes
dividing \(J\). Let \(d_g=g_K\) be the full prime-power part of \(g\)
supported on primes dividing \(K\). Put \(v=f/d_f\), \(w=g/d_g\).
As \(d_f,d_g\) contain the full powers of their
supported primes, \(d_f,d_g,v,w\) are pairwise coprime, and
\(d_fd_gvw=fg=M\).
The four endpoint patterns use the labels
\[
\begin{array}{c|c|c|c}
 \mathcal P(0,1)&\mathcal Q(1,0)&\ell_f&\ell_g\\ \hline
 \ne0&\ne0&\rad(d_f)&\rad(d_g)\\
 \ne0&=0&\rad(d_f)&w\\
 =0&\ne0&v&\rad(d_g)\\
 =0&=0&v&w .
\end{array}                                                   \AutoTag{eq:old:7_16}
\]
Suppose \(\mathcal P(0,1)\ne0\). Every prime dividing \(d_f\) divides \(J\)
and does not divide \(K\) or \(g\). The factor equation modulo that
prime gives \(\mathcal P(f,g)\equiv0\). Homogeneity then gives
\(\mathcal P(0,1)\equiv0\). Hence \(\rad(d_f)\mid\mathcal P(0,1)\).
If \(\mathcal P(0,1)=0\), then \(f\mid\mathcal P(f,g)\). As
\(\gcd(v,J)=1\), the factor equation gives \(v\mid\mathcal Q(f,g)\).
So \(v\mid d_0\mid\operatorname{Res}(\mathcal P,\mathcal Q)\).
The same argument with \(g,K,w,\mathcal Q(1,0)\) proves the two
\(g\)-side entries. These divisibilities include full prime powers in
\(v,w\). The zero endpoint cases are handled by the matching table
entries. So each label divides a fixed nonzero endpoint coefficient or
resultant of polynomial height. Hence there are \(\ll_\eta Y^\eta\)
possible pairs of endpoint labels for any fixed \(\eta>0\).

For each of the four patterns, the sum shown is the lower-corner
base-\(Y\) exponent of the matching full assigned-part product. In
chart 1 these sums are
\[
 P+Q,\quad P+(q-c_1-Q),\quad(c_1-P)+Q,\quad
 (c_1-P)+(q-c_1-Q),
\]
and in chart 2 they are
\[
 R_{\rm LP}+T,\quad R_{\rm LP}+(q-c_2-T),\quad
 (c_2-R_{\rm LP})+T,\quad
 (c_2-R_{\rm LP})+(q-c_2-T),
\]
in the same row order. Here \(P,Q,R_{\rm LP},T\) are the logarithmic sums
in \eqref{eq:old:7_18}. The forms are \(\mathcal P,\mathcal Q\). Put
\(k(x)=x/\rad(x)\). Pairwise coprimality gives
\(k(d_f)k(d_g)k(v)k(w)=M/\rad(M)=K_{\rm res}\). A label is either the
full assigned part or its radical. So in every row
\[
 \ell_f\ell_g\ge
 \frac{\text{corresponding assigned-part product}}{K_{\rm res}}.
\]
By \(K_{\rm res}\le Y^{\Delta_{\rm exp}}\), each matching actual
product is at least \(Y^{E_{\rm pat}}\), where \(E_{\rm pat}\) is its
lower-corner exponent. Thus
\(\ell_f\ell_g\ge Y^{E_{\rm pat}-\Delta_{\rm exp}}\). These are exactly
the four chart pattern terms with loss \(-\Delta_{\rm exp}\) in
\eqref{eq:old:7_19}.
Each of the eight assigned factors is below twice its dyadic lower
endpoint. So \(M\le2^8Y^q\). Also \(U_0V_*\le4M\). Hence
\[
 U_0V_*\le2^{10}Y^q,\qquad
 \frac{U_0V_*}{\ell_f\ell_g}
 \le2^{10}Y^{q-E_{\rm pat}+\Delta_{\rm exp}}.
\]
For the spacing estimate, the conic gives
\[
 {g\over f}=\left({J\over K}\right)^2+{s\over fK^2}.
\]
Over a cell of width \(O(D_0/(E_0N))\) in \(J/K\), the first term
varies by \(O(D_0^2/(E_0^2N))=O(V_*/(U_0N))\). The second term has
size \(O(H_0/(U_0E_0^2))\), smaller by the factor
\[
 {H_0/(U_0E_0^2)\over V_*/(U_0N)}={NH_0\over X}=o(1)
\]
by \eqref{eq:old:7_5}--\eqref{eq:old:7_6} and \(a_{\rm grid}<\tau_0\).
Hence the ratios \(g/f\)
in one cell lie in an interval of length \(O(V_*/(U_0N))\). Take two
distinct primitive coefficient pairs with the same labels. Then
\(\ell_f\ell_g\mid f_1g_2-f_2g_1\). This coefficient determinant is a
nonzero integer, so its absolute value is at least \(\ell_f\ell_g\). As
\(U_0\le f_1,f_2<2U_0\), their ratios differ by at least
\(\ell_f\ell_g/(4U_0^2)\). Therefore
\[
 \#\{\text{coefficient pairs in one cell}\}
 \ll 1+\frac{U_0V_*}{N \ell_f\ell_g}.
\AutoTag{eq:old:7_17}
\]
The graph equation gives at most one positive primitive corner
pair from each coefficient pair. Fix an endpoint pattern.
Summing over the \(3N\) cells and the divisor-bounded labels gives
\[
 O_\eta\!\left(Y^\eta\left(N+
 \frac{U_0V_*}{\ell_f\ell_g}\right)\right).
\]
By the bound above for this pattern, the count is
\[
 O_\eta\!\left(Y^\eta\left(N+
 Y^{q-E_{\rm pat}+\Delta_{\rm exp}}\right)\right).
\]
The \(+1\) term, unit labels, and endpoint equalities are kept. The
eight chart-pattern classes have fixed cardinality. So their sum adds no
further power of \(Y\).
\end{enumerate}
The source corner coordinates are positive. The lattice census in
\eqref{eq:old:7_12} keeps both shortest-vector axes, which contribute
\(U_0+V_*\). In either counting orientation take \(c=c_1\) or
\(c=c_2\) respectively, as defined below. The four assigned factors of
\(f\) and the four complementary factors of \(g\) give
\[
 U_0\le f<2^4Y^c,\qquad V_*\le4g<2^6Y^{q-c}.
\]
These two terms are charged through the first two savings in
\eqref{eq:old:7_19}. The four graph patterns give its other four
savings. The grid and square-root terms are kept separately.

\subsection{The lower-corner linear program and lattice-point count}

The exact finite optimization below controls the source's choice of chart
or corner route. In each source box, let
\(x_0,\ldots,x_7\) be the lower endpoints of the base-\(Y\) logarithms
of the eight uniquely assigned corner factors, in the order
\[
 (a_A,a_B,b_C,b_D,c_A,c_C,d_B,d_D).
\]
Each factor is the full prime-power part of its residual supported on the
eligible corner set. Supported complementary parts are fixed by the complete
coefficient-corner labels and are charged through \eqref{eq:old:6_6}.
Write
\[
\begin{aligned}
x_{AB}&=x_0+x_1,& x_{AC}&=x_4+x_5,& q&=\sum_{i=0}^7x_i,\\
x_{ABCD}&=x_0+x_1+x_2+x_3,& c_1&=x_{AB}+x_{AC},\\
c_2&=x_0+x_1+x_6+x_7,&
P&=x_0+x_4,& Q&=x_3+x_7,\\
R_{\rm LP}&=x_1+x_6,& T&=x_2+x_5.
\end{aligned}
\AutoTag{eq:old:7_18}
\]
The source constraints include the canonical radical, shell, unit, chart,
and corner-capacity inequalities. Their limiting affine logarithmic forms
are listed in Appendix~\ref{app:arithmetic}. Their finite source collars
are recorded below. The source log coordinates
\(\mu,z_u,z_t\) also use their box lower endpoints. Thus the eleven
source coordinates are binned, and the auxiliary \(W\) is evaluated by
\eqref{eq:old:7_21} at those endpoints.
With \(\rho_Y\) as in \eqref{eq:old:7_8h}, every source-row inequality has
lower-corner slack at most
\(\rho_Y\). The 22 affine source rows in Appendix~\ref{app:arithmetic}
have coefficient \(\ell^1\)-norm at most \(8\). So replacing the
eleven source coordinates by reference values costs at most \(8d_2\).
In deriving these rows from \eqref{eq:old:7_8a}, the remaining finite terms
are logarithms of
\(n/Y\), \((m+n+1)/Y\), \((m+1)/m\), \((n+1)/n\), and
\(\kappa_Y=1+1/Y\). The first two have absolute logarithm at most
\(d_2\). The next two together are at most \(\ell_{\rm end}\). Each
source row uses at most three further \(d_2\) units, two
\(\ell_{\rm end}\) units, and two powers of \(\kappa_Y\). This gives the
stated \(\rho_Y\). The count is applied after these source-row comparisons.

For chart 1 put \((c,F,G)=(c_1,P,Q)\), and for chart 2 put
\((c,F,G)=(c_2,R_{\rm LP},T)\). Its six nominal savings are
\[
\begin{aligned}
 &\frac23-c,\quad \frac23-q+c,\\
 &\frac23-q+F+G-\Delta_{\rm exp},\quad
 \frac23-q+F+(q-c-G)-\Delta_{\rm exp},\\
 &\frac23-q+(c-F)+G-\Delta_{\rm exp},\quad
 \frac23-q+(c-F)+(q-c-G)-\Delta_{\rm exp}.
\end{aligned}
\AutoTag{eq:old:7_19}
\]
Let \(A_j\) be their minimum. The four assigned-corner savings are
\[
\begin{aligned}
 G_A&=\frac23-\left(\frac{\mu+z_u-c_1}{2}+x_1+x_5+z_u\right),\\
 G_D&=\frac23-\left(\frac{\mu+z_u-q+c_1}{2}+x_2+x_6+z_u\right),\\
 G_B&=\frac23-\left(\frac12-\frac{z_u}{2}+\frac q4-\frac{c_2}{2}
                   +x_0+x_7+z_t\right),\\
 G_C&=\frac23-\left(\frac12-\frac{z_u}{2}-\frac q4+\frac{c_2}{2}
                   +x_3+x_4+z_t\right).
\end{aligned}
\AutoTag{eq:old:7_20}
\]
At the lower corner set
\[
 W=\max(0,A_1,A_2,G_A,G_D,G_B,G_C).
\AutoTag{eq:old:7_21}
\]
The exact linear program has the twelve nonnegative coordinates
\((x_0,\ldots,x_7,W,z_u,z_t,\mu)\), 22 source rows, two chart rows
and four corner-route rows. The 72 branches come from the selected-chart
choice and the two six-way choices of minimizing chart scores. Dropping the
inequalities that enforce those choices enlarges each branch polytope.
Exact primal and dual elimination on these rational polytopes gives
\[
 W\ge L_{5W}-5\rho_Y,\qquad
 L_{5W}=\frac{1203113}{90000000}.
\AutoTag{eq:old:7_22}
\]
In every branch the counting-row dual mass is \(1\), the source-row dual
mass is at most \(5\), and the total dual mass is at most \(6\). The finite
certificate and its exact row matrix are given in Appendix
~\ref{app:arithmetic}.

Next we tie these nominal rows to counts on the same filtered box.
In chart 1 write \(\alpha=ac\), \(\beta=bd\). For \(G_A\), record
\((A,a_B,c_C,u)\), with \(a_Ac_A=f_J\) paid through
\eqref{eq:old:6_6}. The identity
\(u+\alpha A^2=\beta D^2\) chooses \(D\) by a square-divisor choice and
then fixes \(\beta\). For \(G_D\), record
\((D,b_C,d_B,u)\), with \(b_Dd_D=g_K\) paid in the same way. Then
\(\beta D^2-u=\alpha A^2\) chooses \(A\)
and then fixes \(\alpha\). In both cases split \(\alpha=ac,\beta=bd\)
and reconstruct
\[
 B=\frac{cA+bD}{u},\quad C=\frac{aA+dD}{u},\quad
 t=\frac{cd-ab}{u},\quad m=ABa,\quad n=ACc.
\]
In chart 2 write \(\gamma=ad\), \(\delta=bc\). For \(G_B\), record
\((B,a_A,d_D,t)\), with \(a_Bd_B\) paid through \eqref{eq:old:6_6}.
Then \(t+\gamma B^2=\delta C^2\) chooses \(C\)
and then fixes \(\delta\). For \(G_C\), record \((C,b_D,c_A,t)\), with
\(b_Cc_C\) paid through \eqref{eq:old:6_6}. Then
\(\delta C^2-t=\gamma B^2\) chooses \(B\) and fixes \(\gamma\).
Split \(\gamma=ad,\delta=bc\) and reconstruct
\[
 A=\frac{dB-bC}{t},\quad D=\frac{cC-aB}{t},\quad
 u=\frac{cd-ab}{t},\quad m=ABa,\quad n=ACc.
\]
Each square-divisor choice and supported-prime choice costs \(Y^\eta\)
at the fixed polynomial heights. The residual products must also be split
into their ordered corner factors. For example, splitting
\(\alpha=ac\) and \(\beta=bd\) costs at most
\(\tauDiv(\alpha)\tauDiv(\beta)\ll_\eta Y^\eta\). The same bound holds
for \(\gamma=ad\), \(\delta=bc\). Once these choices are fixed, the
inverse formulas give the other corners, and \(ut=L\) gives
the other shift. Thus the total inverse multiplicity is \(Y^\eta\).

Each encoding freely enumerates one corner, two individual assigned parts,
and the shift of that route. These four data have upper dyadic endpoints
costing \(4d_2\). The supported parts are determined by the full
coefficient-corner labels. They are not extra free variables. Supported
parts use the full prime-power convention of \eqref{eq:old:2_4}. Their
supported-prime choices are paid by the uniform estimate \eqref{eq:old:6_6}.
Each sum of two assigned-part coordinates is a lower bound for the
base-\(Y\) logarithm of their product.
To compare the corner with its LP reference value, put
\[
\nu_n=\log_Y n,\quad H=\log_Y(m+n+1),\quad
h_m=\log_Y(1+1/m),\quad h_n=\log_Y(1+1/n).
\]
Let subscript \(r\) denote the actual source values of \(\mu,z_u,q,c_1,c_2\).
Define the lower-corner reference exponents
\[
\begin{aligned}
\mathsf A&=(\mu+z_u-c_1)/2,&
\mathsf D&=(\mu+z_u-q+c_1)/2,\\
\mathsf B&=1/2-z_u/2+q/4-c_2/2,&
\mathsf C&=1/2-z_u/2-q/4+c_2/2.
\end{aligned}
\]
Take logarithms in \eqref{eq:old:7_8a}, and use \(Z^2M=m(m+1)n(n+1)\) and
\(ZL=(n-m)(m+n+1)\). This gives the exact formulas
\[
\begin{aligned}
\log_Y A&=(\mu_r+z_{u,r}-c_{1,r})/2+(\nu_n-H)/2,\\
\log_Y D&=(\mu_r+z_{u,r}-q_r+c_{1,r})/2
             +(\nu_n-H+h_m+h_n)/2,\\
\log_Y B&=1/2-z_{u,r}/2+q_r/4-c_{2,r}/2
             +(H-1)/2+(h_n-h_m)/4,\\
\log_Y C&=1/2-z_{u,r}/2-q_r/4+c_{2,r}/2
             +(H-1)/2+(h_m-h_n)/4.
\end{aligned}
\AutoTag{eq:old:7_23a}
\]
Each of the eleven independently binned source coordinates lies between
its lower endpoint and that endpoint plus \(d_2\). Apply this
coordinatewise to the linear forms in \eqref{eq:old:7_23a}, before
combining the corner coordinates into \(q,c_1,c_2\). This gives
\[
\begin{aligned}
\log_YA&\le\mathsf A+d_2,&
\log_YD&\le\mathsf D+3d_2+\ell_{\rm end}/2,\\
\log_YB&\le\mathsf B+3d_2/2+\lambda/2,&
\log_YC&\le\mathsf C+3d_2/2+\ell_C/2+\lambda/2.
\end{aligned}
\AutoTag{eq:old:7_23b}
\]
The exact correction in the \(C\) formula is
\((h_m-h_n)/4\le\ell_C/4\). The \(\ell_C/2\) shown is
conservative.
Here \(0<H-\nu_n\le d_2+\lambda\), \(H-1\le d_2+\lambda\),
\(h_m+h_n\le\ell_{\rm end}\), and
\(0\le h_m-h_n\le\ell_C\). The last bound follows from
\(m>Y^{\tau_0}/\sqrt3\), \(n>Y/2\), and \(m<n\). The coefficient
variation from the lower endpoints is at most \(d_2\) for \(\mathsf A\)
and \(3d_2\) for \(\mathsf D\). The \(3d_2\) bound for \(\mathsf D\) is
conservative. Cancellation between \(q\) and \(c_1\) reduces
the positive-coordinate variation in its lower-corner form to at most
\(d_2\). We keep the larger uniform allowance. In each of
\(\mathsf B,\mathsf C\),
the positive coefficients among \(q/4-c_2/2\) (or its negative) sum to
one across the original eight corner coordinates. So the upper variation is
at most \(d_2\), and the other coordinates decrease the
form. The other terms in \eqref{eq:old:7_23a} give the corrections shown.
Adding the \(4d_2\) data endpoints gives the
four route costs
\[
5d_2,\qquad 7d_2+\ell_{\rm end}/2,\qquad
11d_2/2+\lambda/2,\qquad
11d_2/2+\ell_C/2+\lambda/2.
\AutoTag{eq:old:7_23c}
\]
By \eqref{eq:old:7_3}, \(\ell_C\le\ell_{\rm end}\).
The two chart counts use the complete \eqref{eq:old:7_12} and
\eqref{eq:old:7_17}, with both axis terms included in the six chart
costs. They apply to the same filtered source box. Refine it at once
into the dyadic boxes of both counting orientations, keeping its original
selected-chart label. The fixed power of \(\log Y\) refinements is
absorbed in \(Y^\eta\). Their lower-corner loss is at most \(10d_2\).
This includes both axis terms and the four-part products in all graph
patterns. For \(Y\ge2\), every route cost is at most
\(\rho_Y\), by its definition and \(\ell_C\le\ell_{\rm end}\).
These are alternative bounds for the same source box.
Choose a route attaining \eqref{eq:old:7_21}. Its lattice-point count has
geometric saving at least \(W-\rho_Y\), before the arbitrary analytic
loss allocated in Section~\ref{sec:epsilon}. Then \eqref{eq:old:7_22}
gives \(L_{5W}-6\rho_Y\). The geometric,
interval, and square-root terms enter the completed source count
additively. The axis costs are already in the chart bound. The
number of terms is fixed. So the largest term exponent, that is the
minimum saving, controls their sum.

The interval and square-root terms in \eqref{eq:old:7_12} have savings
\[
\begin{aligned}
 s_{\rm int}&=\frac23-a_{\rm grid}\\
 &=\frac{235120131355662893}{9520800917171070000}
   \approx0.0246954,\\
 s_{\rm HD}&=\frac23-\frac{q_{\max}+a_{\rm grid}}{2}\\
 &=\frac{29390272487700347267}{4760400458585535000000}
   \approx0.00617391.
\end{aligned}\AutoTag{eq:old:7_24}
\]
For
\[
 c_N=\log_Y(1+Y^{-a_{\rm grid}}),\qquad E_q=2\lambda+\log_Y18,
\]
the interval term has saving at least \(s_{\rm int}-c_N\). For the
square-root term, \(M<Y^{q_{\max}+2\lambda}\) and
\(N<Y^{a_{\rm grid}+c_N}\) give
\[
 \sqrt{4MN}
 \le Y^{(q_{\max}+a_{\rm grid})/2+\lambda+c_N/2+d_2}.
\]
This is a conservative uniform upper bound for the finite high-degree
correction. It uses the largest allowed source product and the ceiling
bound for \(N\). It uses no correlation between those endpoints.
So the high-degree saving is at least
\(s_{\rm HD}-(E_q+c_N)/2-12d_2\), since
\[
 \frac{E_q+c_N}{2}+12d_2
 =\lambda+\frac{\log_Y18}{2}+\frac{c_N}{2}+12d_2
 \ge\lambda+\frac{c_N}{2}+d_2.
\]
Thus all endpoint terms are kept. Let
\(\eta>0\). Choose all finite determinant, factor, divisor,
Rankin, and logarithmic losses to total at most \(\eta/2\). Then
enlarge one threshold. This gives
\[
\sigma_{5W}=\min\!\left(L_{5W}-6\rho_Y,\,
 s_{\rm int}-c_N,\,
 s_{\rm HD}-(E_q+c_N)/2-12d_2\right).
\]
Thus
\[
 \#\{5W\text{ sources in one box}\}
 \ll_\eta Y^{\,2/3-\sigma_{5W}+\eta/2}.
\AutoTag{eq:old:7_25}
\]
The number of boxes is a fixed power of \(\log Y\). It is absorbed in
the remaining \(\eta/2\). The margins of the three limiting caps above
\(S_0\) are
\[
\begin{aligned}
 L_{5W}-S_0&=\begin{gathered}\frac{647471}{90000000}\\[-2pt]\text{\scriptsize (0.007194)}\end{gathered},\\
 s_{\rm int}-S_0&=\begin{gathered}\frac{88170305326616070517}{4760400458585535000000}\\[-2pt]\text{\scriptsize (0.018522)}\end{gathered},\\
 s_{\rm HD}-S_0&=\begin{gathered}\frac{128034121242821}{1190100114646383750000}\\[-2pt]\text{\scriptsize (1.076$\times10^{-7}$)}\end{gathered}>0.
\end{aligned}
\AutoTag{eq:old:7_26}
\]
The three limiting caps are at least \(S_0\) by \eqref{eq:old:7_26}. Their corrected
versions bound the additive contributions in \eqref{eq:old:7_25}, whose minimum controls
the sum exponent. For \(Y\ge4\), since \(\tau_0>1/2\),
\[
 (1+\sqrt3Y^{-\tau_0})(1+2/Y)(1+1/Y)
 <2\cdot\frac32\cdot\frac54=\frac{15}{4}<4.
\]
Thus \(\ell_{\rm end}+\lambda\le2d_2\), hence
\(\rho_Y\le15d_2\) and \(6\rho_Y\le90\log2/\log Y\).
The inequality \(\log(1+t)\le t\) for \(t\ge0\) and
\(a_{\rm grid}>1/2\) give, for \(Y\ge2\),
\[
 c_N\le\frac{Y^{-a_{\rm grid}}}{\log Y}<\frac1{\log Y},
 \qquad
 \lambda+\frac{c_N}{2}
 \le\frac{Y^{-1}+\tfrac12Y^{-1/2}}{\log Y}
 <\frac1{\log Y}.
\]
Therefore, for \(Y\ge4\),
\[
 \frac{E_q+c_N}{2}+12d_2
 \le\frac{\tfrac12\log18+12\log2+1}{\log Y}
 <\frac{90\log2}{\log Y}.
\]
Thus every correction in \eqref{eq:old:7_25} is at most
\(90\log2/\log Y\), and \eqref{eq:old:7_26} gives
\[
 Y^{S_0-\sigma_{5W}}\le2^{90}.
\]
The finite corrections therefore change only the implied constant.
After the box sum, \eqref{eq:old:7_25} implies the uniform estimate
\[
 \#5W\ll_\eta Y^{2/3-S_0+\eta}.
\AutoTag{eq:old:7_27}
\]
\begin{rowproposition}[\(5W\) row estimate]\label{prop:row-5w}
For every \(\eta>0\), the complete first-match remainder satisfies
\[
 \#5W\ll_\eta Y^{2/3-S_0+\eta}.
\]
\end{rowproposition}
\begin{proof}
The three additive source-box contributions are bounded in
\eqref{eq:old:7_25}. Their limiting caps and positive margins are
\eqref{eq:old:7_26}. This gives \eqref{eq:old:7_27}.
\end{proof}
The threshold in \eqref{eq:old:7_27} depends on \(\eta\) and the fixed
source parameters. It does not depend on any moving cell, shift, label,
factor, or resultant value. For each box and cell the degree set and
height exponents are fixed. For each selected factor the polynomial, its
endpoint coefficients and its resultant are fixed. They may vary between
boxes, cells and factors.

The exact global reserve is
\[
 g_0=S_0-\frac1{165}=\frac{18677}{165000000}>0.
\AutoTag{eq:old:7_28}
\]
It is used in Section~\ref{sec:epsilon} to choose one common loss and
one \(Y_{\rm final}\). The source-to-count argument above gives the
uniform estimate used there.

\section{Exact exponent summary}\label{sec:ledger}

Put
\[
\begin{gathered}
S_0=\epsilon_0=k_0=\frac{30869}{5000000},\quad
l_0=2S_0,\quad j_0=p_0=3S_0,\quad \xi=\frac1{10000},\\
a_0=\frac6{125},\quad
b_0=\frac{4166537}{30000000},\quad
\Delta_{\rm exp}=3S_0.
\end{gathered}\AutoTag{eq:old:8_1}
\]
All entries below are limiting savings from \(2/3\) at this fixed
rational parameter vector. Sections~\ref{sec:long} and
\ref{sec:epsilon} handle the finite source collars and the arbitrary
analytic losses. The row estimates in Sections~\ref{sec:outer}--\ref{sec:long}
give the exact values in the second column. The symbolic counting exponents
and their evaluations are displayed in Sections~\ref{sec:joint}--\ref{sec:intermediate}
(for \(3u\), \(3t\), \(4H\), and \(4P\)). Appendix~\ref{app:arithmetic}
collects the parameter vector, exact savings, and margins above \(S_0\).

\setlength{\LTleft}{0pt}
\setlength{\LTright}{0pt}
\renewcommand{\arraystretch}{1.35}
\footnotesize
\unnumtable
The table expands the savings used in the preceding row estimates. The
minimum row and its margin are read from the exact fractions. The rounded
decimals are annotations only.
\begin{longtable}{@{}p{0.30\textwidth}p{0.53\textwidth}p{0.12\textwidth}@{}}
\toprule
Contribution & Limiting saving & Decimal\\
\midrule
\endfirsthead
\toprule
Contribution & Limiting saving & Decimal\\
\midrule
\endhead
\(O_R\) & \(30869/5000000\) & --\\
\(O_M\) & \(30869/5000000\) & --\\
\(O_{\operatorname{rad}M}\) & \(30869/5000000\) & --\\
\(Q0\) & \(2/3\) & --\\
\(Q1\) & \(30869/5000000\) & --\\
\(Q2t_{\rm int}\) & \(972179/7500000\) & --\\
\(Q2t_{\rm harm}\) & \(30869/5000000\) & --\\
\(Q2u_{\rm int}\) & \(2222179/15000000\) & --\\
\(Q2u_{\rm harm}\) & \(30869/1250000\) & --\\
\(3u_{\rm sqrt}\) & \(5102312049347/600000000000000\) & --\\
\(3u_{\rm int}\) & \(23206592049347/300000000000000\) & --\\
\(3u_{\rm core}\) & \(2473/12000\) & --\\
\(3u_{\rm small,1}\) & \(1619393/6000000\) & --\\
\(3u_{\rm small,2}\) & \(8003/500000\) & --\\
\(3t_{\rm sqrt}\) & \(388751/60000000\) & --\\
\(3t_{\rm int}\) & \(2199179/30000000\) & --\\
\(3t_{\rm axes}\) & \(2726537/15000000\) & --\\
\(4E\) & \(30869/5000000\) & --\\
\(4H\) & \(140441/5000000\) & --\\
\(4P_{\rm graph}\) & \(30869/5000000\) & --\\
\(4P_{\rm HD}\) & \(107393/15000000\) & --\\
\(4P_{\rm cover}\) & \(2/75\) & --\\
\(4P_{\rm axes}\) & \(4166537/30000000\) & --\\
\(5W_{\rm HD}\) & \(\dfrac{29390272487700347267}{4760400458585535000000}\) & \(0.006174\)\\[4pt]
\(5W_{\rm int}\) & \(\dfrac{235120131355662893}{9520800917171070000}\) & \(0.024695\)\\[4pt]
\(L_{5W}\) & \(1203113/90000000\) & \(0.013368\)\\
\bottomrule
\end{longtable}
\renewcommand{\arraystretch}{1}
\normalsize

Some entries add up within one row estimate. These are the five \(3u\)
entries, the three \(3t\) entries, the four \(4P\) entries, the two
\(Q2\) terms per chart and the three \(5W\) entries. In each row the
largest exponent controls the sum. Equivalently, the smallest saving does.
In each \(5W\) source box, Section~\ref{sec:long} selects the best chart
and assigned-corner route. The three displayed terms are the geometric,
interval and high-degree contributions that follow.
In the \(Q2\) entries, \(Q2t_{\rm int},Q2t_{\rm harm}\) are the first and
second terms of the \(Q2t\) estimate \eqref{eq:old:5_3}. Likewise,
\(Q2u_{\rm int},Q2u_{\rm harm}\) are the terms of the \(Q2u\) estimate.
Each pair is summed within its own first-match
population.
The table has 26 scalar terms. After grouping the three \(5W\)
contributions, 24 effective terms remain.

The minimum row saving is \(S_0\). It is attained by
\(O_R,O_M,O_{\operatorname{rad}M},Q1,Q2t_{\rm harm},4E\), and
\(4P_{\rm graph}\).
In the final row the smallest of the three savings is \(5W_{\rm HD}\), and
\[
5W_{\rm HD}-S_0=
\frac{128034121242821}{1190100114646383750000}
\quad\text{(about }1.076\times10^{-7}\text{)}>0.
\AutoTag{eq:old:8_2}
\]
The 72-branch minimum used for the geometric cap is
\[
L_{5W}=\frac{1203113}{90000000},\qquad
L_{5W}-S_0=\frac{647471}{90000000}>0.
\AutoTag{eq:old:8_3}
\]
Section~\ref{sec:long} proves the branch construction and the
comparison needed for \(5W\). No LP value is used as a count by itself.

\begin{proposition}[Reserve]\label{prop:reserve}
Put
\[
g_0=S_0-\frac1{165}=\frac{18677}{165000000},\qquad
\beta_0=\frac{g_0}{16}=\frac{18677}{2640000000}.
\]
Then
\[
S_0-\beta_0-\frac{g_0}{4}
=\frac1{165}+\frac{11g_0}{16}
=\frac{16205447}{2640000000}>\frac1{165}.
\AutoTag{eq:old:8_4}
\]
\end{proposition}

\begin{proof}
Put the displayed fractions over the common denominator
\(2640000000\). The loss \(\beta_0\) is charged once to each complete
analytic row estimate. In the last row, the \(g_0/4\) term covers the maximum of the
geometric, interval, and high-degree corrections.
\end{proof}

\section{One loss and one threshold}\label{sec:epsilon}

Fix the rational vector in \eqref{eq:old:8_1} and the complete analytic loss
\(\beta_0=g_0/16\) from \eqref{eq:old:8_4}. In every determinant application, the positive exponent gap is fixed
first. The degree and the matrix dimensions are chosen after it. This
gives a finite family of degrees and dimensions. A moving box may select
one member of it. The coefficients, the polynomial and the resultant may
vary with the box and with \(Y\). Only their degrees and height exponents
are uniform. Nonvanishing comes from the source determinant, from
coprimality, or from the endpoint alternative in the relevant row. The
height bound plays no part in it.

Fix a selected member. The cofactor solution and the explicit
integerizer give a polynomial height \(Y^{C_i}\). Here \(C_i\) depends
only on the fixed degree, the dimension and the allocated loss. Hence a
nonzero resultant is an integer of size at most \(Y^{C_i'}\)
for a fixed exponent \(C_i'\). The integer itself may change as the
source labels move.

We allocate the loss as follows. Suppose one complete row proof contains
\(N_i\) divisor, Heath--Brown, square-fibre, or lattice estimates. Let
the \(j\)-th estimate apply to a parameter bounded by \(Y^{h_{ij}}\).
Here \(N_i\) is the fixed finite number of estimate stages in the complete
row proof. A stage repeated over moving boxes keeps its one uniform
exponent choice. It does not incur a new loss for each box.
Choose \(h_{ij}\) at least as large as every polynomial-height exponent in
that invocation. This covers the auxiliary form, its integerization, any
nonzero resultant and the size exponent of the counted parameter.
For a product of height factors, add their exponents when choosing
\(h_{ij}\). Then it bounds the total contribution of the product.
The degrees and height exponents are uniform over the complete moving
source family in that row.
Choose the internal exponent \(\rho_{ij}\) so that
\[
0<\rho_{ij}\le
\frac{\beta_0}{8N_i(1+h_{ij})}.
\AutoTag{eq:old:9_1}
\]
Then
\[
\sum_{j=1}^{N_i}h_{ij}\rho_{ij}\le\frac{\beta_0}{8}.
\]
Allocate the other fixed portions of \(\beta_0\) to the finitely many
dyadic sums, logarithms and divisor multiplicities in the same complete
row proof. This is one allocation for the whole row. After a row
estimate is stated, no factor or inverse count is charged again.

The thirteen first-match estimates are stated in
Propositions~\ref{prop:outerrows} and~\ref{prop:q2}, and in the row
Propositions~\ref{prop:row-3u}, \ref{prop:row-3t},
\ref{prop:row-4e}, \ref{prop:row-4h}, \ref{prop:row-4p}, and
\ref{prop:row-5w}. They use the finite cutoff data summarized below. Each cutoff is fixed
by the corresponding proof and depends only on \(\beta_0\) and the
rational vector. The table records conditions and the estimates that
supply them; it does not give numerical values for all row thresholds
or analytic implied constants. Equation \eqref{eq:old:9_5} takes the
maximum of the thresholds used in those proofs, including all displayed
row conditions.

\begin{table}[!ht]
\centering
\caption{Threshold conditions for the thirteen row estimates.}
\label{tab:threshold-conditions}
\begin{tabular}{@{}>{\raggedright\arraybackslash}p{0.14\textwidth}>{\raggedright\arraybackslash}p{0.79\textwidth}@{}}
Row & Threshold conditions included in $Y_i(\beta_0)$\\ \hline
$O_R$ & $Y\ge64$, $H=Y/2\ge2$, $V=Y^{5/3-\epsilon_0}<H^2$; fixed-height partner and low-radical estimates\\
$O_M$ & Ordered residual-factor divisor bound and the fixed-height quadratic-fibre estimate\\
$O_{\operatorname{rad}M}$ & One fixed positive Rankin parameter and the same quadratic-fibre estimate\\
$Q0$ & $Y\ge256$ and $3Y^{q_{\max}-1}<1$\\
$Q1$ & Fixed-height bound for all $(P,L)$ labels\\
$Q2u$ & $1/3-5\epsilon_0>0$; unit-root count, one $u$-sum, and the dyadic sector count\\
$Q2t$ & $1/3-6\epsilon_0-l_0>0$; corresponding $t$-chart count and one $t$-sum\\
$3u$ & Positive determinant gap, finite factor family, all-aspect lattice count, axes, and inverse bounds\\
$3t$ & Determinant cutoff, fixed grid condition $N\ge40$, and
later lattice, factor, census, and inverse cutoffs\\
$4E$ & Four supported-factor sums, their fixed-height bounds, and the single shift sum\\
$4H$ & Reducible, content, and primitive cubic branches, including the primitive-form estimate\\
$4P$ & Graph, high-degree, cover, and axis bounds, with simultaneous shifts counted once\\
$5W$ & Lemma~\ref{lem:ratio-cell-determinant}, finite chart comparisons,
original-grid census, factor branches, endpoint bounds, assigned-corner
cases, and the geometric envelope; in particular $6E_Y<\xi/2$ and
$576Y^{q_{\max}+2\lambda-\tau_0}<1$ in \eqref{eq:old:7_15}\\
\end{tabular}
\end{table}

\begin{samepage}
In the \(5W\) row, the geometric cutoff is
\[
T_{\rm geom}=\frac{252}{g_0}+1.
\AutoTag{eq:old:9_3}
\]
\end{samepage}%
For \(\log Y>T_{\rm geom}\), \(\log2<7/10\) gives
\[
\frac{90\log2}{\log Y}
<\frac{63}{252/g_0}=\frac{g_0}{4}.
\AutoTag{eq:old:9_4}
\]
The high-degree correction in \eqref{eq:old:7_25} is separate from the geometric
correction. Since \(a_{\rm grid}>1/2\), for \(Y\ge2\),
\[
\lambda+\frac{c_N}{2}
\le \frac{Y^{-1}+\tfrac12Y^{-1/2}}{\log Y}
<\frac1{\log Y}.
\]
Also, \(c_N=\log_Y(1+Y^{-a_{\rm grid}})
\le Y^{-a_{\rm grid}}/\log Y\), so the interval correction is smaller
than \(1/\log Y\).
Consequently
\[
\frac{E_q+c_N}{2}+12d_2
\le\frac{\tfrac12\log18+12\log2+1}{\log Y}
<\frac{90\log2}{\log Y}<\frac{g_0}{4}.
\]
Here \(\tfrac12\log18+12\log2=9.763\ldots\). Thus each of the geometric
correction \(6\rho_Y\), the interval correction \(c_N\), and the high-degree
correction \((E_q+c_N)/2+12d_2\) is at most \(g_0/4\). The final row estimate
is charged their maximum, not their sum.
We retain this conservative charge in the reserve \eqref{eq:old:8_4};
it is unnecessary for the exponent in Proposition~\ref{prop:row-5w},
where the three finite corrections are absorbed into the implied constant.
The determinant, census and geometric-envelope thresholds in that row
are the conditions used in Section~\ref{sec:long}. The
contribution \(90\) is the source-mass \(5\) plus the one geometric
collar mass \(1\), each times the coefficient \(15\) in the
row margin.

Let \(Y_i(\beta_0)\) be the maximum of the finitely many thresholds
used in the proof of row \(i\), after applying \eqref{eq:old:9_1}.
Define
\[
T_0=\max\{3,T_{\rm geom},\log Y_1(\beta_0),\ldots,
                         \log Y_{13}(\beta_0)\},\qquad
Y_{\rm final}=\lceil e^{T_0}\rceil+1.
\AutoTag{eq:old:9_5}
\]
Since \(T_0\ge3\), \(Y_{\rm final}\ge\lceil e^3\rceil+1=22>4\).
Thus the elementary conditions \(Y\ge4\) of the 5W accounting
hold at the common cutoff.
The maximum runs over the thirteen fixed row estimates and their fixed
finite degree families. It does not run over moving source labels. It
includes the separate post-determinant thresholds of the \(3t\) row.

\begin{proposition}[Common-loss assembly]\label{prop:commonloss}
For \(Y\ge Y_{\rm final}\), each row \(i\ne5W\) is
\[
\#\mathrm{row}_i\ll_{\beta_0}
Y^{2/3-r_i+\beta_0},
\qquad r_i\ge S_0,
\]
and the last row satisfies
\[
\#5W\ll_{\beta_0}
Y^{2/3-S_0+\beta_0+g_0/4}.
\]
Consequently \(\#\Omega_Y\ll_\varepsilon
Y^{109/165+\varepsilon}\).
\end{proposition}

\begin{proof}
The first twelve row estimates are complete counts. They include the
moving labels, shifts, coefficient factorizations and inverse
multiplicities. By Section~\ref{sec:ledger}, their minimum saving is
\(S_0\). By \eqref{eq:old:7_26}, the three limiting values
\(L_{5W},s_{\rm int},s_{\rm HD}\) in the definition of \(\sigma_{5W}\)
are each at least \(S_0\). Hence the corrected minimum satisfies
\[
 \sigma_{5W}\ge S_0-\max\!\left(6\rho_Y,\ c_N,\
       \frac{E_q+c_N}{2}+12d_2\right).
\]
The three corresponding source-box contributions are added in
\eqref{eq:old:7_25}, so this minimum controls their sum exponent.
The corrections have separate sources. First, \(6\rho_Y\) is the finite
source-to-lower-corner allowance in \eqref{eq:old:7_22}--\eqref{eq:old:7_23c}.
Next, \(c_N\) is the ceiling correction from \(N=\lceil Y^{a_{\rm grid}}\rceil\)
in \eqref{eq:old:7_6}. Last, \((E_q+c_N)/2+12d_2\) bounds the finite
square-root term \(\sqrt{4MN}\). This uses the source-product and
grid-ceiling bounds just before \eqref{eq:old:7_25}. The margins add no
losses beyond these three terms.
By \eqref{eq:old:9_4} and the correction bounds above, each of the three
corrections is at most \(g_0/4\) for \(Y\ge Y_{\rm final}\). Hence
\(\sigma_{5W}\ge S_0-g_0/4\). Apply \eqref{eq:old:7_25} with
\(\eta=\beta_0\), using the remaining \(\beta_0/2\) for the box sum.
Thus the total count is
\[
\ll_{\beta_0}
Y^{2/3-S_0+\beta_0+g_0/4}
=Y^{109/165-11g_0/16}\ll Y^{109/165}.
\]
There are thirteen rows. Their finite sum changes only the implied
constant.
\end{proof}

\section{Dyadic shells}\label{sec:conclusion}

For \(Y=2^i\), the shells \(2^{i-1}<n\le2^i\) are disjoint. By
Proposition~\ref{prop:commonloss}, for every \(i\) with \(2^i\ge
Y_{\rm final}\),
\[
\#\Omega_{2^i}\ll_\varepsilon 2^{i(109/165+\varepsilon)}.
\AutoTag{eq:old:10_1}
\]
Let \(I\) be the least integer with \(X\le2^I\). Then \(2^I<2X\) and
\[
\sum_{i\le I}2^{i(109/165+\varepsilon)}
\ll_\varepsilon (2X)^{109/165+\varepsilon}.
\AutoTag{eq:old:10_2}
\]
The finitely many shells below \(Y_{\rm final}\) are absorbed by
\[
\#\Omega_Y\le \frac{Y(Y-1)}2
\le\frac12\max(1,Y_{\rm final}^{2-\alpha})Y^\alpha,
\qquad \alpha=109/165+\varepsilon.
\]
Each pair counted by \(F(X)\) lies in exactly one shell. This proves
Theorem~\ref{thm:main}.

\appendix
\section{Boundary rules for the first-match partition}\label{app:equality}
\begin{samepage}
All cuts in the partition apply to the original pair. The table gives
the first-match row for each equality, after all earlier rows are tested.
The notation is that of Section~\ref{sec:partition}.
\end{samepage}
\[
\begin{array}{c|c}
\text{boundary}&\text{first-match row}\\ \hline
R=Y^{5/3-\epsilon_0}&O_R\\
M=Y^{2/3-\epsilon_0}&O_M\\
\rad(M)=Y^{2/3-k_0}&O_{\rad M}\\
K_{\rm res}=Y^{\Delta_{\rm exp}}&\text{fails }Q0;\text{ continue testing}\\
L=Y^{1/3-l_0}&\text{fails }Q1;\text{ continue testing}\\
\min(J,K)=Y^{1/3-j_0}&\text{fails strict }Q2;\text{ test later rows}\\
u=t&u\text{-chart}\\
s=Y^{a_0}&3u\text{ or }3t,\text{ as selected}\\
s=Y^{b_0}&5W\\
\min(v,w)=Y^{p_0}&4E\\
\Pi=Y^{p_0},\quad v,w>Y^{p_0}&4H\\
\Pi>Y^{p_0},\quad v,w>Y^{p_0}&4P\\
\end{array}
\AutoTag{eq:old:A_1}
\]
The rows for \(4E,4H,4P\) require that all earlier predicates fail and
that \(Y^{a_0}<s<Y^{b_0}\). Equality at the lower-corner boundary fails the
strict \(Q2\) cut. It then goes to the first later row whose predicate
holds. The weak lower-corner inequalities describe the complements of the
strict \(Q2\) cuts.

The row predicates are Boolean functions of the original pair. A
counting argument may round an interval endpoint outward to an integer.
This only enlarges the counted set. No real rounding enters the
definition of a source row. Take the explicit inverse in
\eqref{eq:old:2_2}. Then test positivity, integrality, the gcd, the shell
and the mask. Each accepted reconstructed tuple then gives at most one
original pair.

Fix the ordered residual tuple that determines \(L\). Then the
identity \(ut=L\) determines the opposite shift. Hence a sum over \(u\) or
over \(t\) is a sum over one source coordinate. It is not a product
of independent shift ranges. In the \(4H\) and \(4P\) rows this identity
forces the opposite shift. In \(5W\) both shifts stay in the
original-cell source. Its geometric envelope adds no new shift sum.

The four $Q2$ count contributions fall into two populations:
\[
\begin{array}{c|c|c}
\text{population}&\text{contributions}&\text{selected chart}\\ \hline
Q2u&Q2u_{\rm int}+Q2u_{\rm harm}&u\le t\\
Q2t&Q2t_{\rm int}+Q2t_{\rm harm}&t<u
\end{array}
\]
Section~\ref{sec:ledger} computes their savings. The two terms on a
line are added to bound one chart population. They are not separate rows.
Their exponents and savings are
\[
\begin{array}{c|c|c|c}
\text{term}&\text{chart contribution}&\text{exponent}&\text{saving}\\ \hline
Q2t_{\rm int}&\text{interval term in }Q2t&1/2+6S_0&1/6-6S_0=972179/7500000\\
Q2t_{\rm harm}&\text{harmonic term in }Q2t&2/3-S_0&S_0\\
Q2u_{\rm int}&\text{interval term in }Q2u&1/2+3S_0&1/6-3S_0=2222179/15000000\\
Q2u_{\rm harm}&\text{harmonic term in }Q2u&2/3-4S_0&4S_0=30869/1250000
\end{array}
\AutoTag{eq:old:A_2}
\]
In each chart the larger exponent controls the sum. This is the smaller
saving. No term is an additional source population.

The partition is disjoint and exhaustive. Hence
\begin{align}
\#\Omega_Y={}&\#O_R+\#O_M+\#O_{\rad M}+\#Q0+\#Q1
 +\#Q2u+\#Q2t \notag\\
 &+\#3u+\#3t+\#4E+\#4H+\#4P+\#5W.
\AutoTag{eq:old:A_3}
\end{align}
The first-match convention is essential. A relaxed set that bounds one
row may overlap another relaxed counting set. Still, every original pair
goes to exactly one row. So the populations satisfy the exact
identity \eqref{eq:old:A_3}.

\section{A uniform square-fibre bound}\label{app:fibre}

Tao's Lemma~2.10 \cite[Lemma~2.10]{Tao} gives a related square-fibre
bound. Appendix~\ref{app:sources} states its hypotheses. The self-contained
theorem below gives the
fixed-\(\varepsilon\), fixed-polynomial-height version used here. It assumes
no uniform Pell-orbit or regulator bound.

\begin{theorem}[Uniform square fibres]\label{thm:fibre}
Fix \(A\ge1\). For every \(\varepsilon>0\) there is a constant
\(C_{A,\varepsilon}\) such that, for \(x\ge2\), positive integers
\(a,b\le x^A\), and a nonzero integer \(h\) with \(|h|\le x^A\),
\[
 \#\{(n,m)\in\mathbb Z_{>0}^2:
       an^2+h=bm^2,\ n\le x\}
 \le C_{A,\varepsilon}x^\varepsilon.
\AutoTag{eq:old:B_1}
\]
The constant is independent of \(a,b,h\). The exponent \(A\) is fixed.
\end{theorem}

\begin{lemma}[Divisor bound]\label{lem:tau}
For every \(\delta>0\), \(\tauDiv(N)\ll_\delta N^\delta\) for \(N\ge1\).
\end{lemma}
\begin{proof}
For primes \(p\ge2^{1/\delta}\), \(j+1\le2^j\le p^{\delta j}\).
There are finitely many smaller primes. For each of them
\(\sup_{j\ge0}(j+1)p^{-\delta j}<\infty\). Multiply over the
prime factorization.
\end{proof}

\begin{lemma}[Roots of a quadratic congruence]\label{lem:qroots}
If \(D\ge2\) is squarefree and \(M\ge1\), then
\[
\#\{r\bmod M:r^2\equiv D\pmod M\}\le4\cdot2^{\omega(M)}.
\AutoTag{eq:old:B_2}
\]
For \(M=1\) there is one residue class.
\end{lemma}
\begin{proof}
It suffices to count roots modulo prime powers. For an odd prime \(p\mid D\),
there is one root modulo \(p\) and none modulo \(p^e\), \(e\ge2\). Here
a root would have valuation at least one. Its square would then have
valuation at least two. If \(p\nmid D\), a root \(r_0\) is a unit. For any other
root \(r\), \(p^e\mid(r-r_0)(r+r_0)\). The two factors are not both
divisible by \(p\). So \(r\equiv r_0\) or \(-r_0\pmod{p^e}\).
There are at most two roots.

Let \(p=2\). If \(D\) is even, squarefreeness gives \(D\equiv2\pmod4\).
Then there are no roots modulo \(2^e\) for \(e\ge2\), and at most one modulo
2. Let \(D\) be odd and let \(r_0\) be one root. Multiplication by
\(r_0^{-1}\) injects the roots into solutions of \(z^2\equiv1\pmod{2^e}\).
For \(e=1,2\) there are at most two. Let \(e\ge3\). Then \(z\) is odd, and
exactly one of \(z-1,z+1\) is divisible by \(2\) but not \(4\).
Their product is divisible by \(2^e\). So the other is divisible by
\(2^{e-1}\). Hence \(z\equiv1\) or \(-1\pmod{2^{e-1}}\). This gives at most
four classes modulo \(2^e\). The Chinese remainder theorem multiplies
these local bounds, and \eqref{eq:old:B_2} follows.
\end{proof}

\begin{lemma}[Spacing in one norm class]\label{lem:normspacing}
Let \(D\ge2\) be squarefree, \(N\in\mathbb Z\setminus\{0\}\), and \(H\ge1\). The number of
positive integral solutions of
\[
 X^2-DY^2=N,\qquad X+Y\sqrt D\le H
\]
is at most
\[
 4\tauDiv(|N|)^2
 \left(1+\frac{\log H}{\log(2+\sqrt3)}\right).
\AutoTag{eq:old:B_3}
\]
\end{lemma}
\begin{proof}
Put \(g=\gcd(X,Y)\). Then \(g^2\mid N\). For each of at most
\(\tauDiv(|N|)\) choices of \(g\), set \(U=X/g\), \(V=Y/g\),
and \(N_0=N/g^2\). The pair \((U,V)\) is primitive. Also
\(\gcd(V,|N_0|)=1\). Otherwise, by the norm equation, a common prime
would divide \(U^2\), hence \(U\). This contradicts primitivity. Thus
\(r=UV^{-1}\pmod{|N_0|}\) is defined and satisfies \(r^2\equiv D\).
For modulus one use its unique class. By \eqref{eq:old:B_2}, there are at most
\(4\tauDiv(|N|)\) such classes.

Take two solutions in one class. Write
\(\alpha_i=U_i+V_i\sqrt D>1\). Direct multiplication gives
\[
 \frac{\alpha_1}{\alpha_2}=s+t\sqrt D,\quad
 s=\frac{U_1U_2-DV_1V_2}{N_0},\quad
 t=\frac{V_1U_2-U_1V_2}{N_0}.
\AutoTag{eq:old:B_4}
\]
Here \(s,t\) are the integer coefficients of this norm-one quotient.
This local \(t\) is unrelated to either source shift.
Both numerators are divisible by \(N_0\). To see this, work modulo \(|N_0|\) and use
\(U_i\equiv rV_i\) and \(r^2\equiv D\). Hence \(s,t\) are integers
even when \(N_0<0\). The quotient has norm one. Let
\(\alpha_1/\alpha_2>1\). Its conjugate is the positive reciprocal. Thus
\(s=(\alpha_1/\alpha_2+\alpha_2/\alpha_1)/2>1\). Since \(s\) is an
integer, \(s\ge2\), and
\[
 \frac{\alpha_1}{\alpha_2}=s+\sqrt{s^2-1}\ge2+\sqrt3.
\]
Distinct pairs have distinct \(\alpha_i\), as \(\sqrt D\) is irrational.
Hence each class contains at most
\(1+\log H/\log(2+\sqrt3)\) solutions. Multiply by the choices of
\(g\) and of the residue class to get \eqref{eq:old:B_3}.
\end{proof}

\begin{proof}[Proof of Theorem~\ref{thm:fibre}]
Write \(ab=Dc^2\) with \(D\) squarefree and put
\[
 X=bm,\qquad Y_0=cn,\qquad N=bh.
\]
Then
\[
 X^2-DY_0^2=b^2m^2-abn^2=bh=N\ne0.
\AutoTag{eq:old:B_5}
\]
This map is injective. From the original equation and \(n\le x\),
\[
 bm\le\sqrt{abx^2+b|h|}\le\sqrt2\,x^{A+1},\qquad
 cn\sqrt D=\sqrt{ab}\,n\le x^{A+1}.
\]
Thus \(X+Y_0\sqrt D\le3x^{A+1}\) and \(|N|\le x^{2A}\).

If \(D\ge2\), apply Lemma~\ref{lem:normspacing}. Then Lemma~\ref{lem:tau}
absorbs the two divisor factors and the logarithm into \(x^\varepsilon\). Its
exponent is chosen in terms of \(A,\varepsilon\). If \(D=1\), \eqref{eq:old:B_5} factors
as \((X-Y_0)(X+Y_0)=N\). Each positive divisor \(X+Y_0\) of
\(|N|\) fixes the signed factor \(X-Y_0=N/(X+Y_0)\). So it gives at
most one pair. The divisor bound then proves \eqref{eq:old:B_1} in this case too.
\end{proof}

\begin{remark}
Every use has positive coefficients and a nonzero shift. The
coefficients and the shift have fixed polynomial height. One variable is
bounded, and the equation then bounds the other. No
discriminant-dependent constant or fundamental-unit estimate occurs.
\end{remark}

In the conic uses of Sections~\ref{sec:intermediate} and~\ref{sec:long},
put $(a,b,h,n,m)=(f,g,s,J,K)$. The source bounds give
$1\le f,g,s\le (Y+1)^2$ and $J,K\le Y+1$. Take $x=2Y$ and the fixed
height exponent $A_{\rm ht}=2$, with theorem parameter $A=A_{\rm ht}$. All
height hypotheses hold, since $(Y+1)^2\le(2Y)^2$ and $Y+1\le2Y$.
The coefficients are positive and the selected shift satisfies $s>0$.
Suppose the factorization does not fix the shift. Then we first fix
the ordered residual factors, which determine $L$. Next we enumerate $s\mid L$.
As $L\le(Y+1)^2$, this costs at most $\tauDiv(L)\ll_\delta Y^\delta$.
Thus the theorem is applied once per fixed nonzero shift, with its uniform
constant. The shift divisor count is part of the allocated loss.

\section{Application of the low-radical estimate}\label{app:lowrad}

Theorem~3 of \cite{Paper1}, stated in Appendix~\ref{app:sources},
holds uniformly for real $H\ge2$ and $1\le V<H^2$.
Its determinant-method input is Reuss's Theorem~1 \cite{Reuss}; we use
the consequence stated in Paper 1.

For the outer row $O_R$, set
\[
 H=Y/2,\qquad V=Y^{5/3-S_0},\qquad
 S_0=30869/5000000.
\]
For $Y\ge64$, $H\ge2$, $V\ge1$, and
\[
 \frac{V}{H^2}=4Y^{-1/3-S_0}<1.
\]
Hence that theorem applies. Its two terms satisfy
\[
 \frac VH=2Y^{2/3-S_0},\qquad
 H^{7/12}=2^{-7/12}Y^{7/12}
 \le Y^{2/3-S_0},
\]
because
\[
 \frac23-S_0-\frac7{12}
 =\frac{1157393}{15000000}>0.
\]
Keep the common factor \(H^\eta=(Y/2)^\eta\). Then
\[
 (Y/2)^\eta\left(2Y^{2/3-S_0}
       +2^{-7/12}Y^{7/12}\right)
 \ll_\eta Y^{2/3-S_0+\eta}.
\]
Consequently
\[
 \#\{Y/2<t\le Y:\operatorname{rad}(t(t+1))\le Y^{5/3-S_0}\}
 \ll_\eta Y^{2/3-S_0+\eta},
\]
with one constant for every $Y\ge64$. This is the low-radical input of
Section~\ref{sec:outer}. The partner count there separately uses
Lemma~2.1 of Paper~1 \cite{Paper1}.

\section{Exact exponents and finite optimization}
\label{app:arithmetic}

All numbers in this appendix are exact rationals. Set
\[
\begin{gathered}
S_0=\epsilon_0=k_0=\frac{30869}{5000000},\qquad
l_0=2S_0=\frac{30869}{2500000},\qquad
j_0=p_0=3S_0=\frac{92607}{5000000},\\
a_0=\frac6{125},\qquad b_0=\frac{4166537}{30000000},\qquad
\Delta_{\rm exp}=3S_0=\frac{92607}{5000000},\\
\mu_{\rm HP}=\frac{16}{25},\qquad \nu=\frac1{64},\qquad
\xi=\frac1{10000}.
\end{gathered}
\AutoTag{eq:old:D_1}
\]
The target saving and reserve are
\[
S_{\rm target}=\frac1{165},\qquad
g_0=S_0-S_{\rm target}=\frac{18677}{165000000}.
\AutoTag{eq:old:D_2}
\]

The table gives every source-count contribution. The margin is
the displayed saving minus $S_0$. A zero margin marks a binding contribution.
Two chart terms of one population are both listed, since their counts
add. The smaller saving controls their sum. The three
displayed $5W$ savings are additive contributions to the complete
count. Their minimum controls the exponent of the sum. The chart and
assigned-corner routes are selected in each source box as in
Section~\ref{sec:long}.
\small
\unnumtable
\newcommand{\fivewdecimal}[1]{\text{\scriptsize (#1)}}
\begin{longtable}{@{}p{0.25\textwidth}p{0.36\textwidth}p{0.32\textwidth}@{}}
\toprule
Contribution & Saving & Margin above $S_0$\\
\midrule
\endhead
$O_R$ & $30869/5000000$ & $0$\\
$O_M$ & $30869/5000000$ & $0$\\
$O_{\operatorname{rad}M}$ & $30869/5000000$ & $0$\\
$Q0$ & $2/3$ & $9907393/15000000$\\
$Q1$ & $30869/5000000$ & $0$\\
$Q2t_{\rm int}$ & $972179/7500000$ & $1851751/15000000$\\
$Q2t_{\rm harm}$ & $30869/5000000$ & $0$\\
$Q2u_{\rm int}$ & $2222179/15000000$ & $532393/3750000$\\
$Q2u_{\rm harm}$ & $30869/1250000$ & $92607/5000000$\\
$3u_{\rm sqrt}$ & $5102312049347/600000000000000$ & $1398032049347/600000000000000$\\
$3u_{\rm int}$ & $23206592049347/300000000000000$ & $21354452049347/300000000000000$\\
$3u_{\rm core}$ & $2473/12000$ & $2998643/15000000$\\
$3u_{\rm small,1}$ & $1619393/6000000$ & $7911751/30000000$\\
$3u_{\rm small,2}$ & $8003/500000$ & $49161/5000000$\\
$3t_{\rm sqrt}$ & $388751/60000000$ & $18323/60000000$\\
$3t_{\rm int}$ & $2199179/30000000$ & $402793/6000000$\\
$3t_{\rm axes}$ & $2726537/15000000$ & $263393/1500000$\\
$4E$ & $30869/5000000$ & $0$\\
$4H$ & $140441/5000000$ & $27393/1250000$\\
$4P_{\rm graph}$ & $30869/5000000$ & $0$\\
$4P_{\rm HD}$ & $107393/15000000$ & $7393/7500000$\\
$4P_{\rm cover}$ & $2/75$ & $307393/15000000$\\
$4P_{\rm axes}$ & $4166537/30000000$ & $3981323/30000000$\\
$5W_{\rm geom}$ &
$\begin{gathered}\dfrac{1203113}{90000000}\\[-2pt]\fivewdecimal{0.013368}\end{gathered}$ &
$\begin{gathered}\dfrac{647471}{90000000}\\[-2pt]\fivewdecimal{0.007194}\end{gathered}$\\
$5W_{\rm int}$ &
$\begin{gathered}\dfrac{235120131355662893}{9520800917171070000}\\[-2pt]\fivewdecimal{0.024695}\end{gathered}$ &
$\begin{gathered}\dfrac{88170305326616070517}{4760400458585535000000}\\[-2pt]\fivewdecimal{0.018522}\end{gathered}$\\
$5W_{\rm HD}$ &
$\begin{gathered}\dfrac{29390272487700347267}{4760400458585535000000}\\[-2pt]\fivewdecimal{0.006174}\end{gathered}$ &
$\rule{0pt}{3.2ex}\begin{gathered}\dfrac{128034121242821}{1190100114646383750000}\\[-2pt]\text{\scriptsize (1.076\ensuremath{\times10^{-7}})}\end{gathered}$\\
\bottomrule
\end{longtable}
\normalsize

Decimals in parentheses are rounded annotations. All comparisons
and margins use the exact rational entries.

Thus the seven binding source contributions are
\[
O_R,\ O_M,\ O_{\operatorname{rad}M},\ Q1,\ Q2t_{\rm harm},\ 4E,\ 4P\text{ graph},
\]
each with saving exactly $S_0$. The high-degree term limits the $3u$
aggregate. The high-degree term also limits the $3t$ aggregate.
The graph term limits the $4P$ aggregate. For $5W$, the minimum
saving among the geometric, interval and high-degree contributions controls
the combined count in \eqref{eq:old:7_25}.

\subsection*{The limiting linear program}

Let $x_0,\ldots,x_7$ be the base-$Y$ logarithms of the eight assigned
corner factors, in the order
\[
(a_A,a_B,b_C,b_D,c_A,c_C,d_B,d_D).
\]
Put
\[
\begin{gathered}
\alpha=x_0+x_1+x_2+x_3,\quad
\gamma=x_4+x_5+x_6+x_7,\quad q=\alpha+\gamma,\\
c_1=x_0+x_1+x_4+x_5,\quad c_2=x_0+x_1+x_6+x_7,\\
P=x_0+x_4,\quad Q=x_3+x_7,\quad R_{\rm LP}=x_1+x_6,\quad T=x_2+x_5.
\end{gathered}
\AutoTag{eq:old:D_3}
\]
The other variables are $W,z_u,z_t,\mu\ge0$. Here $W$ is a
candidate saving, $z_u=\log_Yu$, $z_t=\log_Yt$, and
$\mu=\log_Ym$. Every branch uses the limiting source constraints below.
They are not finite-source identities. Section~\ref{sec:long} records the
finite relations and their collars.
\[
\begin{gathered}
\frac{2407393}{7500000}\le\alpha\le\frac{5092607}{15000000},\qquad
\frac{9907393}{30000000}\le\gamma\le\frac{2092607}{6000000},\\
\frac{9907393}{15000000}\le q\le\frac{5092607}{7500000},\qquad
\frac{2407393}{7500000}\le z_u+z_t\le\gamma,\\
\mu+\frac{\gamma-\alpha}{2}=1,\qquad
\mu\le1,\qquad \mu\ge\frac{9907393}{15000000}+\frac q2,\\
z_u,z_t\ge b_0,\qquad z_u+z_t\le\frac{5092607}{7500000}.
\end{gathered}
\AutoTag{eq:old:D_4}
\]
The four assigned-corner capacity inequalities are
\begin{align}
\frac32x_0+\frac12x_1+\frac32x_4+\frac12x_5
 -\frac12z_u-\frac12\mu&\le\frac{92607}{5000000},\label{eq:old:D_capA}\\
\frac12x_2+\frac32x_3+\frac12x_6+\frac32x_7
 -\frac12z_u-\frac12\mu&\le\frac{92607}{5000000},\label{eq:old:D_capD}\\
\frac12x_0+\frac32x_1-\frac12x_4-\frac12x_5+x_6
 +\frac12z_u-\frac12\mu&\le\frac{92607}{5000000},\label{eq:old:D_capB}\\
-\frac12x_0-\frac12x_1+x_2+\frac12x_4+\frac32x_5
 +\frac12z_u+\frac12\mu&\le\frac{5092607}{5000000}.\label{eq:old:D_capC}
\end{align}
The selected-chart condition is $z_u\le z_t$ in chart 1 and
$z_t\le z_u$ in chart 2. In chart 1 also impose
\[
\mu+z_u-c_1\ge\frac{4722179}{7500000},\qquad
\mu+z_u-(x_2+x_3+x_6+x_7)\ge\frac{4722179}{7500000};
\AutoTag{eq:old:D_5}
\]
in chart 2 impose
\[
x_0+x_1-x_4-x_5+z_u-\mu\le-\frac{4722179}{7500000},\qquad
-x_0-x_1+x_4+x_5+z_u+\mu\le\frac{10277821}{7500000}.
\AutoTag{eq:old:D_6}
\]

The four assigned-corner route scores are the $G_A,G_D,G_B,G_C$ in
\eqref{eq:old:7_20}. A branch also imposes $W\ge0$ and
$W\ge G_A,G_D,G_B,G_C$. For $i,j\in\{0,1,2,3,4,5\}$ impose
$W\ge C_{1,i}$ and $W\ge C_{2,j}$. The six values in each chart are
defined by
{\small
\[
\begin{array}{c|l|c|l}
i&C_{1,i}&j&C_{2,j}\\ \hline
0&2/3-c_1&0&2/3-c_2\\
1&2/3-q+c_1&1&2/3-q+c_2\\
2&9722179/15000000-x_1-x_2-x_5-x_6&2&9722179/15000000-x_0-x_3-x_4-x_7\\
3&9722179/15000000-x_1-x_3-x_5-x_7&3&9722179/15000000-x_0-x_2-x_5-x_7\\
4&9722179/15000000-x_0-x_2-x_4-x_6&4&9722179/15000000-x_1-x_3-x_4-x_6\\
5&9722179/15000000-x_0-x_3-x_4-x_7&5&9722179/15000000-x_1-x_2-x_5-x_6
\end{array}
\AutoTag{eq:old:D_7}
\]
}
\subsection*{Derivation of the 22 limiting source rows}
The source rows are R1--R6 and R9--R24 in the certificate order
below. The LP linear forms \(P,Q\) are those of \eqref{eq:old:D_3}. For the
source products write \(P_{\rm src}=ab\) and \(Q_{\rm src}=cd\), and
write \(L=Q_{\rm src}-P_{\rm src}\). Set
\(\ell_{\rm src}=\log_Y L\),
\(\alpha_{\rm src}=\log_Y P_{\rm src}\),
\(\gamma_{\rm src}=\log_Y Q_{\rm src}\), and
\(q_{\rm src}=\alpha_{\rm src}+\gamma_{\rm src}=\log_Y M\).
Write \(\ell=z_u+z_t\) for the limiting LP coordinate. The finite logarithmic collars come from
\(\kappa_Y=1+Y^{-1}\), from factors of two and from the endpoint terms in
\(Z^2M=m(m+1)n(n+1)\). They are included in
\eqref{eq:old:7_23b}--\eqref{eq:old:7_23c}. The LP rows are their
limiting affine inequalities. The table writes them in limiting LP
endpoint coordinates \(q,\ell,z_u,z_t\), distinct from the finite source
logarithms with subscript \({\rm src}\). The limiting coordinates are
denoted by \(q\) and \(\ell\); \(z_u,z_t\)
are the LP variables corresponding to the finite logarithms
\(\log_Yu,\log_Yt\). The subscript \({\rm src}\) marks source quantities
before taking limits. Labels R1--R28 refer throughout to the certificate
row numbering below.
\small
\unnumtable
\begin{longtable}{@{}p{0.09\textwidth}p{0.84\textwidth}@{}}
\toprule
Limiting LP row & Source derivation \\
\midrule
\endhead
R1 & If \(P_{\rm src}\le L\), then \(P_{\rm src}\le L<\kappa_YY^{1/3+S_0}\).
If \(P_{\rm src}>L\), then \(P_{\rm src}<\sqrt M<\kappa_YY^{q_{\max}/2}\).
Since \(q_{\max}/2=1/3+S_0\), these give the displayed limiting upper bound. \\
R2 & If \(P_{\rm src}\le L\), then \(Q_{\rm src}\le2L\), so
\(P_{\rm src}=M/Q_{\rm src}>Y^{q_{\min}}/(2\kappa_YY^{1/3+S_0})\),
whose limiting exponent is \(1/3-2S_0\). If \(P_{\rm src}>L\), then
\(Q_{\rm src}<2P_{\rm src}\) and \(M<2P_{\rm src}^2\), giving the
stronger limiting lower exponent \(q_{\min}/2=1/3-S_0/2\). \\
R3 & If \(P_{\rm src}\le L\), then \(Q_{\rm src}\le2L\); otherwise
\(Q_{\rm src}<2P_{\rm src}<2\sqrt M\). Thus the source gives
\(\gamma_{\rm src}\le1/3+S_0+\log_Y(2\kappa_Y)\); for sufficiently large \(Y\)
this implies the weaker \(\gamma_{\rm src}\le1/3+5S_0/2\). \\
R4 & Since \(P_{\rm src}<Q_{\rm src}\), we have
\(\alpha_{\rm src}<\gamma_{\rm src}\); with
\(q_{\rm src}=\alpha_{\rm src}+\gamma_{\rm src}\ge q_{\min}\), this gives
\(\gamma_{\rm src}\ge q_{\min}/2\). \\
R5 & The failed \(O_R\) cut gives \(M<\kappa_Y^2Y^{q_{\max}}\), hence
\(q_{\rm src}\le q_{\max}+2\lambda\), where \(\lambda=\log_Y\kappa_Y\);
the LP uses the limiting bound \(q\le q_{\max}\). \\
R6 & The failed \(O_M\) cut gives \(M>Y^{q_{\min}}\), hence
\(q_{\rm src}\ge q_{\min}\). \\
R9 & Failure of \(Q1\) gives \(L\ge Y^{1/3-l_0}\), so
\(\ell_{\rm src}\ge1/3-l_0\). \\
R10 & Since \(L=Q_{\rm src}-P_{\rm src}<Q_{\rm src}\), taking logarithms gives
\(\ell_{\rm src}\le\gamma_{\rm src}\). \\
R11 & The mask gives \(a_Ac_A\le A K_{\rm res}\). Taking logarithms
and replacing the finite \(\log_Y A\) by its lower-corner reference
\(\mathsf A=(\mu+z_u-c_1)/2\) from \eqref{eq:old:7_23a}--\eqref{eq:old:7_23b}
gives \(x_0+x_4-\mathsf A\le\Delta_{\rm exp}\), which is
\eqref{eq:old:D_capA}. \\
R12 & The mask gives \(b_Dd_D\le D K_{\rm res}\). Replacing the finite
\(\log_YD\) by its lower-corner reference
\(\mathsf D=(\mu+z_u-q+c_1)/2\) from \eqref{eq:old:7_23a}--\eqref{eq:old:7_23b}
gives \(x_3+x_7-\mathsf D\le\Delta_{\rm exp}\), which is
\eqref{eq:old:D_capD}. \\
R13 & The source B-corner capacity is \(a_Bd_B\le B K_{\rm res}\).
Taking logarithms and using the limiting corner form
\(\mathsf B=1/2-z_u/2+q/4-c_2/2\), where
\(R_{\rm LP}=x_1+x_6\) and \(c_2=x_0+x_1+x_6+x_7\), gives
the limiting inequality below. Here \(q\) denotes the LP coordinate
corresponding to the source exponent \(q_{\rm src}\) defined above; its
finite endpoint error is included in the source collar:
\[
R_{\rm LP}+z_u/2-q/4+c_2/2\le1/2+\Delta_{\rm exp}.
\]
Use \(1/2=\mu/2+(\gamma-\alpha)/4\), from the
limiting source equality, and expand. This gives row R13, namely
\eqref{eq:old:D_capB}. \\
R14 & The source C-corner capacity is \(b_Cc_C\le C K_{\rm res}\), so
\(T+z_u/2+q/4-c_2/2\le1/2+\Delta_{\rm exp}\), where
\(T=x_2+x_5\). Using the limiting equality first rewrites this as
\[
T+z_u/2+\alpha/2-c_2/2-\mu/2\le\Delta_{\rm exp}.
\]
Now eliminate \(\alpha=\gamma+2\mu-2\) and expand.
This gives row R14, namely \eqref{eq:old:D_capC} with right side
\(1+\Delta_{\rm exp}=5092607/5000000\); its finite endpoint error is
included in the source collar. Using \eqref{eq:old:7_23a}, after the
\(h_m-h_n\) terms cancel, the finite right sides in rows R13 and R14
are respectively
\[
 \Delta_{\rm exp}+\frac{H-\nu_n}{2},\qquad
 1+\Delta_{\rm exp}+\frac{H+\nu_n-2}{2}.
\]
After removing the limiting row constants, each endpoint residual is at
most \(d_2/2\) before lower-corner binning; that binning is included in
\(\rho_Y\). \\
R15--R16 & Put \(\nu_n=\log_Yn\) and
\(h_r=\log_Y(1+1/r)\). The exact source identity is
\[
\mu+(\gamma_{\rm src}-\alpha_{\rm src})/2
=\nu_n+(h_n-h_m)/2,
\]
obtained by taking logarithms of
\(Q_{\rm src}m(m+1)=P_{\rm src}n(n+1)\). Since \(Y/2<n\le Y\),
\(-d_2<\nu_n-1\le0\), and \(|h_n-h_m|\le\ell_{\rm end}\).
The finite discrepancy is absorbed in the source collar. In the limiting
LP, row R15 is
\(\mu+(\gamma-\alpha)/2\le1\), and row R16 is
the reverse inequality
\(\mu+(\gamma-\alpha)/2\ge1\). \\
R17 & The source has \(m<n\le Y\), hence \(\mu\le1\). \\
R18 & Since \(m+1\le2m\), \(n+1\le2Y\), and
\(Z>Y^{5/3-S_0}\), the identity
\(Z^2M=m(m+1)n(n+1)\) gives
\(Y^{10/3-2S_0+q_{\rm src}}<4m^2Y^2\), hence the finite bound
\(\mu>q_{\min}+q_{\rm src}/2-d_2\), where \(d_2=\log_Y2\). The LP uses the
limiting weak inequality \(\mu\ge q_{\min}+q/2\); the finite \(d_2\)
loss is included in the source collar \(\rho_Y\). \\
R19--R20 & The 5W predicate has \(u,t\ge Y^{b_0}\), hence
\(z_u,z_t\ge b_0\). \\
R21 & The source has \(L<\kappa_YY^{1/3+S_0}\), so its limiting
bound is \(\ell_{\rm src}\le1/3+S_0\). The displayed LP bound
\(\ell\le q_{\max}\)
is a redundant, weaker relaxation; it only enlarges the feasible set and
can lower the LP minimum. \\
R22 & The selected-chart rule \eqref{eq:old:2_3} gives \(z_u\le z_t\) in
chart 1 and \(z_t\le z_u\) in chart 2. \\
R23--R24 & The strict \(Q2\) cuts have failed. In chart 1, the complement
is the lower-corner inequality \eqref{eq:old:D_5}; in chart 2, it is
\eqref{eq:old:D_6}. \\
\bottomrule
\end{longtable}
\normalsize

Rows R11--R14 are the four limiting logarithmic capacity inequalities.
They come from the two-sided assignment of each residual prime to its
unique eligible corner. Rows R13--R14 use equality rows R15--R16.
Rows R3 and R21 weaken stronger source bounds. So the
relaxed feasible set only grows. The six counting rows R7, R8 and
R25--R28 come from the chart costs \eqref{eq:old:7_19}, the route scores
\eqref{eq:old:7_20}, the endpoint patterns
\eqref{eq:old:6_18a}, \eqref{eq:old:7_16}, and the route-score and
endpoint bounds \eqref{eq:old:7_23a}--\eqref{eq:old:7_23c}. They are not
extra source premises.

Fix an actual source. Keep its selected chart. Then choose a
minimizing form for each of the two counting orientations. Its endpoint
data then place it in one of the $2\cdot6\cdot6=72$ cases.
Dropping the inequalities that enforce the minimizing choices enlarges
that case. Hence a lower bound on every displayed relaxation applies to
the original source.

The certificates number the inequalities as follows.
\begin{itemize}
\item R1--R6: the upper and lower bounds for $\alpha$, then for
$\gamma$, then for $q$.
\item R7--R8: $W\ge C_{1,i}$ and $W\ge C_{2,j}$.
\item R9--R10: the lower bound for $z_u+z_t$ and its upper bound by $\gamma$.
\item R11--R14: the four capacity inequalities as displayed.
\item R15--R16: the two directions of the equality in \eqref{eq:old:D_4}.
\item R17--R18: the upper and lower bounds for $\mu$.
\item R19--R20: the lower bounds for $z_u,z_t$.
\item R21: the upper bound for $z_u+z_t$.
\item R22: the selected-chart inequality.
\item R23--R24: the two core inequalities for that chart.
\item R25--R28: $W\ge G_A,G_D,G_B,G_C$, in this order.
\end{itemize}
The six counting rows are R7, R8, and R25--R28.

Each branch has 28 inequalities. These are the displayed source bounds,
the two selected chart score rows from \eqref{eq:old:D_7} and the four
assigned-corner rows. The table below gives their exact
dual certificates. Write the rows as
$a_r\cdot x\le b_r$. The listed nonnegative multipliers satisfy
$\sum_r\lambda_r a_r\ge-e_W$ coordinatewise,
$\sum_{r\in\mathcal C}\lambda_r=1$ over the six counting rows, and
$-\sum_r\lambda_r b_r=W_{\sigma i j}$. As $x\ge0$, these relations
give $-W\le\sum_r\lambda_r a_r\cdot x\le\sum_r\lambda_r b_r$.
Thus every branch satisfies
\[
W\ge L_{5W},\qquad L_{5W}=\frac{1203113}{90000000}.
\AutoTag{eq:old:D_8}
\]
The value is attained in the relaxed branch $(1,2,5)$ by
\[
\begin{aligned}
&(x_0,\ldots,x_7,W,z_u,z_t,\mu)=\\
&\left(0,\frac{24907393}{90000000},\frac{5648249}{90000000},
0,0,\frac{2092607}{9000000},\right.\\
&\left.\frac{5648249}{90000000},\frac{3981323}{90000000},
\frac{1203113}{90000000},\frac{4166537}{30000000},
\frac{6018677}{30000000},1\right).
\end{aligned}
\AutoTag{eq:old:D_9}
\]
Exchange the two corner pairs. This gives a feasible point of the
chart-swapped branch $(2,5,2)$ with the same value. In
the coordinate order of \eqref{eq:old:D_3}, the permutation is
\[
(x_0,\ldots,x_7,z_u,z_t)\longmapsto
(x_1,x_0,x_3,x_2,x_6,x_7,x_4,x_5,z_t,z_u),
\]
with \(W\) and \(\mu\) fixed.
Consequently the exact optimization margin is
\[
\frac{1203113}{90000000}-S_0=\frac{647471}{90000000}>0.
\AutoTag{eq:old:D_10}
\]

The common-loss reserve is
\[
S_0-\frac{g_0}{16}-\frac{g_0}{4}
=\frac1{165}+\frac{11g_0}{16}>\frac1{165}.
\]
The threshold in Section~\ref{sec:epsilon} bounds each of the
$5W$ geometric, interval, and high-degree corrections by $g_0/4$.
The assembly charges their maximum, not their sum. The exact fractions and the certificate table give
the finite optimization explicitly.
\subsection*{Table of branch certificates}\label{app:branchcert}
In the table, $W_{\sigma,i,j}$ is the lower bound for branch
$(\sigma;i,j)$ in the lower-bound column. An entry $r:c$ means
$\lambda_r=c$ for the inequality row numbered above. Omitted
multipliers are zero.
A multiplier list may wrap onto continuation lines with blank branch
and lower-bound columns; all such lines belong to the same certificate.
\setlength{\tabcolsep}{2pt}
\scriptsize
\begin{longtable}{@{}p{0.15\textwidth}p{0.20\textwidth}p{0.60\textwidth}@{}}
\caption{Exact dual certificates for the 72 chart and score branches.}\label{tab:branchcert}\\
\toprule
Branch $(\mathrm{chart};i,j)$ & lower bound & Nonzero entries of $\lambda$\\
\midrule
\endfirsthead
\toprule
Branch $(\mathrm{chart};i,j)$ & lower bound & Nonzero entries of $\lambda$\\
\midrule
\endhead
$1;0,0$ & $4722179/30000000$ & \texttt{1:1/2; 5:1/2; 7:1/2; 8:1/2} \\
$1;0,1$ & $4722179/30000000$ & \texttt{5:1/2; 7:1/2; 8:1/2; 15:1/2; 18:1/2} \\
$1;0,2$ & $21943969/150000000$ & \texttt{7:3/5; 8:2/5; 10:1/10; 11:2/5; 15:1/10; 17:1; 18:9/10; 22:1/10} \\
$1;0,3$ & $3796109/30000000$ & \texttt{5:1/4; 7:1; 10:1/2; 15:1/2; 17:1/2; 22:1/2; 23:1} \\
$1;0,4$ & $3796109/30000000$ & \texttt{5:1/4; 7:1; 10:1/2; 15:1/2; 17:1/2; 22:1/2; 23:1} \\
$1;0,5$ & $3796109/30000000$ & \texttt{5:1/4; 7:1; 10:1/2; 15:1/2; 17:1/2; 22:1/2; 23:1} \\
$1;1,0$ & $4722179/30000000$ & \texttt{5:1/2; 7:1/2; 8:1/2; 15:1/2; 18:1/2} \\
$1;1,1$ & $4722179/30000000$ & \texttt{1:1/2; 5:1/2; 7:1/2; 8:1/2} \\
$1;1,2$ & $21943969/150000000$ & \texttt{1:1/10; 5:2/5; 7:3/5; 8:2/5; 10:1/10; 12:2/5; 15:1/5; 22:1/10} \\
$1;1,3$ & $3796109/30000000$ & \texttt{5:1/4; 7:1; 10:1/2; 15:1/2; 17:1/2; 22:1/2; 24:1} \\
$1;1,4$ & $3796109/30000000$ & \texttt{5:1/4; 7:1; 10:1/2; 15:1/2; 17:1/2; 22:1/2; 24:1} \\
$1;1,5$ & $3796109/30000000$ & \texttt{5:1/4; 7:1; 10:1/2; 15:1/2; 17:1/2; 22:1/2; 24:1} \\
$1;2,0$ & $21017899/150000000$ & \texttt{7:2/5; 8:3/5; 13:2/5; 15:1/5; 17:1; 18:1; 19:1/5} \\
$1;2,1$ & $21017899/150000000$ & \texttt{1:1/5; 7:2/5; 8:3/5; 14:2/5; 17:3/5; 18:4/5; 19:1/5} \\
$1;2,2$ & $1157393/3750000$ & \texttt{7:1/2; 8:1/2; 17:1; 18:1} \\
$1;2,3$ & $18424903/180000000$ & \texttt{1:1/6; 7:1/2; 8:1/2; 14:1/3; 17:5/6; 18:1; 19:1/6} \\
$1;2,4$ & $18424903/180000000$ & \texttt{7:1/2; 8:1/2; 13:1/3; 15:1/6; 17:7/6; 18:7/6; 19:1/6} \\
$1;2,5$ & $1203113/90000000$ & \texttt{8:1/3; 10:2/3; 15:2/3; 17:2/3; 18:4/3; 19:1; 27:1/3; 28:1/3} \\
$1;3,0$ & $3796109/90000000$ & \texttt{7:2/9; 8:1/3; 10:5/18; 13:2/9; 15:7/18; 17:7/9; 18:5/6; 22:5/18; 26:4/9} \\
$1;3,1$ & $3796109/90000000$ & \texttt{7:2/9; 8:1/3; 10:5/18; 14:2/9; 15:1/6; 17:7/9; 18:5/6; 22:5/18; 26:4/9} \\
$1;3,2$ & $4413489/40000000$ & \texttt{7:1/4; 8:1/4; 10:3/8; 15:3/8; 17:3/4; 18:7/8; 22:3/8; 26:1/2} \\
$1;3,3$ & $4166537/30000000$ & \texttt{7:1/2; 8:1/2; 15:1/2; 17:1; 18:3/2} \\
$1;3,4$ & $4166537/30000000$ & \texttt{7:1/2; 8:1/2; 16:1/2; 17:2; 18:3/2} \\
$1;3,5$ & $5462257/135000000$ & \texttt{7:1/3; 8:4/9; 10:2/9; 15:2/9; 17:11/9; 18:10/9; 23:1/3; 27:1/9; 28:1/9} \\
$1;4,0$ & $3796109/90000000$ & \texttt{7:2/9; 8:1/3; 10:5/18; 13:2/9; 15:7/18; 17:7/9; 18:5/6; 22:5/18; 25:4/9} \\
$1;4,1$ & $3796109/90000000$ & \texttt{7:2/9; 8:1/3; 10:5/18; 14:2/9; 15:1/6; 17:7/9; 18:5/6; 22:5/18; 25:4/9} \\
$1;4,2$ & $4413489/40000000$ & \texttt{7:1/4; 8:1/4; 10:3/8; 15:3/8; 17:3/4; 18:7/8; 22:3/8; 25:1/2} \\
$1;4,3$ & $4166537/30000000$ & \texttt{7:1/2; 8:1/2; 16:1/2; 17:2; 18:3/2} \\
$1;4,4$ & $4166537/30000000$ & \texttt{7:1/2; 8:1/2; 15:1/2; 17:1; 18:3/2} \\
$1;4,5$ & $5462257/135000000$ & \texttt{7:1/3; 8:4/9; 10:2/9; 15:2/9; 17:11/9; 18:10/9; 24:1/3; 27:1/9; 28:1/9} \\
$1;5,0$ & $3981323/90000000$ & \texttt{7:1/3; 10:1/2; 15:1/2; 17:2/3; 18:5/6; 22:1/2; 25:1/3; 26:1/3} \\
$1;5,1$ & $3981323/90000000$ & \texttt{7:1/3; 10:1/2; 15:1/2; 17:2/3; 18:5/6; 22:1/2; 25:1/3; 26:1/3} \\
$1;5,2$ & $3981323/90000000$ & \texttt{8:1/3; 10:1/2; 15:1/2; 17:2/3; 18:5/6; 22:1/2; 25:1/3; 26:1/3} \\
$1;5,3$ & $3981323/90000000$ & \texttt{7:1/3; 10:1/2; 15:1/2; 17:2/3; 18:5/6; 22:1/2; 25:1/3; 26:1/3} \\
$1;5,4$ & $3981323/90000000$ & \texttt{7:1/3; 10:1/2; 15:1/2; 17:2/3; 18:5/6; 22:1/2; 25:1/3; 26:1/3} \\
$1;5,5$ & $1157393/3750000$ & \texttt{7:1/2; 8:1/2; 17:1; 18:1} \\
$2;0,0$ & $4722179/30000000$ & \texttt{7:1/2; 8:1/2; 16:1/2; 17:2; 18:3/2} \\
$2;0,1$ & $4722179/30000000$ & \texttt{7:1/2; 8:1/2; 15:1/2; 17:1; 18:3/2} \\
$2;0,2$ & $21017899/150000000$ & \texttt{7:3/5; 8:2/5; 10:1/5; 11:2/5; 15:1/5; 17:1; 18:1; 20:1/5} \\
$2;0,3$ & $3796109/90000000$ & \texttt{7:1/3; 8:2/9; 10:1/6; 11:2/9; 15:1/6; 17:7/9; 18:5/6; 22:5/18; 28:4/9} \\
$2;0,4$ & $3796109/90000000$ & \texttt{5:7/18; 7:1/3; 8:2/9; 10:1/6; 11:2/9; 15:1/6; 18:1/18; 22:5/18; 27:4/9} \\
$2;0,5$ & $3981323/90000000$ & \texttt{8:1/3; 10:1/6; 15:1/6; 17:2/3; 18:5/6; 22:1/2; 27:1/3; 28:1/3} \\
$2;1,0$ & $4722179/30000000$ & \texttt{7:1/2; 8:1/2; 15:1/2; 17:1; 18:3/2} \\
$2;1,1$ & $4722179/30000000$ & \texttt{1:1/2; 5:1/2; 7:1/2; 8:1/2} \\
$2;1,2$ & $21017899/150000000$ & \texttt{5:1/2; 7:3/5; 8:2/5; 10:1/5; 12:2/5; 15:1/5; 20:1/5} \\
$2;1,3$ & $3796109/90000000$ & \texttt{5:7/18; 7:1/3; 8:2/9; 10:1/6; 12:2/9; 15:1/6; 18:1/18; 22:5/18; 28:4/9} \\
$2;1,4$ & $3796109/90000000$ & \texttt{5:7/18; 7:1/3; 8:2/9; 10:1/6; 12:2/9; 15:1/6; 18:1/18; 22:5/18; 27:4/9} \\
$2;1,5$ & $3981323/90000000$ & \texttt{8:1/3; 10:1/6; 15:1/6; 17:2/3; 18:5/6; 22:1/2; 27:1/3; 28:1/3} \\
$2;2,0$ & $5416537/37500000$ & \texttt{7:2/5; 8:3/5; 9:1/10; 13:2/5; 15:1/5; 17:1; 18:1; 22:1/10} \\
$2;2,1$ & $5416537/37500000$ & \texttt{7:2/5; 8:3/5; 9:1/10; 14:2/5; 16:1/5; 17:1; 18:1; 22:1/10} \\
$2;2,2$ & $1157393/3750000$ & \texttt{7:1/2; 8:1/2; 17:1; 18:1} \\
$2;2,3$ & $4413489/40000000$ & \texttt{7:1/4; 8:1/4; 10:1/8; 15:1/8; 17:3/4; 18:7/8; 22:3/8; 28:1/2} \\
$2;2,4$ & $4413489/40000000$ & \texttt{7:1/4; 8:1/4; 10:1/8; 15:1/8; 17:3/4; 18:7/8; 22:3/8; 27:1/2} \\
$2;2,5$ & $3981323/90000000$ & \texttt{8:1/3; 10:1/6; 15:1/6; 17:2/3; 18:5/6; 22:1/2; 27:1/3; 28:1/3} \\
$2;3,0$ & $219893/1875000$ & \texttt{5:1/2; 8:1; 9:1/2; 15:1; 22:1/2; 23:1} \\
$2;3,1$ & $219893/1875000$ & \texttt{5:1/2; 8:1; 9:1/2; 16:1; 22:1/2; 24:1} \\
$2;3,2$ & $18424903/180000000$ & \texttt{7:1/2; 8:1/2; 10:1/6; 12:1/3; 15:1/6; 17:7/6; 18:7/6; 20:1/6} \\
$2;3,3$ & $4166537/30000000$ & \texttt{7:1/2; 8:1/2; 15:1/2; 17:1; 18:3/2} \\
$2;3,4$ & $4166537/30000000$ & \texttt{1:1/2; 5:1/2; 7:1/2; 8:1/2} \\
$2;3,5$ & $3981323/90000000$ & \texttt{8:1/3; 10:1/6; 15:1/6; 17:2/3; 18:5/6; 22:1/2; 27:1/3; 28:1/3} \\
$2;4,0$ & $219893/1875000$ & \texttt{5:1/2; 8:1; 9:1/2; 15:1; 22:1/2; 23:1} \\
$2;4,1$ & $219893/1875000$ & \texttt{5:1/2; 8:1; 9:1/2; 16:1; 22:1/2; 24:1} \\
$2;4,2$ & $18424903/180000000$ & \texttt{7:1/2; 8:1/2; 10:1/6; 11:1/3; 15:1/6; 17:7/6; 18:7/6; 20:1/6} \\
$2;4,3$ & $4166537/30000000$ & \texttt{1:1/2; 5:1/2; 7:1/2; 8:1/2} \\
$2;4,4$ & $4166537/30000000$ & \texttt{5:1/2; 7:1/2; 8:1/2; 15:1/2; 18:1/2} \\
$2;4,5$ & $3981323/90000000$ & \texttt{8:1/3; 10:1/6; 15:1/6; 17:2/3; 18:5/6; 22:1/2; 27:1/3; 28:1/3} \\
$2;5,0$ & $219893/1875000$ & \texttt{8:1; 9:1/2; 15:1; 17:1; 18:1; 22:1/2; 23:1} \\
$2;5,1$ & $219893/1875000$ & \texttt{1:1/2; 8:1; 9:1/2; 16:1/2; 18:1/2; 22:1/2; 24:1} \\
$2;5,2$ & $1203113/90000000$ & \texttt{5:1/3; 8:1/3; 10:1; 15:1; 18:2/3; 20:1; 25:1/3; 26:1/3} \\
$2;5,3$ & $5462257/135000000$ & \texttt{1:1/9; 5:1/2; 7:4/9; 8:1/3; 15:4/9; 23:1/3; 25:1/9; 26:1/9} \\
$2;5,4$ & $5462257/135000000$ & \texttt{1:1/3; 5:5/18; 7:4/9; 8:1/3; 18:2/9; 24:1/3; 25:1/9; 26:1/9} \\
$2;5,5$ & $1157393/3750000$ & \texttt{7:1/2; 8:1/2; 17:1; 18:1} \\
\bottomrule
\end{longtable}
\setlength{\tabcolsep}{6pt}
\normalsize


\section{Cited statements and their applications}\label{app:sources}

This appendix lists the inputs from earlier work and the
parameter substitutions at their applications. Results proved here are not
attributed to these sources.

For integers \(k,n\ge1\), define
\[
 P_k(n)=\prod_{j=0}^{k-1}(n+j),\qquad
 q_k(n)=\frac{P_k(n)}{\operatorname{rad}(P_k(n))}.
\]
For real \(X\ge2\), define
\[
 F_{k,k}(X)=\#\{1\le m<n\le X:
     \operatorname{rad}(P_k(m))=\operatorname{rad}(P_k(n))\}.
\]

\subsection*{The author's earlier papers}

Theorem~1 of \cite{Paper1} states that for every fixed integer \(k\ge2\), every
\(\varepsilon>0\), and \(X\ge2\),
\[
 F_{k,k}(X)\ll_{k,\varepsilon}X^{2/3+\varepsilon}.
\]
At \(k=2\), this is the earlier \(2/3\) baseline for \(F(X)\). It is
given for comparison only. No row below uses it.

Lemma~2.1 of \cite{Paper1} states that for squarefree \(Q\) and \(X\ge2\),
\[
 \#\{1\le z\le X:\operatorname{rad}(z(z+1))\mid Q\}
 \ll 2^{\omega(Q)}\log(2X),
\]
with an absolute implied constant. For the \(O_R\) row take
\(Q=\operatorname{rad}(n(n+1))\) and \(X=n\). The source equality of
radicals gives the divisibility condition for every partner \(m\).
Also \(Q\) is squarefree and \(Q\le Y(Y+1)\). Thus
\(2^{\omega(Q)}\ll_\eta Y^{2\eta}\). Choose the exponents of the
logarithm and divisor loss below the requested row loss. Then both are
absorbed. No other uniformity in \(Q\) is needed.

Lemma~2.3 of \cite{Paper1} gives, for fixed \(k\ge2\), \(X\ge2\),
\(K\ge\max(k,2)\), and \(0<\eta<1\),
\[
 \#\{n\le X:q_k(n)>K\}
 \ll_{k,\eta} XK^{-1+\eta}
      +K^2(1+\log K)^{k-1}+X^{1/2}.
\]
Lemma~2.4 of \cite{Paper1} gives
\[
 \#\{1\le m<n\le X:q_k(m),q_k(n)\le K,
       \operatorname{rad}(P_k(m))=\operatorname{rad}(P_k(n))\}
 \ll_k K^2\log(2X).
\]
Here $X\ge2$ and $K\ge1$, and $k\ge2$ is fixed. The stronger lower
restriction $K\ge\max(k,2)$ belongs only to Lemma~2.3 above.
These quotient-tail and Farey-spacing estimates are stated for context
only. Neither is a premise of the proof.

Theorem~3 of \cite{Paper1} gives, for every \(\eta>0\), uniformly for real
\(H\ge2\) and \(1\le V<H^2\),
\[
 \#\{H<t\le2H:\operatorname{rad}(t(t+1))\le V\}
 \ll_\eta H^\eta\left(\frac VH+H^{7/12}\right).
\]
For the outer low-radical estimate put \(H=Y/2\) and
\(V=Y^{5/3-S_0}\). For all sufficiently large \(Y\), \(H\ge2\),
\(V\ge1\), and
\(V/H^2=4Y^{-1/3-S_0}<1\). The right side becomes
\[
 (Y/2)^\eta\left(2Y^{2/3-S_0}+(Y/2)^{7/12}\right).
\]
The common factor \((Y/2)^\eta\) is kept. Fix a requested row
loss \(\varepsilon>0\). Choose the theorem parameter below both
\(\varepsilon/2\) and half the fixed gap \(2/3-S_0-7/12>0\). Then the
display is \(O_\varepsilon(Y^{2/3-S_0+\varepsilon})\). These are
the hypotheses and bound used in \eqref{eq:old:4_2} and
Appendix~\ref{app:lowrad}.

Theorems~1.1 and~1 of the author's earlier papers
\cite{Zhang712,ZhangE20} state, respectively,
\[
 N_1(X)\ll_\varepsilon X^{7/12+\varepsilon},\qquad
 N_1(X)\ll_\varepsilon X^{20/37+\varepsilon},
\]
where \(N_1(X)\) counts primitive positive triples \(x+y=z\) with
\(\operatorname{rad}(xyz)<z\) and \(\max(x,y,z)\le X\). These theorems
are cited for context only. The source rows below count equal-radical
pairs of consecutive-product arguments. They do not count primitive
\(abc\) triples. So neither theorem enters a row estimate.

The four-corner parametrization and the identities of
Section~\ref{sec:notation} follow
from the four gcd definitions. They give the coordinates used below. Their
derivation uses no further cited result.

\subsection*{External results}

We use Heath--Brown's published \cite[Lemma~2, p.~8]{HB2012}.
Let \(G\) be an absolutely irreducible
bihomogeneous integer form of bidegree \((a,b)\), with \(a,b\ge1\), and
let \(X\ge1\). Then the number of primitive pairs \((x_1,x_2)\),
\((y_1,y_2)\) with \(G=0\) and \(\max_i|x_i|\le X\) is
\[
 O_{a,b,\varepsilon}\bigl(X^{2/b+\varepsilon}\|G\|^\varepsilon\bigr).
\]
There is no size condition on the \(y\)-pair. Section~\ref{sec:intermediate} treats a factor
with corner degree \(a_c\ge2\) and coefficient degree \(b_q\ge1\).
There the coefficient-lattice pair is first and the primitive corner pair is
second. The reductions in Section~\ref{sec:intermediate} make both pairs primitive. The
first has radius \(O(T_1+1)\). The transformed factor height is at
most \(Y^C\) for a fixed \(C\). The source also gives \(T_1+1\le Y^{C'}\)
for a fixed \(C'\). The lemma with bidegree \((b_q,a_c)\) gives
\[
 O((T_1+1)^{2/a_c+\varepsilon}Y^{C\varepsilon}).
\]
Fix a requested total loss \(\eta>0\). Choose a fixed exponent
\(A_{\rm HB}\ge\max(C,C',1)\). Apply Heath--Brown with
\(0<\varepsilon<\eta/(4A_{\rm HB})\). Then its per-cell loss is at most
\(Y^{\eta/2}(T_1+1)\). Use the lattice census with its own loss
parameter \(\eta/2\). As \(2/a_c\le1\),
\[
 (T_1+1)^{2/a_c+\varepsilon}Y^{C\varepsilon}
 \le Y^{2A_{\rm HB}\varepsilon}(T_1+1)
\le Y^{\eta/2}(T_1+1).
\]
Summation by \eqref{eq:old:6_19} then gives the requested total
factor \(Y^\eta\). Section~\ref{sec:joint} uses the same pair order and degree
conditions. Its lattice census gets the separate \(\eta/2\) loss.
Thus the hypotheses \(X\ge1\), \(a,b\ge1\), primitivity and one-sided
size hold in Sections~\ref{sec:joint} and \ref{sec:intermediate}.

Section~\ref{sec:long} uses the same theorem for each absolutely irreducible factor
there with \(a_c\ge2\) and \(b_q\ge1\). The pair order and the
primitivity are the same. Let \(H_7\) be the fixed exponent in the factor-height
bound \(\lVert F\rVert\le Y^{H_7}\). The source bounds give
\(U_0,V_*,N\le Y\) for sufficiently large \(Y\), since
\(U_0\le f\le M\), \(V_*\le4g\le4M\),
\(M\ll Y^{q_{\max}}\), \(q_{\max}<1\), and \(N=\lceil Y^{a_{\rm grid}}\rceil\)
with \(a_{\rm grid}<1\). Set the temporary census loss parameter
\(\eta_{\rm census}=1\) in \eqref{eq:old:7_12}. This choice serves only
the coarse height bound below. It differs from the later allocated
loss \(\eta\). It yields
\(\sum_i(T_{1,i}+1)\ll Y^{5/2}\). After raising the fixed lower
threshold, every \(T_{1,i}+1\le Y^3\). This coarse bound serves only to
give this pointwise height cap. Set
\(A_{\rm HB}=H_7+3\). For a
requested total loss \(\eta>0\), take
\(0<\varepsilon<\eta/(2A_{\rm HB})\) in the Heath--Brown estimate. Then
\[
 (T_{1,i}+1)^{2/a_c+\varepsilon}\lVert F\rVert^\varepsilon
 \le Y^{(H_7+3)\varepsilon}(T_{1,i}+1)
 \le Y^{\eta/2}(T_{1,i}+1),
\]
and \eqref{eq:old:7_12} with loss parameter \(\eta/2\) pays for
the cell sum. The combined loss is \(Y^\eta\). Section~\ref{sec:long} treats the
cases \(a_c=0\), \(a_c=1\), and \(b_q=0\) separately. The same result
applies there.

The unnumbered theorem of Bombieri--Schmidt \cite[pp.~69--70]{BS1987} reads as follows. Let
\(F\in\mathbb Z[x,y]\) be an irreducible binary form of degree
\(r\ge3\) and let \(h\ge1\). Then the number of primitive integer solutions of
\(|F(x,y)|=h\) is at most \(C_{\rm BS}r^{1+\omega(h)}\), with absolute
\(C_{\rm BS}\). The 1989 correction to their proof of Lemma~7
\cite{BS1989} changes no
hypothesis of this theorem. In Section~\ref{sec:intermediate} the form is the irreducible
binary cubic \(c_1x^3-c_2z^3\). Its content and the
common divisor of \(x,z\) are removed first. The right side is nonzero. The theorem gives
\(O(3^{1+\omega(|h|)})\) primitive solutions. The reducible case
occurs when both coprime coefficients are cubes. Section~\ref{sec:intermediate} counts it
separately by factoring the difference of cubes.

Tao's Lemma~2.10 \cite[Lemma~2.10]{Tao} reads as follows. Let \(x\) be large. Let \(a,b\) be
positive integers and \(h\) a nonzero integer, all of polynomial size in
\(x\). Then
\[
 \#\{(n,m)\in\mathbb Z_{>0}^2:an^2+h=bm^2,\ n\le x\}
 \ll x^{o(1)}.
\]
Sections~\ref{sec:outer}--\ref{sec:long} use the fixed-height Theorem~\ref{thm:fibre},
proved in Appendix~\ref{app:fibre}. Tao's statement is cited as the
precedent. The applications are as follows. In Section~\ref{sec:outer}, the
outer fibre \eqref{eq:old:4_5} uses
\[
(a,b,h,n,m)=(P,Q_{\rm prod},Q_{\rm prod}-P,2n+1,2m+1),
\qquad x=2Y+1,\qquad A_{\rm ht}=4.
\]
Here \(P,Q_{\rm prod},|Q_{\rm prod}-P|\le4Y^4\),
\(2n+1\le2Y+1\), and the shift is nonzero.
In Section~\ref{sec:joint}, the two applications use, respectively,
\[
(a,b,h,n,m)=(\alpha,\beta,u,A,D),\qquad
(a,b,h,n,m)=(\gamma,\delta,t,B,C),\qquad x=2Y,\quad A_{\rm ht}=2.
\]
The equations are \(\alpha A^2+u=\beta D^2\) and
\(\gamma B^2+t=\delta C^2\). Section~\ref{sec:joint} gives
\(\alpha,\beta,u,\gamma,\delta,t\le(Y+1)^{21/20}\le(2Y)^2\) and
\(A,B\le Y+1\le2Y\) for \(Y\ge2\). Thus the fixed height exponent is
\(A_{\rm ht}=2\). The bounded variables are \(A\) and \(B\).
In Sections~\ref{sec:intermediate}--\ref{sec:long}, every corner-degree-zero use takes
\[
(a,b,h,n,m)=(f,g,s,J,K),\qquad x=2Y,\qquad A_{\rm ht}=2.
\]
The source gives positive \(f,g,s\), nonzero \(s\),
\(f,g,s\le(Y+1)^2\), and \(J,K\le Y+1\). In the corner-degree-zero
branches the ordered residual factors first fix \(L\). The selected
shift is then counted among \(s\mid L\). As \(L\) has polynomial
height in \(Y\), for every \(\eta>0\) the divisor bound gives the
uniform estimate \(\tauDiv(L)\ll_\eta Y^\eta\). The proof here uses the self-contained
Theorem~\ref{thm:fibre}.

Reuss's Theorem~1 is the determinant-method input behind this low-radical
estimate; we use the consequence stated in Theorem~3 of \cite{Paper1},
as applied in Appendix~\ref{app:lowrad}.

\section*{Use of AI tools}

The author devised the main argument. Codex (OpenAI, GPT-6 Astra) carried
out numerical checks of identities and of the cubic family, and searched
the literature. Claude Code (Anthropic, Claude Opus 5.5) reviewed the
arguments and references. The author takes responsibility for the proof
and references.

\section*{Statements and declarations}
\textbf{Competing interests.} The author has no competing interests.

\textbf{Funding.} No funding was received for this work.

\begin{thebibliography}{20}
\bibitem{Paper1} R.~Zhang, \emph{Equal radicals of consecutive products}
(2026), \url{https://doi.org/10.5281/zenodo.23758120}.
\bibitem{Zhang712} R.~Zhang, \emph{A primitive-ratio cut and a $7/12$
bound for the abc exceptional set}, preprint (2026),
\url{https://doi.org/10.5281/zenodo.23247752}.
\bibitem{ZhangE20} R.~Zhang, \emph{Cubic sieves, a primitive-ratio cut
and a $20/37$ bound for the abc exceptional set}, preprint (2026),
\url{https://doi.org/10.5281/zenodo.23266656}.
\bibitem{HB2012} D.~R. Heath-Brown, \emph{Square-free values of $n^2+1$},
Acta Arith. \textbf{155} (2012), no.~1, 1--13,
\url{https://doi.org/10.4064/aa155-1-1}.
\bibitem{BS1987} E.~Bombieri and W.~M. Schmidt, \emph{On Thue's equation},
Invent. Math. \textbf{88} (1987), no.~1, 69--81,
\url{https://doi.org/10.1007/BF01405092}.
\bibitem{BS1989} E.~Bombieri and W.~M. Schmidt, \emph{Correction to
On Thue's equation}, Invent. Math. \textbf{97} (1989), no.~2, 445,
\url{https://doi.org/10.1007/BF01389050}.
\bibitem{Reuss} T.~Reuss, \emph{Pairs of $k$-free numbers, consecutive square-full
numbers}, preprint (2014), arXiv:1212.3150v2.
\bibitem{Tao} T.~Tao, \emph{Products of consecutive integers with unusual anatomy},
preprint (2026), arXiv:2603.27990v3.
\end{thebibliography}
\end{document}
